%%%%%%%%%%%%%%%%%%%%%%%%%%%%%%%%%%%%%%%%%%%%%%%%%%%%%%%%%%%%%%%%%%%%%%%%%%%
%                                                                         %
%                      Chapter 4 - Section 1                              %
%                                                                         %
%%%%%%%%%%%%%%%%%%%%%%%%%%%%%%%%%%%%%%%%%%%%%%%%%%%%%%%%%%%%%%%%%%%%%%%%%%%


\chapter{Differential Equations}


The rate equations with which we began our study of calculus are
called {\bf differential equations}\index{differential equation} when
we identify the rates of change that appear within them as derivatives
of functions.  Differential equations are essential tools in many area
of mathematics and the sciences.  In this chapter we explore three of
their important uses:
\begin{itemize}

\item {\bf Modelling} problems using differential equations;

\item {\bf Solving} differential equations, both through numerical
  techniques like Euler's method and, where possible, through finding
  formulas which make the equations true;

\item {\bf Defining} new functions by differential equations.

\end{itemize}
We also introduce two important functions---the {\bf exponential
  function} and the {\bf logarithmic function}---which play central
roles in the theory of solving differential equations.  Finally, we
introduce the operation of {\bf antidifferentiation} as an important
tool for solving some special kinds of differential equations.

\section{Modelling with Differential Equations}
\index{model}

To analyze the way an infectious disease spreads through a population,
we asked how three quantities $S$, $I$, and $R$ would vary over time.
This was difficult to answer; we found no simple, direct relation
between $S$ (or $I$ or $R$) and $t$.  What we {\em did\/} find,
though, was a relation between the variables $S$, $I$, and $R$ and
their rates $S'$, $I'$, and $R'$.  We expressed the relation as a set
of rate equations.  Then, given the rate equations and initial values
for $S$, $I$, and $R$, we used Euler's method to estimate the values
at any time in the future.  By constructing a sequence of successive
approximations, we were able to make these estimates as accurate as we
wished.

There are two ideas here.  The first is that we could write down
equations for the rates of change that reflected important features of
the process we sought to model.  The second is that these equations
{\em determined\/} the variables as functions of time, so we could
make predictions about the real process we were modelling.  Can we
apply these ideas to other processes?

To answer this question, 
%
\mar{Differential equations and initial value problems}
%
it will be helpful to introduce some new terms.  What we have been
calling rate equations are more commonly called {\bf differential
  equations}.  (The name is something of an historical accident.
Since the equations involve functions and their derivatives, we might
better call them {\em derivative\/} equations.)  Euler's method treats
the differential equations for a set of variables as a prescription
for finding future values of those variables.  However, in order to
get started, we must always specify the initial values of the
variables---their values at some given time.  We call this
specification an {\bf initial condition}.\index{initial condition} The
differential equations together with an initial condition is called an
{\bf initial value problem}.\index{initial value problem} Each initial
value problem determines a set of functions which we find by using
Euler's method.

\aside{If we use Leibniz's notation for derivatives, a differential
  equation like $S' = -aSI$ takes the form $dS/dt = -aSI$.  If we then
  treat $dS/dt$ as a quotient of the individual {\em
    differentials\/}\index{differential} $dS$ and $dt$ (see
  page~\pageref{differential}), we can even write the equation as $dS
  = -aSI\,dt$.  Since this expresses the differential $dS$ in terms of
  the differential $dt$, it was natural to call it a differential
  equation.  Our approach is similar to Leibniz's, except that we
  don't need to introduce infinitesimally small quantities, which
  differentials were for Leibniz.  Instead, we write $\Delta S \approx
  -aSI\,\Delta t$ and rely on the fact that the accumulated error of
  the resulting approximations can be made as small as we~like.}

To illustrate how differential equations can be used to describe a
wide range of processes in the physical, biological, and social
sciences, we'll devote this section to a number of ways to model and
analyze the long-term behavior of animal populations.  To be specific,
we will talk about rabbits and foxes, but the ideas can be adapted to
the population dynamics of virtually all living things (and many
non-living systems as well, such as chemical reactions).

In each model, we will begin by identifying variables that describe
what is happening.  Then, we will try to establish how those variables
change over time.  Of course, no model can hope to capture every
feature of the process we seek to describe, so we begin simply.  We
choose just one or two elements that seem particularly important.
After examining the predictions of our simple model and checking how
well they correspond to reality, we make modifications.
%
\mar{Models can\\ provide successive approximations\\ to reality}\index{successive approximation}
%
We might include more features of the population dynamics, or we might
describe the same features in different ways.  Gradually, through a
succession of refinements of our original simple model, we hope for
descriptions that come closer and closer to the real situation we are
studying.

\subsection{Single-species Models: Rabbits}
\index{population model!single-species}

\noindent
{\bf The problem}.  If we turn 2000 rabbits\index{rabbits} loose on a
large, unpopulated island that has plenty of food for the rabbits, how
might the number of rabbits vary over time?  If we let $R=R(t)$ be the
number of rabbits at time $t$ (measured in months, let us say), we
would like to be able to make some predictions about the function
$R(t)$.  It would be ideal to have a formula for $R(t)$---but this is
not usually possible.  Nevertheless, there may still be a great deal
we can say about the behavior of $R$.  To begin our explorations we
will construct a model of the rabbit population that is obviously too
simple.  After we analyze the predictions it makes, we'll look at
various ways to modify the model so that it approximates reality more
closely.

\medskip
\noindent{\bf The first model}.  Let's assume that, at any time $t$,
%
\mar{\vspace{16pt} Constant\\ per capita growth}
%
the rate at which the rabbit population changes is simply proportional
to the number of rabbits present at that time.  For instance, if there
were twice as many rabbits, then the {\em rate\/} at which new rabbits
appear will also double. In mathematical terms, our assumption takes
the form of the differential equation
\[
     \dfrac{dR}{dt} = k\,R 
		\ \frac{\mbox{rabbits}}{\mbox{month}}.
		\leqno(1) 
\]
The multiplier $k$ is called the {\bf per capita growth
  rate}\index{growth rate!per capita} (or the {\bf
  reproductive\index{reproductive rate} rate}), and its units are
rabbits per month per rabbit.  Per capita growth is discussed in
exercise~22 in chapter~1, section~2.

For the sake of discussion, let's suppose that $k = .1$ rabbits per
month per rabbit.  This assumption means that, on the average, one
rabbit will produce .1 new rabbits every month.  In the $S$-$I$-$R$
model of chapter 1, the reciprocals of the coefficients in the
differential equations had natural interpretations.  The same is true
here for the per capita growth rate.  Specifically, we can say that
$1/k = 10$ months is the average length of time required for a rabbit
to produce one new rabbit.

Since there are 2000 rabbits at the start, we can now state a clearly
defined initial value problem for the function $R(t)$:
\[
	\dfrac{dR}{dt} = .1\,R \qquad \qquad R(0) = 2000.
\]
By modifying the program SIRPLOT, 
%
\mar{Use Euler's method\\ to find $R(t)$}\index{Euler's method}
%
we can readily produce the graph of the function that is determined by
this problem.  Before we do that, though, let's first consider some of
the implications that we can draw out of the problem without the
graph.

Since $R'(t) = .1\, R(t)$ rabbits per month and $R(0)=2000$ rabbits,
we see that the initial rate of growth is $R'(0)=200$ rabbits per
month.  If this rate were to persist for 20 years (= 240 months), $R$
would have increased by
\[
	\Delta R = 240 \ \hbox{months} \times  200 \ 	
	\frac{\mbox{rabbits}}{\mbox{month}}  =  
	48000 \ \mbox{rabbits}, 
\]
yielding altogether
\[
	R(240) = R(0) + \Delta R = 2000 + 48000 = 
		50000 \ \mbox{rabbits}
\]
at the end of the 20 years.  However, since the population $R$ is
always getting larger, the differential equation tells us that the
growth rate $R'$ will {\em also\/} always be getting larger.
Consequently, 50,000 is actually an underestimate of the number of
rabbits predicted by this model.

Let's restate our conclusions in a graphical form.  If $R'$ were
always 200 rabbits per month, the graph of $R$ plotted against $t$
would just be a straight line whose slope is 200 rabbits/month.  But
$R'$ is always getting bigger,
%
\mar{The graph of $R$\\ curves up}
%
so the slope of the graph should increase from left to right.  This
will make the graph curve upward.  In fact, SIRPLOT will produce the
following graph of $R(t)$:

\vspace*{100pt}
\special{psfile=../ps/calc1/pic178.ps}

\newpage

Later, we will see that the function $R(t)$ determined by this initial
value problem is actually an exponential
function\index{function!exponential} of $t$, and we will even be able
to write down a formula for $R(t)$, namely
\[
R(t) = 2000 \, (1.10517)^{\displaystyle t}.
\]

	This model is too simple to be able to describe what happens to a rabbit population very well.  One of the obvious difficulties is that it predicts the rabbit population just keeps growing---forever.  For example, if we used the formula for $R(t)$ given above, our model would predict that after 20 years ($t = 240$) there will be more than 50 {\em trillion\/} rabbits!  While rabbit populations can, under good conditions, grow at a nearly constant per capita rate for a surprisingly long time (this happened in Australia during the 19th century), our model is ultimately  unrealistic.

\aside{It is a good idea to think qualitatively about the functions determined by a differential equation and make some rough estimates before doing extensive calculations.  Your sketches may help you see ways in which the model doesn't correspond to reality.  Or, you may be able to catch errors in your computations if they differ noticeably from what your estimates led you to expect.}

\noindent\label{logistic-start}
{\bf The second model}.  One way out of the problem of unlimited growth is to modify equation (1) to take into account the fact that any given ecological system 
	\mar{The carrying capacity of the environment}
can support only some finite number of creatures over the long term.  This number is called the {\bf carrying capacity}\index{carrying capacity} of the system.  We expect that when a population has reached the carrying capacity of the system, the population should neither grow nor shrink.  At carrying capacity, a population should hold steady---its rate of change should be zero.   For the sake of specificity, let's suppose that in our example the carrying capacity of the island is 25,000 rabbits.

	What we would like to do, then, is to find an expression for $R'$ which is similar to equation (1) when the number of rabbits $R$ is near 2000, but which approaches 0 as $R$ approaches 25,000.  One model which captures these features is the {\bf logistic equation},\index{logistic equation} first proposed by the Belgian mathematician Otto Verhulst \index{Verhulst, O.} in 
	\mar{\vspace{22pt} Logistic growth}
1845:
	$$
	R' = k\,R \left(1-\frac{R}{b} \right) 
	\ \dfrac{\mbox{rabbits}}{\mbox{month}}. \leqno(2)
	$$
In this equation, the coefficient $k$ is called the {\bf natural
  growth rate}\index{growth rate!natural}.  It plays the same role as
the per capita growth rate in equation (1), and it has the same
units---rabbits per month per rabbit. The number $b$ is the {\bf
  carrying capacity}\index{carrying capacity}; it is measured in
rabbits.  (We first saw the logistic equation on pages
\pageref{logistic-begin}--\pageref{logistic-end}.)  Notice also that
we have written the derivative of $R$ in the simpler form $R'$, a
practice we will continue for the rest of the section.

	If the carrying capacity of the island is 25,000 rabbits, and if we keep the natural growth rate at .1 rabbits per month per rabbit, then the logistic equation for the rabbit population is
	$$
	R' = .1\,R \left(1-\frac{R}{25000} \right) 
		\ \dfrac{\mbox{rabbits}}{\mbox{month}}.
	$$
Check to see that this equation really does have the behavior claimed for it---namely, that a population of 25,000 rabbits neither grows or declines.  Notice also that $R'$ is positive as long as $R$ is less than 25000, so the population increases.  However, as $R$ approaches 25000, 
	\mar{The graph of $R(t)$ levels off near $R = 25000$}
$R'$ will get closer and closer to 0, so the graph will become nearly horizontal.  (What would happen if the island ever had more than 25,000 rabbits?)  

	These observations about the qualitative behavior of $R(t)$ are consistent with the following graph, produced by a modified version of the program SIRPLOT.  For comparison, we have also graphed the exponential function produced by the first model.  Notice that the two graphs ``share ink'' when $R$ near 2000, but diverge later on. 

\vspace*{3in}
\special{psfile=../ps/calc1/pic180.ps}

	By modifying the program SIRVALUE, we can even get numerical answers to specific questions about the two models.  For example, after 30 months under constant per capita growth, the rabbit population will be more than 40,000---well beyond the carrying capacity of the island.  Under logistic growth, though, the population will be only about 16,000.  

	In the following figure we display several functions that are determined by the logistic equation 
	$$
	R' = .1\,R \left(1-\frac{R}{25000} \right) 
	$$
when different initial conditions are given.  Each graph therefore predicts the future for a different initial population $R(0)$.   One of the graphs is just the $t$-axis itself.  What does this graph predict about the rabbit population?  What other graph is just a straight line, and what initial population will lead to this line?

\aside{While the logistic equation above was developed to model a physical problem in which only values of $R$ with $R\geq 0$ have any meaning, the mathematical problem of finding solutions for the resulting differential equation makes sense for all values of $R$.  We have drawn three graphs resulting from initial values $R(0) < 0$.  While this growth behavior of `anti-rabbits' is of little practical interest in this case, there may well be other physical problems of an entirely different sort which lead to the same mathematical model, and in which the solutions below the $t$--axis are crucial.}

\vspace*{3in}
\special{psfile=../ps/calc1/pic181.ps}

\begin{center}
  Solutions to the logistic equation $R'=.1R\,(1-R/25000)$ 
\end{center}
\label{logistic-finish}

\subsection{Two-species Models: Rabbits and Foxes}
\index{population model!two-species}

	No species lives alone in an environment, 
	\mar{Introduce predators}\index{predator}
and the same is true of the rabbits on our island.  The rabbit population will probably have to deal with predators of various sorts.  Some are microscopic---disease organisms, for 
example---while others loom as obvious threats.  We will enrich our population model by adding a second 
species---foxes---that will prey on the rabbits.\index{foxes}  We will continue to suppose that the rabbits live on abundant native vegetation, and we will now assume that the rabbits are the sole food supply of the foxes.  Can we say what will happen?  Will the number of foxes and rabbits level off and reach a ``steady state'' where their numbers don't vary?  Or will one species perhaps become extinct?

	Let $F$ denote the number of foxes, and $R$ the number of rabbits.  As before, measure the time $t$ in months.  Then $F$ and $R$ are functions of $t$: $F(t)$ and $R(t)$.  We seek differential equations that describe how the growth rates $F'$ and $R'$ are related to the population sizes $F$ and $R$.  We make the following assumptions.
\begin{itemize}
	\item In the absence of foxes, the rabbit population grows logistically.
	\item  The population of rabbits declines at a rate proportional to the product $R\cdot F$.  This is reasonable if we assume rabbits never die of old age---they just get a little too slow.  Their death rate, which depends on the number of fatal encounters between rabbits and foxes, will then be approximately proportional to both $R$ and $F$---and thus to their product.  (This is the same kind of interaction effect we used in our epidemic model to predict the rate at which susceptibles become infected.)
	\item In the absence of rabbits, the foxes die off at a rate proportional to the number of foxes present.  
	\item  The fox population increases at a rate proportional to the number of encounters between rabbits and foxes.  To a first approximation, this says that the birth rate in the fox population depends on maternal fox nutrition, and this depends on the number of rabbit-fox encounters, which is proportional to $R\cdot F$.
\end{itemize}

Our assumptions are about birth and death rates, so we can convert
them quite naturally into differential equations.  Pause here and
check that the assumptions translate into these differential
equations:

\newpage

\mbox{}
\vspace{-35pt}

\label{L-V bounded growth-begin}  
\begin{align*}
 R' & =  a \, R \left(1 - \frac{R}{b}\right) -cRF = a\,R  -  \frac{a}{b}\,R^2 - c\,RF   \\
 F' & =  d\,RF - e\,F   
\end{align*}
These are 
%
\mar{Lotka--Volterra equations with bounded growth} 
%
the {\bf Lotka--Volterra equations with bounded growth}.\index{Lotka--Volterra equations}  The coefficients $a$, $b$, $c$, $d$, and $e$  are {\bf parameters}\index{parameter}---constants that have to be determined through field observations in particular circumstances.  

\medskip
\noindent {\bf An example}.  To see what kind of predictions the Lotka--Volterra equations make, we'll work through an example with specific values for the parameters.  Let 
	$$
	\begin{array}{rcrl}
  	a &=& .1 & \mbox{rabbits per month per rabbit} \\
  	b &=& 10000 & \mbox{rabbits} \\
	c &=& .005 & \mbox{rabbits per month per rabbit-fox} \\ 
	d &=& .00004 & \mbox{foxes per month per rabbit-fox} \\ 
	e &=& .04 & \mbox{foxes per month per fox} 
	\end{array}
	$$
(Check that these five parameters have the right units.)  These choices give us the specific differential equations
	\begin{align*}
         R' & =  .1\,R  - .00001\,R^2 - .005\,RF   \\
         F' & =  .00004\,RF - .04\,F
	\end{align*}

	To use this model to follow $R$ and $F$ into the future, we need to know the initial sizes of the two populations.  Let's suppose that there are 2000 rabbits and 10 foxes at time $t = 0$.  
	\mar{The graphs \\ of $R$ and $F$}
Then the two populations will vary in the following way over the next 250 months.

%
% put figure at bottom of this page
%

\vspace{150pt}

\special{psfile=../ps/calc1/rf1.ps}

\newpage

%\vspace*{3in}
%\special{psfile=../ps/calc1/pic185.ps}

%\begin{center}
% Rabbit and fox populations as a function of time 
%\end{center}

	A variant of the program SIRPLOT was used to produce these graphs.  Notice that it plots $100F$ rather than $F$ itself.  This is because the number of foxes is about 100 times smaller than the number of rabbits.  Consequently, $100F$ and $R$ are about the same size, so their graphs fit nicely together on the same screen. 

	The graphs have several interesting features.  There are different scales for the $R$ and the $F$ values, because the program plots $100F$ instead of $F$.  The peak fox population is about 30, while the peak rabbit population is about 2300.  The rabbit and fox populations rise and fall in a regular manner.  They rise and fall less with each repeat, though, and if the graphs were continued far enough into the future we would see $R$ and $F$ level off to nearly constant values.  

	The illustration below shows what happens 
	\mar{How rabbits respond to changes in the initial fox population}
to an initial rabbit population of 2000 in the presence of three different initial fox populations $F(0)$.  Note that the peak rabbit populations are different, and they occur at different times.  The size of the intervals between peaks also depends on $F(0)$.  

\vspace*{3in}
\special{psfile=../ps/calc1/rf2.ps}


\begin{center}
Rabbit populations for different initial fox populations 
\end{center}

We have looked at three models, each a refinement of the preceding
one.  The first was the simplest.  It accounted only for the rabbits,
and it assumed the rabbit population grew at a constant per capita
rate.  The second was also restricted to rabbits, but it assumed
logistic growth to take into account the carrying capacity of the
environment.  The third introduced the complexity of a second species
preying on the rabbits.  In the exercises you will have an opportunity
to explore these and other models.  Remember that when you use Euler's
method to find the functions determined by an initial value problem,
you must construct a sequence of successive approximations, until you
obtain the level of accuracy desired.
\label{L-V bounded growth-end}  

\subsection{Exercises}

\subsubsection{Single-species models}

\vspace{-10pt}
\num  {\bf Constant per capita Growth}.  This question considers the initial value problem given in the text:
	$$
	R' = 0.1\,R \ \mbox{rabbits per month}; \qquad
		R(0) = 2000 \ \mbox{rabbits}.
	$$

\sub  Use Euler's method to determine how many rabbits there are after 6 months.  Present a table of successive approximations from which you can read the exact value to whole-number accuracy.  

\sub  Determine, to whole-number accuracy, how many rabbits there are after 24 months.

\sub  How many months does it take for the rabbit population to reach 25,000?

\num {\bf Logistic Growth}.\index{logistic growth}  The following questions concern  a rabbit population described by the logistic model
	$$
	R' = 0.1\,R \left( 1 - \dfrac{R}{25000} \right) 
		\ \mbox{rabbits per month}.
	$$

\sub  What happens to a population of 2000 rabbits after 6 months, after 24 months, and after 5 years?   To answer each question, present a table of successive approximations that allows you to give the exact value to the nearest whole number.  

\sub Sketch the functions determined by the logistic equation if you start with either 2000 or 40000 rabbits.  (Suggestion: you can modify the program SIRPLOT to answer this question.)  Compare the two functions.  How do they differ?  In what ways are they similar?

\num {\bf Seasonal Factors} Living conditions for most wild
populations are not constant throughout the year---due to factors like
drought or cold, the environment is less supportive during some parts
of the year than at others.  Partially in response to this, most
animals don't reproduce uniformly throughout the year.  This problem
explores ways of modifying the logistic model to reflect these facts.

\sub For the eastern cottontail rabbit, most young are born during the
months of March--May, with reduced reproduction during June--August,
and virtually no reproduction during the other six months of the year.
Write a program to generate the solution to the differential equation
$R'=k(1-R/25000)$, where $k=.2$ during March, April, and May; $k=.05$
during June, July, and August; and $k=0$ the rest of the year.  Start
with an initial population of 2000 rabbits on January 1.  You may find
that using the {\tt IF ... THEN} construction in your program is a
convenient way to incorporate the varying reproductive rate.  \sub How
would you modify the model to take into account the fact that rabbits
don't reproduce during their first season?

\num {\bf World population}.\index{population!world} The world's
population in 1990 was about 5 billion, and data show that birth rates
range from 35 to 40 per thousand per year and death rates from 15 to
20.  Take this to imply a net annual growth rate of 20 per thousand.
One model for world population assumes constant per capita growth,
with a per capita growth rate of 20/1000 = 0.02.

\sub  Write a differential equation for $P$ that expresses this assumption.  Use $P$ to denote the world population, measured in billions.  

\sub  According to the differential equation in (a), at what rate (in billions of persons per year) was the world population growing in 1990?

\sub  By applying Euler's method to this model, using the initial value of 5 billion in 1990, estimate the world population in the years 1980, 2000, 2040, and 2230.  Present a table of successive approximations that stabilizes with one decimal place of accuracy (in billions).  What step size did you have to use to obtain this accuracy?

\num  {\bf Supergrowth}.\index{supergrowth}  Another model for the world population, one that actually seems to fit recent population data fairly well, assumes ``supergrowth''---the rate $P'$ is proportional to a {\em higher power} of $P$, rather than to $P$ itself.  The model is
	$$
	P'=.015 \, P^{1.2}.
	$$
As in the previous exercise, assume that $P$ is measured in billions, and the population in 1990 was about 5 billion.

\sub  According to this model, at what rate (in billions of persons per year) was the population growing in 1990?

\sub  Using Euler's method, estimate the world population in the years 1980, 2000 and 2040.  Use successive approximations until you have one decimal place of accuracy (in billions).  What step size did you have to use to obtain this accuracy?

\sub  Use an Euler approximation with a step size of 0.1 to estimate the world population in the year 2230.  What happens if you repeat your calculation with a step size of 0.01?  [Comment:  Something strange is going on here.  We will look again at this model in the next section.]


\subsubsection{Two-species models}

Here are some other differential equations that model a predator-prey
interaction between two species.

\medskip

\num {\bf The May Model}.\label{May model} This model has been
proposed by the contemporary ecologist, R.M. May,\index{May, R.} to
incorporate more realistic assumptions about the encounters between
predators (foxes) and their prey (rabbits).  So that you can work with
quantities that are about the same size (and therefore plot them on
the same graph), let $y$ be the number of foxes and let $x$ be the
number of rabbits {\em divided by\/} 100---we are thus measuring
rabbits in units of ``hectorabbits".

\aside{While a term like ``hectorabbits'' is deliberately whimsical,
  it echoes the common and sensible practice of choosing units that
  allow us to measure things with numbers that are neither too small
  nor too large.  For example, we wouldn't describe the distance from
  the earth to the moon in millimeters, and we wouldn't describe the
  mass of a raindrop in kilograms.}



In his model, May makes the following assumptions.
\begin{itemize}

\item In the absence of foxes, the rabbits grow logistically.

\item The number of rabbits a single fox eats in a given time period
  is a function $D(x)$ of the number of rabbits available.  $D(x)$
  varies from 0 if there are no rabbits available to some value $c$
  (the {\bf saturation value}) if there is an unlimited supply of
  rabbits.  The total number of rabbits consumed in the given time
  period will thus be $D(x)\cdot y$.

\item The fox population is governed by the logistic equation, and the
  carrying capacity is proportional to the number of rabbits.

 \end{itemize}

\sub Explain why $D(x) =\dfrac{c\,x}{x+d}$ ($d$ some constant) might
be a reasonable model for the function $D(x)$.  Include a sketch of
the graph of $D$ in your discussion.  What is the role of the
parameter $d$?  That is, what feature of rabbit -- fox interactions is
reflected by making $d$ smaller or larger?

\sub  Explain how the following system of equations incorporates May's assumptions.  
\begin{align*}
x' &= a\,x\left(1-\frac{x}{b}\right) - \frac{c\,xy}{x+d} \nonumber \\
y' &= e\,y\left(1-\frac{y}{fx}\right)
\end{align*}
The parameters $a$, $b$, $c$, $d$, $e$ and $f$ are all positive.
 
\sub Assume you begin with 2000 rabbits and 10 foxes.  (Be careful: $x(0) \neq 2000$.)  What does May's model predict will happen to the rabbits and foxes over time if the values of the parameters are $a = .6$, $b = 10$, $c = .5$, $d = 1$, $e = .1$ and $f = 2$?  Use a suitable modification of the program SIRPLOT.

\sub  Using the same parameters, describe what happens if you begin with 2000 rabbits and 20 foxes; with 1000 rabbits and 10 foxes; with 1000 rabbits and 20 foxes.  Does the eventual long-term behavior depend on the initial condition?  How does the long-term behavior here compare with the long-term behavior of the two populations in the Lotka--Volterra model of the text?

\sub Using 2000 rabbits and 20 foxes as the initial values, let's see how the behavior of the solutions is affected by changing the values of the parameter $c$, the saturation value for the number of rabbits (measured in centirabbits, remember) a single fox can eat in a month.  Keeping all the other parameters ($a$, $b$, $d$, \ldots ) fixed at the values given above, get solution curves for $c=.5$, $c=.45$, $c=.4$,  \ldots , $c=.15$, and $c=.1$\,.   The solutions undergo a qualitative change somewhere between $c=.3$ and $c=.25$.  Describe this change.  Can you pinpoint the crucial value of $c$ more closely?  This phenomenon is an example of {\bf Hopf bifurcation}\index{Hopf bifurcation}, which we will look at more closely in chapter 8.  The May model undergoes a Hopf bifurcation as you vary each of the other parameters as well.  Choose a couple of them and locate approximately the associated bifurcation values.


\num {\bf The Lotka--Volterra Equations}.\index{Lotka--Volterra
  equations}\label{L-V} This model for predator and prey interactions is slightly
simpler than the ``bounded growth'' version we consider in the text.
It is important historically, though, because it was one of the first
mathematical population models, proposed as a way of understanding why
the harvests of certain species of fish in the Adriatic Sea exhibited
cyclical behavior over the years.  For the sake of variety, let's take
the prey to be hares and the predators to be lynx.\index{lynx}

	Let  $H(t)$ denote the number of hares\index{hare} at time $t$ and $L(t)$ the number of lynx.  This model, the basic Lotka--Volterra model, differs from the bounded growth model in only one respect: it assumes the hares would experience constant per capita growth if there were no lynx.  

\sub  Explain why the following system of equations incorporates the assumptions of the basic model.  (The parameters $a$, $b$, $c$,  and  $d$ are all positive.)
\begin{align*}
 H' & =  a\,H  -  b\,HL  \\
 L' & =  c\,HL - d\,L
\end{align*}
(These are called the {\bf Lotka--Volterra equations}.  They were developed independently by the Italian mathematical physicist Vito Volterra\index{Volterra, V.} in 1925--26, and 
by the mathematical ecologist and demographer Alfred James Lotka\index{Lotka, A.} a few years earlier.  Though simplistic, they form one of the principal starting points in ecological modeling.)

\sub Explain why $a$ and $b$ have the units hares per month per hare and hares per month per hare-lynx, respectively.  What are the units of $c$ and $d$?  Explain why.
\smallskip

\noindent
Suppose time $t$ is measured in months, and suppose the parameters
have values
\begin{alignat*}{2}
a &= .1 &     &\mbox{hares per month per hare} \\
b &= .005 &   &\mbox{hares per month per hare-lynx} \\ 
c &= .00004 &\quad  &\mbox{lynx per month per hare-lynx}\\ 
d &= .04 &    &\mbox{lynx per month per lynx} 
\end{alignat*}
This leads to the system of differential equations 
\begin{align*}
	H' & =  .1\,H - .005\,HL   \\
	L' & =  .00004\,HL - .04\,L.
 \end{align*}

\sub  Suppose that you start with 2000 hares and 10 
lynx---that is, $H(0)$ = 2000 and $L(0)$ = 10.  Describe what happens to the two populations.  A good way to do this is to draw graphs of the functions $H(t)$ and $L(t)$.  It will be convenient to have the Hare scale run from 0 to 3000, and the Lynx scale from 0 to 50.  If you modify the program SIRPLOT, have it plot $H$ and $60L$.

\vspace*{2.5in}
\special{psfile=../ps/calc1/hare1.ps}
\begin{center}
 Hare and lynx populations as a function of time 
\end{center}
You should get graphs like those above. Notice that the hare and lynx populations rise and fall in a fashion similar to the rabbits and foxes, but here they oscillate---returning {\em periodically\/} to their original values.  

\sub What happens if you keep the same initial hare population of
2000, but use different initial lynx populations?  Try $L(0) = 20$ and
$L(0) = 50$.  (In each case, use a step size of .1 month.)

\sub Start with 2000 hares and 10 lynx.  From part (c), you know the
solutions are periodic.  The goal of this part is to analyze this
periodic behavior.  You can do this with your program in part (c), but
you may prefer to replace the FOR-NEXT loop in your program by a
variety of DO-WHILE loops (see page~\pageref{DO-WHILE}).  First find
the maximum number of hares.  What is the length of one {\bf period}
for the hare population?  That is, how long does it take the hare
population to complete one cycle (e.g., to go from one maximum to the
next)?  Find the length of one period for the lynx.  Do the hare and
lynx populations have the same periods?

\sub  Plot the hare populations over time when you start with 2000 hares and, successively, 10, 20, and 50 lynx.  Is the hare population periodic in each case?  What is the period?  Does it vary with the size of the initial lynx population?


\subsubsection{Fermentation}\index{fermentation}

\label{ferment-begin}
	Wine is made by yeast; yeast\index{yeast} digests the sugars in grape juice and produces alcohol\index{alcohol} as a waste product.  This process is called fermentation.  The alcohol is toxic to the yeast, though, and the yeast is eventually killed by the alcohol.  This stops fermentation, and the liquid has become wine, with about 8--12\% alcohol.  

	Although alcohol isn't a ``species,'' it acts like a predator on yeast.  Unlike the other predator-prey problems we have considered, though, the yeast does not have an unlimited food supply.  The following exercises develop a sequence of models to take into account the interactions between sugar, yeast, and alcohol.  

\numa  In the first model assume that the sugar supply is not depleted, that no alcohol appears, and that the yeast simply grows {\em logistically}.  Begin by adding 0.5 lb of yeast to a large vat of grape juice whose carrying capacity is 10 lbs of yeast.  Assume that the natural growth rate of the yeast is 0.2 lbs of yeast per hour, per pound of yeast.  Let $Y(t)$ be the number of pounds of live yeast present after $t$ hours; what differential equation describes the growth of $Y$?  

\sub  Graph the solution $Y(t)$, for example by using a suitable modification of the program SIRPLOT.  Indicate on your graph approximately when the yeast reaches one-half the carrying capacity of the vat, and when it gets to within 1\% of the carrying capacity.

\sub  Suppose you use a second strain of yeast whose natural growth rate is only half that of the first strain of yeast.  If you put 0.5 lb of {\em this\/} yeast into the vat of grape juice, when will it reach one-half the carrying capacity of the vat, and when will it get to within 1\% of the carrying capacity?  Compare these values to the values produced by the first strain of yeast: are they larger, or smaller?  Sketch, on the same graph as in part (b), the way this yeast grows over time.


\numa  Now consider how the yeast produces alcohol.  Suppose that waste products are generated at a rate proportional to the amount of yeast present; specifically, suppose each pound of yeast produces 0.05 lbs of alcohol per hour.  (The other major waste product is carbon dioxide gas, which bubbles out of the liquid.)  Let $A(t)$ denote the amount of alcohol generated after $t$ hours.  Write a differential equation that describes the growth of $A$.

\sub  Consider the toxic effect of the alcohol on the yeast.  Assume that yeast cells die at a rate proportional to the amount of alcohol present, and also to the amount of yeast present.  Specifically, assume that, in each pound of yeast, a pound of alcohol will kill 0.1 lb of yeast per hour.  Then, if there are $Y$ lbs of yeast and $A$ lbs of alcohol, how many pounds of yeast will die in one hour?  Modify the original logistic equation for $Y$ (strain 1) to take this effect into account.  The modification involves subtracting off a new term that describes the rate at which alcohol kills yeast.  What is the new differential equation?

\sub  You should now have two differential equations describing the rates of growth of yeast and alcohol.  The equations are {\bf coupled}, in the sense that the yeast equation involves alcohol, and the alcohol equation involves yeast.  Assuming that the vat contains, initially, 0.5 lb of yeast and no alcohol, describe by means of a graph what happens to the yeast.  How close does the yeast get to carrying capacity, and when does this happen?  Does the fermentation end?  If so, when; and how much alcohol has been produced by that time?  (Note that since $Y$ will never get all the way to 0, you will need to adopt some convention like $Y\leq .01$ to specify the end of fermentation.)

\num  What happens if the rates of toxicity and alcohol production are different?  Specifically, increase the rate of alcohol production by a factor of five---from 0.05 to 0.25 lbs of alcohol per hour, per pound of yeast---and at the same time reduce the toxicity rate by the same factor---from 0.10 to 0.02 lb of yeast per hour, per pound of alcohol and pound of yeast.  How do these changes affect the time it takes for fermentation to end?  How do they affect the amount of alcohol produced?  What happens if only the rate of alcohol production is changed?  What happens if only the toxicity rate is reduced?

\numa  The third model will take into account that the sugar in the grape juice is consumed.  Suppose the yeast consumes .15 lb of sugar per hour, per lb of yeast.  Let $S(t)$ be the amount of sugar in the vat after $t$ hours.  Write a differential equation that describes what happens to $S$ over time.

\sub  Since the carrying capacity of the vat depends on the amount of sugar in it, the carrying capacity must now vary.  Assume that the carrying capacity of $S$ lbs of sugar is $.4\,S$ lbs of yeast.  How much sugar is needed to maintain a carrying capacity of 10 lbs of yeast?  How much is needed to maintain a carrying capacity of 1 lb of yeast?  Rewrite the logistic equation for yeast so that the carrying capacity is $.4\,S$ lbs, instead of 10 lbs, of yeast.  Retain the term you developed in 9.b to reflect the toxic impact of alcohol on the yeast.

\sub  There are now three differential equations.  Using them, describe what happens to .5 lbs of yeast that is put into a vat of grape juice that contains 25 lbs of sugar at the start.  Does all the sugar disappear?  Does all the yeast disappear?  How long does it take before there is only .01 lb of yeast?  How much sugar is left then?  How much alcohol has been produced by that time?
\label{ferment-end}

\subsubsection{Newton's law of cooling}  
\index{Newton's law of cooling}
\label{cooling coffee}

Suppose that we start off with a freshly brewed cup of coffee at
\mbox{$90^{\circ}$C} and set it down in a room where the temperature
is \mbox{$20^{\circ}$C}.  What will the temperature of the coffee be
in 20 minutes?  How long will it take the coffee to cool to
\mbox{$30^{\circ}$C}?

If we let the temperature of the coffee be $Q$ (in $^{\circ}$C), then
$Q$ is a function of the time $t$, measured in minutes.  We have $Q(0)
= 90^{\circ}$C, and we would like to find the value $t_1$ for which
$Q(t_1) = 30^{\circ}$C.

It is not immediately apparent how to give $Q$ as a function of $t$.
However, we can describe the {\em rate\/} at which a liquid cools off,
using {\bf Newton's law of cooling}:\index{Newton's law of cooling}
the rate at which an object cools (or warms up, if it's cooler than
its surroundings) is proportional to the {\em difference\/} between
its temperature and that of its surroundings.

\num In our example, the temperature of the room is
\mbox{$20^{\circ}$C}, so Newton's law of cooling states that $Q'(t)$
is proportional to $Q - 20$, the difference between the temperature of
the liquid and the room.  In symbols, we have
	$$
	Q' = -k\,(Q -20)
	$$
where $k$ is some positive constant.  

\sub  Why is there a minus sign in the equation?

\smallskip
\noindent
The particular value of $k$ would need to be determined
experimentally.  It will depend on things like the size and shape of
the cup, how much sugar and cream you use, and whether you stir the
liquid.  Suppose that $k$ has the value of \mbox{$.1^{\circ}$} per
minute per \mbox{$^{\circ}$C} of temperature difference. Then the
differential equation becomes:
	$$
	Q' = -.1(Q-20) \quad \mbox{$^{\circ}$C per minute}.
	$$

\sub\label{ex:coffee1} Use Euler's method to determine the temperature
$Q$ after 20 minutes.  Write a table of successive approximations with
smaller and smaller step sizes.  The values in your table should
stabilize to the second decimal place.

\sub\label{ex:coffee2} How long does it take for the temperature $Q$
to drop to \mbox{$30^{\circ}$C}?  Use a DO-WHILE loop\index{loop!DO-WHILE} to construct a table of successive approximations that
stabilize to the second decimal place.

\num On a hot day, a cold drink warms up at a rate approximately
proportional to the difference in temperature between the drink and
its surroundings.  Suppose the air temperature is 90$^\circ$F and the
drink is initially at 36$^\circ$F.  If $Q$ is the temperature of the
drink at any time, we shall suppose that it warms up at the rate
	$$
	Q' = -0.2(Q - 90) \quad 
			\mbox{$^\circ$F per minute}.
	$$ 
According to this model, what will the temperature of the drink be
after 5 minutes, and after 10 minutes.\label{ex:cold drink} In both
cases, produce values that are accurate to two decimal places.

\num  In our discussion of cooling coffee, we assumed that the coffee did not heat up the room.  This is reasonable because the room is large, compared to the cup of coffee.  Suppose, in an effort to keep it warmer, we put the coffee into a small insulated container---such as a microwave oven (which is turned off).  We must assume that the coffee {\em does\/} heat up the air inside the container.  Let $A$ be the air temperature in the container and $Q$ the temperature of the coffee.  Then both $A$ and $Q$ change over time, and Newton's law of cooling tells us the {\em rates\/} at which they change.  In fact, the law says that both $Q'$ and $A'$ are proportional to $Q-A$.  Thus,
	\begin{align*}
	Q' &= -k_1(Q - A) \\
	A' &=  k_2(Q - A),
	\end{align*}
where $k_1$ and $k_2$ are positive constants.

\sub  Explain the signs that appear in these differential equations.

\sub  Suppose  $k_1 = .3$ and $k_2 = .1$.  If $Q(0) = \mbox{$90^{\circ}$C}$  and  $A(0) = \mbox{$20^{\circ}$C}$, when will the temperature of the coffee be \mbox{$40^{\circ}$C}?  What is the temperature of the air at this time?  Your answers should be accurate to one decimal place.

\sub  What does the temperature of the coffee become eventually?  How long does it take to reach that temperature?


\subsubsection{$S$-$I$-$R$ revisited}

Consider the spread of an infectious disease that is modelled by the
$S$-$I$-$R$ differential equations
\begin{alignat*}{2}
S'  & =   -.00001\,SI,  &      &    \\
I'  & = \hphantom{-}.00001\,SI  &   - \, &  .08\,I,  \\
R'  & =           		&      &  .08\,I.
\end{alignat*}
Take the initial condition of the three populations to be
	\begin{align*}
     S(0) &= 35,\!400 \ \hbox{persons}, \\
	I(0) &= 13,\!500 \ \hbox{persons}, \\
	R(0) &= 22,\!100 \ \hbox{persons}.
	\end{align*}

\num How many susceptibles are left after 40 days?  When is the
largest number of people infected?  How many susceptibles are there at
that time?  Explain how you could determine the last number {\em
  without\/} using Euler's method.

\num What happens as the epidemic ``runs its course''?  That is, as
more and more time goes by, what happens to the numbers of infecteds
and susceptibles?

\num One of the principal uses of a mathematical model is to get a
qualitative idea how a system will behave with different initial
conditions.  For instance, suppose we introduce 100 infected
individuals into a population.  How will the spread of the infection
depend on the size of the population?  Assume the same $S$-$I$-$R$
differential equations that were used in the previous exercise, and
draw the graphs of $S(t)$ for initial susceptible population sizes
$S(0)$ ranging from 0 to 45,000 in increments of 5000 (that is, take
$S(0)$ = 0, 5000, 10000, \ldots , 45000).  In each case assume that
$R(0) = 0$ and $I(0) = 100$.  Use these graphs to argue that the
larger the initial susceptible population, the more rapidly the
epidemic runs its course.

\newcommand{\smax}{\mbox{\scriptsize max}}

\num Draw the graphs of $I(t)$ for the same initial conditions as in
the previous problem.  Using these graphs you can demonstrate that the
larger the susceptible population, the larger will be the fraction of
the population that is infected during the worst stages of the
epidemic.  Do this by constructing a table displaying $I_{\smax}$,
$t_{\smax}$, and $P_{\smax}$, where $I_{\smax}$ is the maximum value
of $I(t)$, $t_{\smax}$ is the time at which this maximum occurs (that
is, $I_{\smax} = I(t_{\smax})$), and $P_{\smax}$ is the ratio of
$I_{\smax}$ to the initial susceptible population:
$P_{\smax}=I_{\smax}/S(0)$.  The table below gives the first three
sets of values.  \addtolength{\arraycolsep}{3pt}
	$$
\begin{array}{cccc}
	S(0) & I_{\smax} & P_{\smax} & t_{\smax} \\[2pt] \hline
	\begin{array}{r}
	\rule{0pt}{14pt} 5\,000 \\ 10\,000 \\ 15\,000 
	\end{array}
	&
	\begin{array}{r}
	\rule{0pt}{14pt} 100 \\ 315 \\ 2071 
	\end{array}
	&
	\begin{array}{r}
	\rule{0pt}{14pt} 0.02 \\ 0.03 \\ 0.14 
	\end{array}
	&
	\begin{array}{r}
	\rule{0pt}{14pt} 0 \\ >100 \\ 66 
	\end{array}
	\\[-2pt]
	\vdots & \vdots & \vdots & \vdots 
\end{array}
	$$ 
\addtolength{\arraycolsep}{-3pt}% 
%
Your table should show that there is a time when over half the
population is infected if $S(0) = 45000$, while there is never a time
when more than one-fourth of the population is infected if $S(0) =
20000$.


\subsubsection{Constructing models}
\index{model}

Systems in which we know a number of quantities at a given time and
would like to know their values at a future time (or know at what
future time they will attain given values) occur in many different
contexts.  The following are some systems for discussion.  Can any of
these be modelled as initial value problems?  What information would
you need to resolve the question?  Make some reasonable assumptions
about the missing information and write down an initial value problem
which would model the system.

\num We deposit a fixed sum of money in a bank, and we'd like to know
how much will be there in ten years.

\num We know the diameter of the mold spot growing on a cheese
sandwich is 1/4 inch, and we'd like to know when its diameter will be
one inch.

\num We know the fecal bacterial and coliform concentrations in a
local swimming hole, and we'd like to know when they fall below
certain prescribed levels (which the Board of Health deems safe).

\num We know what the temperature and rainfall is today, and we'd like
to know what both will be one week from today.

\num We know what the winning lottery number was yesterday, and we'd
like to know what the winning number will be the day after tomorrow.

\num We know where the earth, sun, and moon are in relation to each
other now, and how fast and in what direction they are moving.  We
would like to be able to predict where they are going to be at any
time in the future.  We know the gravity of each affects the motions
of the others by determining the way their velocities are changing.

%%%%%%%%%%%%%%%%%%%%%%%%%%%%%%%%%%%%%%%%%%%%%%%%%%%%%%%%%%%%%%%%%%%%%%%%%%%
%                                                                         %
%                                Section 2                                %
%                                                                         %
%%%%%%%%%%%%%%%%%%%%%%%%%%%%%%%%%%%%%%%%%%%%%%%%%%%%%%%%%%%%%%%%%%%%%%%%%%%

\section{Solutions of Differential Equations}

\subsection{Differential Equations are Equations}

Until now, we have viewed a system of differential equations
%
\mar{Differential equations give instructions for Euler's method}
%
as a set of {\em instructions\/} for ``stepping into the future'' (or
the past).  Put another way, an initial value problem was treated as a
prescription for using Euler's method\index{Euler's method} to
determine a set of functions which were then given either graphically
or in tabular form.

In this section we take a new point of view: we will think of
differential equations as {\em equations\/}\index{equation} for which
we would like to find {\em solutions}\index{solution} in terms of
functions which can be given by explicit formulas.  While it is
unfortunately the case that most differential equations do not have
solutions which can be given by formulas, there are enough important
classes of equations where such solutions do exist to make them worth
studying.  When such solutions can be found, we have a very powerful
tool for examining the behavior of the phenomenon being modelled.

To see what this means, let's look first at equations in algebra.
%
\mar{Equations have {\em solutions}}
%
Consider the equation $x^2 = x + 6$.  As it stands, this is neither
true nor false.  We make it true or false, though, when we substitute
a particular number for $x$.  For example, $x = 3$ makes the equation
true, because $3^2 = 3 + 6$.  On the other hand, $x = 1$ makes the
equation false, because $1^2 \neq 1 + 6$.
Any number that makes an equation true is called a {\em solution\/} to
that equation.  In fact, $x^2 = x + 6$ has exactly two solutions: $x =
3$ and $x = -2$.

We can view differential equations the same way.  Consider, for
example, the differential equation
	$$
	\dfrac{dy}{dt} = \dfrac{1}{2y}.
	$$
Because it involves the expression $dy/dt$, we understand that $y$ is a function of $t$.  As it stands, the differential equation is neither true nor false.  We make it true or false, though, when we substitute a particular {\em function\/} for $y$.  For example, $y = \sqrt{t} = t^{1/2}$ makes the differential equation true.  To see this, first look at the left-hand side of the 
%
\mar{\vspace{9pt}Substitute $y = \sqrt{t}$\\ into the differential equation}
%
equation:
	$$
	\dfrac{dy}{dt} = \tfrac{1}{2} t^{-1/2} = 
		\dfrac{1}{2 \sqrt{t}}.
	$$
Now look at the right-hand side:
	$$
	\dfrac{1}{2y} = \dfrac{1}{2\sqrt{t}}.
	$$
The two sides of the equation are equal, so the substitution $y = \sqrt{t}$ makes the equation true.

	The function $y = t^2$, however, makes the differential equation {\em false}.  The left-hand side is
	$$
	\dfrac{dy}{dt} = 2t,
	$$
but the right-hand side is
	$$
	\dfrac{1}{2y} = \dfrac{1}{2t^2}.
	$$
Since $2t$ and $1/2t^2$ are different functions, the two sides are unequal and the equation is therefore false.  

	We say that $y = \sqrt{t}$ is a 
	\mar{A {\em solution\/} makes\\ the equation true}
{\bf solution} to this differential equation.  The function $y = t^2$ is {\em not\/} a solution.  To decide whether a particular function is a solution when the function is given by a formula, notice that we need to be able to differentiate the formula.  

\aside{If we view differential equations simply as instructions for carrying out Euler's method, we need only the microscope equation $\Delta y \approx y'\cdot \Delta t$ in order to find functions.  However, if we want to find functions that are solutions to differential equations from our new point of view, we first need to introduce the idea of the derivative and the rules for differentiating functions.}

	Just as an algebraic equation can have more than one solution, so can a differential equation.  In fact, we can show that $y = \sqrt{t + C}$ is a solution to the differential equation 
	$$
	\dfrac{dy}{dt} = \dfrac{1}{2y},
	$$
for any value of the constant $C$.  To evaluate the left-hand side $dy/dt$, we need the chain rule (chapter~3.6).  Let's write
	$$
	y = \sqrt{u} \qquad \mbox{where} \qquad u = t + C.
	$$
Then the left-hand side is the function
	$$
	\dfrac{dy}{dt} = \dfrac{dy}{du} \cdot \dfrac{du}{dt} =
	\dfrac{1}{2\sqrt{u}} \cdot 1 = \dfrac{1}{2\sqrt{t + C}}.
	$$
Since the right-hand side of the differential equation is
	$$
	\dfrac{1}{2y} = \dfrac{1}{2 \sqrt{t + C}},
	$$
the two sides are equal---no matter what value $C$ happens to have.  
	\mar{A differential equation can have infinitely many solutions}
This proves that every function of the form $y = \sqrt{t + C}$ is a solution to the differential equation.  Since there are infinitely many values that $C$ can take, the differential equation has infinitely many different solutions!  
\smallskip

	If a differential equation arises in modelling a physical or biological process, the variables involved must also satisfy an initial condition.  Suppose we add an initial condition to our differential equation:
	$$
	\dfrac{dy}{dt} = \dfrac{1}{2y} \qquad \mbox{and}
		\qquad y(0) = 5.
	$$
Does {\em this\/} problem have a solution---that is, can we find a function $y(t)$ that is a solution to the differential equation and also satisfies the condition $y(0) = 5$?  

	Notice $y = \sqrt{t}$ is not a solution to this new problem.  Although it satisfies the differential equation, it fails to satisfy the initial condition:
	$$
	y(0) = \sqrt{0} = 0 \neq 5.
	$$
Perhaps one of the other solutions to the differential equation will work.  When we evaluate the solution $y = \sqrt{t + C}$ at $t = 0$ we get
	$$
	y(0) = \sqrt{0 + C} = \sqrt{C}.
	$$

\pagebreak

\noindent
We want this to equal 5, and it will if $C = 25$.  
%
\mar{An initial value problem has\\ only one solution}
%
Thus, $y = \sqrt{t + 25}$ is a solution to the initial value problem.
Furthermore, the only value of $C$ which will make $y(0) = 5$ is $C =
25$, so the initial value problem has only one solution of the form
$\sqrt{t + C}$.  Here is the graph of this solution:
\begin{center}
	\begin{picture}(230,90)(-15,-10)

	\put(-15,0){\vector(1,0){210}}
	\put(100,0){\vector(0,1){70}}
	\put(193,-4){\makebox(0,0)[t]{\small $t$}}
	\put(96,70){\makebox(0,0)[r]{\small $y$}}

	\multiput(0,-1)(20,0){10}{\line(0,1){2}}
	\multiput(99,8.94)(0,8.94){7}{\line(1,0){2}}
	\put(-3,-3){\makebox(0,0)[t]{\footnotesize $-25$}}
	\put(100,44.7){\circle*{3}}
	\put(96,47){\makebox(0,0)[br]{\footnotesize $5$}}

	\begin{thicklines}
	\bezier{250}(0,0)(0,30)(180,60)
	\end{thicklines}
	\put(186,63){\makebox(0,0)[l]
		{\small $y = \sqrt{t + 25}$}}

	\end{picture}
\end{center}
As always, you can use Euler's method to find the function determined
by an initial value problem, and you can graph that function using the
program SIRPLOT, for example.  How will {\em that\/} graph compare
with this one?  In the exercises you can explore this question.
\medskip

\noindent
{\bf Checking solutions versus finding solutions}.  Notice that we
have only {\em checked\/} whether a given function solves an initial
value problem; we have not {\em constructed\/} a formula to solve the
problem.  By this point you are probably wondering where the given
solutions came from.

It is helpful to continue exploring the parallels with solutions of
algebraic equations.  In the case of the equation $x^2 = x + 6$, there
are, of course, methods to find solutions.  One possibility is to
rewrite $x^2 = x + 6$ in the form $x^2 - x - 6 = 0$.  By factoring
$x^2 - x - 6$ as
	$$
	x^2 - x - 6 = (x - 3)(x + 2)
	$$
we can see that either $x - 3 = 0$ (so $x = 3$), or $x + 2 = 0$ (so $x
= -2$).  Another method is to use the {\bf quadratic
  formula}\index{quadratic formula}
	$$
	x = \dfrac{-b \pm \sqrt{b^2 - 4ac}}{2a}
	$$
for the roots of the quadratic function $ax^2 + bx + c$.  In our case,
the quadratic formula yields
	$$
	x = \dfrac{-(-1) \pm \sqrt{1 - 4 \cdot 1 \cdot (-6)}} 
			{2 \cdot 1} 
		= \dfrac{1 \pm \sqrt{1 + 24\vphantom{(-6)}}}{2} 
		= \dfrac{1 \pm 5}{2},
	$$
so again we find that $x$ must be either 3 or $-2$.  

\pagebreak

Thus we have at least two different methods 
%
\mar{There are special methods to solve particular equations, but other methods work generally}
%
for finding solutions to this particular equation.  The methods we use
to solve an algebraic equation depend very much on the equation we
face.  For example, there is no way to find a solution to $\sin{x} =
2^x$ by factoring, or by using a ``magic formula'' like the quadratic
formula.  Nevertheless, there are methods that {\em do\/} work.  In
chapters 1 and 2 we dealt with similar problems by using a computer
graphing utility that could zoom in on the point of intersection of
two graphs.  In chapter 6 we will introduce another tool, the
Newton--Raphson method, for finding roots.  These are both powerful
methods, because they will work with nearly all algebraic equations.
It is important to recognize that these numerical methods really do
solve the problem, even though they do not give solutions in closed
form the way the quadratic formula does.

The situation is entirely analogous in dealing with differential
equations.  The methods we use to solve a differential equation depend
on the equation we face.  A course in differential equations provides
methods for finding formulas that solve many different kinds of
differential equations.  The methods are like the quadratic formula in
algebra, though---they give a complete solution, but they work only
with differential equations that have a very specific form.  This
course will not attempt to survey the methods that find such formulas,
although in the next sections we will see effective methods for
dealing with some special subcases.


It is important to realize, though, that Euler's method is always
there if we can't think of anything cleverer, and it really does
provide solutions.
%
\mar{Most initial value problem have one,\\ and only one, solution.  Euler's method\\ will find it.}
%
In fact, most initial value problems have one, {\em and only one},
solution, and Euler's method can be used to determine this unique
solution. If we can also find a formula for the solution, then it must
be the same solution as that produced by Euler's method.  In more
advanced courses you will see a proof that this is true in general,
provided some mild conditions are satisfied.  To emphasize the
importance of this idea, we give it a name:
\begin{center}
\begin{picture}(290,60)
\begin{thicklines}
\put(0,0){\framebox(290,60)
	{\shortstack
	{
	{\bf Existence and Uniqueness Principle\index{existence and uniqueness principle}} \\
	Under most conditions, an initial value problem  \\
	has one and only one solution.
	}}}
\end{thicklines}
\end{picture}
\end{center}
The existence and uniqueness principle \index{existence and uniqueness
  principle} is one of the most important mathematical results in the
theory of differential equations.

We will continue to rely primarily on Euler's method, which generates
solutions for nearly all differential equations.

However, there are clear benefits to having a formula for the solution
to a differential equation, allowing us to investigate questions that
we can't answer very well if we only have solutions given by Euler's
method.  In this section, we will look at some of those benefits.


\subsection{World Population Growth}

\subsubsection{Two models}

In the exercises in the last section, we looked at two different
models that seek to describe how the world population will
grow.\index{population!world} One model assumed constant per capita
growth---rate of change proportional to population size.  The other
assumed ``supergrowth"---rate of change proportional to a {\em
  higher\/} power of the population size.  Let's write $P$ for the
population size in the constant per capita growth model and $Q$ for
the population size in the supergrowth model.  In both cases, the
population is expressed in billions of persons and time is measured in
years, with $t=0$ in 1990.  In this notation, \mar{\vspace{30pt}
  Models for world population growth} the two models are
	$$
	\begin{array}{rll}
	\mbox{constant per capita:}  \qquad\quad &
		\dfrac{dP}{dt} =.02\,P
		& P(0) = 5; \\[12pt]
	\mbox{supergrowth:} & 
		\dfrac{dQ}{dt} = .015\,Q^{1.2}  
		& Q(0) = 5. 
	\end{array}
	$$
\index{growth!constant per capita}
\index{supergrowth}

By using Euler's method, we discover that the two models predict
fairly similar results over sixty years, although the supergrowth
model lives up to its name by predicting larger populations than the
constant per capita growth model as time passes:
	$$
	\begin{array}{ccc}
	\hphantom{-}t  & P & Q \\ [2pt] \hline
	\begin{array}{r}
	\rule{0pt}{14pt} -10 \\ 0 \\ 10 \\ 50 
	\end{array}
	&
	\begin{array}{r}
	\rule{0pt}{14pt} 4.09 \\ 5.00 \\ 6.11 \\ 13.59 
	\end{array}
	&
	\begin{array}{r}
	\rule{0pt}{14pt} 4.08 \\ 5.00 \\ 6.18 \\ 15.94
	\end{array}
	\end{array}
	$$
These estimates are accurate to one decimal place, and that level of
accuracy was obtained with the step size $\Delta t=.1$.

However, the predictions made by the models differ widely over longer
time spans.  If we use Euler's method to estimate the populations
after 240 years, we get
	$$
	\begin{array}{ccc}
	\Delta t	& P(240)	& Q(240)\label{Q(240)} \\
		[2pt] \hline
	\begin{array}{l}
	.1 \\ .01 \\ .001 
	\end{array}
	&
	\begin{array}{r}
	6.046 \times 10^2 \\ 
		6.073 \times 10^2 \\ 6.075 \times 10^2
	\end{array}
	&
	\begin{array}{r}
	\rule{0pt}{15pt} 1.979 \times 10^{10} \\ 
		2.573 \times 10^{11} \\ 3.825 \times 10^{11}
	\end{array}
	\end{array}
	$$
As the step size decreases from 0.1 to 0.01 to 0.001, the estimates of
the constant per capita growth model $P(240)$ behave as we have come
to expect: already three digits have stabilized.  But in the estimates
of the supergrowth model, not even one digit of $Q(240)$ has
stabilized.

In this section we will see that there are actually {\em formulas\/}
for the functions $P(t)$ and $Q(t)$. These formulas will illuminate
the reason behind the differences in speed of stabilization in the
estimates.

\subsubsection{A formula for the supergrowth model}
\index{formula}

	Without asking how the following formula might have been derived, let's check that it is indeed a solution to the supergrowth initial value problem.
 	$$
	Q(t) = \left( \frac{1}{\sqrt[5]{5}} - .003\,t
		\right)^{-5}
	$$
First of all, 
	\mar{Checking the\\ initial condition}
the formula satisfies the initial condition $Q(0)=5$:\index{initial condition}
	$$
	Q(0) = \left( \frac{1}{\sqrt[5]{5}} \right)^{-5} =
		(\sqrt[5]{5})^5 =5.
	$$

	To check that it 
	\mar{Checking the differential equation}
also satisfies the differential equation, we must evaluate the two sides of the differential equation	
	$$
	\dfrac{dQ}{dt} = .015\,Q^{1.2}.
	$$
Let's begin by evaluating the left-hand side.  
	\mar{Left-hand side}
To differentiate $Q(t)$, we will write $Q$ as a {\em chain\/} of functions:
	$$
	Q=u^{-5} \qquad \mbox{where} \qquad  
		u = \frac{1}{\sqrt[5]{5}}-.003 \,t.
	$$
Since $Q=u^{-5}$, $dQ/du = (-5)u^{-6}$.  Also, since $u$ is just a linear function of $t$ in which the multiplier is $-.003$, we have $du/dt = -.003$.  Consequently,
	\begin{align*}
	\dfrac{dQ}{dt} &= \dfrac{dQ}{du} \cdot \dfrac{du}{dt} \\
		&= (-5)\,u^{-6} \cdot (-.003) \\
		&= .015 \, u^{-6} 
	\end{align*}
Ordinarily, we would ``finish the job'' by substituting for $u$ its formula in terms of $t$.  However, in this case it is clearer to just leave the left-hand side in this form.

	To evaluate the right-hand side of the differential equation
	\mar{Right-hand side}
(which is the expression $.015\,Q^{1.2}$), we would expect to substitute for $Q$ its formula in terms of $t$.  But since in evaluating the left-hand side, we expressed things in terms of $u$,  let's do the same thing here.  Since $Q = u^{-5}$,
	$$
	Q^{1.2} = Q^{6/5} = (u^{-5})^{6/5} = u^{-5 \cdot 6/5}
		= u^{-6}.
	$$
Therefore, the right-hand side is equal to $.015\,u^{-6}$.  But so is the left-hand side, so $Q(t)$ is indeed a solution to the differential equation
	$$
	\dfrac{dQ}{dt} = .015\,Q^{1.2}.
	$$

\aside{Notice two things about this result.  First, when we work with
  formulas we have greater need for algebra to manipulate them.  For
  example, we needed one of the laws of exponents, $(a^b)^c = a^{bc}$,
  to evaluate the right-hand side.  Second, we found it simpler to
  express $Q$ in terms of the intermediate variable $u$, instead of
  the original input variable $t$.  In another computation, it might
  be preferable to replace $u$ by its formula in terms of $t$.  You
  need to choose your algebraic strategy to fit the circumstances.}

\subsubsection{Behavior of the supergrowth solution}

	It was convenient to use a negative exponent in the formula for $Q(t)$ when we wanted to differentiate $Q$.  However, to understand what the formula tells us about supergrowth, it will be more useful to write $Q$ as
	$$
	Q(t)= \left( \frac{1}{1/\sqrt[5]{5}-.003\,t} \right)^5
		\!\!.
	$$
This way makes it clear that $Q$ is a fraction, and we can see its denominator.  In particular, this fraction is not defined when the denominator is zero---that is, when 
	$$
	\frac{1}{\sqrt[5]{5}} - .003\,t = 0, \quad \mbox{or}
		\quad t = \frac{1/\sqrt[5]{5}}{.003} 
		= 241.6\ldots \quad  \mbox{years after 1990}.
	$$
Consider what happens, though, as $t$ approaches this special value 241.6\ldots .  The 
	\mar{$Q$ ``blows up'' as $t \to 241.6\ldots$}
denominator isn't {\em yet\/} zero, but it is {\em approaching\/} zero, so the fraction $Q$ is becoming infinite.   This means that the supergrowth model predicts the world population will become infinite in about 240 years!

	Let's see what the predicted population size is when $t = 240$ (which is the year 2230 {\sc a.d.}), shortly before $Q$ becomes infinite.  We have 
	$$
	Q(240) = \left( \frac{1}{\sqrt[5]{5}} - (.003)(240)
		\right)^{-5} \approx 4.0088 \times 10^{11}.
	$$
Remember that $Q$ expresses the population in {\em billions\/} of people, so the supergrowth model predicts about $4 \times 10^{20}$ people (i.e.\  400 quintillion!) in the year 2230.    Refer back to our estimates of $Q(240)$ using Euler's method (page~\pageref{Q(240)}). Although not even one digit of the estimates had stabilized, at least the final one (with a step size of .001) had reached the right power of ten.  In fact, estimates made with still smaller step sizes do eventually approach the value given by the formula for $Q$:
	\addtolength{\arraycolsep}{3pt}
	$$
	\begin{array}{cc}
	\mbox{step size} & Q(240) \\ [2pt] \hline
	\rule{0pt}{16pt}  .1\hphantom{00000} \rule{0pt}{16pt} 
		& 0.1979 \times 10^{11}  \\
	.01\hphantom{0000} & 2.5727 \times 10^{11} \\
	.001\hphantom{000} & 2.8249 \times 10^{11} \\
	.0001\hphantom{00} & 3.9999 \times 10^{11} \\
	.00001\hphantom{0} & 4.0069 \times 10^{11} 
	\end{array}
	$$
	\addtolength{\arraycolsep}{-3pt}

	Let's look at the relationship between the Euler approximations of $Q$ and the formula for $Q$ graphically.  Here are graphs produced by a modification of the program SEQUENCE.

\vspace*{2.6in}
\special{psfile=../ps/calc1/pic204a.ps}

\noindent
The range of values of $Q$ for $0 \leq t \leq 240$ is so immense that
these graphs are useless.  In a case like this, it is helpful to
rescale the vertical axis so that the space between one power of 10
and the next is the same.  In other words, instead of seeing 1, 2,
3, \ldots , we see $10^1$, $10^2$, $10^3$, \ldots .  This is called a
{\bf logarithmic} scale.\index{logarithmic scale} Here's what happens
to the graphs if we put a logarithmic scale on the vertical axis:

%\newpage

%\mbox{}

\vspace*{2.45in}
\special{psfile=../ps/calc1/pic204b.ps}

\begin{center}
Euler approximations and the formula for $Q$
\end{center}


\noindent
The second graph makes it clearer that the Euler approximations do
indeed approach the graph of the function given by our formula, but
they approach more and more slowly, the closer $t$ approaches
241.6\ldots .

\newpage

\aside{Graphs are made with logarithmic scales particularly when the
  numbers being plotted cover a wide range of values.  When just one
  axis is logarithmic, the result is called a {\em semi-log\/} plot;
  when both axes are logarithmic, the result is called a {\em
    log--log\/} plot.}\index{log--log plot}\index{semi-log plot}

Since $Q(t)$ becomes infinite when $t = 241.6\ldots$ , we must
conclude that the solution to the original initial value problem
%
\mar{The initial value problem has a solution only for $t < 241.6\ldots$}
%
is meaningful only for $t < 241.6\ldots$ .  Of course, the formula for
$Q$ works quite well when $t > 241.6\ldots$ .  It just has no meaning
as the size of a population.  For instance, when $t = 260$ we get
	$$
	Q(260) = \left( \frac{1}{\sqrt[5]{5}} 
			- (.003)(260) \right)^{-5} 
		\approx (-.05522)^{-5} 
		\approx -1.948 \times 10^6.
	$$
In other words, the function determined by the initial value problem is defined only on intervals around $t=0$ that do not contain $t = 241.6\ldots$~.

The formula for $Q'(t)$ is informative too:
\[
Q'(t) = .015 \left( \frac{1}{1/\sqrt[5]{5}-.003t} \right)^6
\]
Since $Q'$ has the same denominator as $Q$, 
%
\mar{$Q'$ blows up as $t \to 241.6\ldots$}
%
it becomes infinite the same way $Q$ does: $Q'(t) \to \infty$ as $t
\to 241.6\ldots$ .  Because Euler's method uses the microscope
equation $\Delta Q \approx Q' \cdot \Delta t$ to predict the next
value of $Q$, we can now understand why the estimates of $Q(240)$ were
so slow to stabilize: as $Q' \to \infty$, $\Delta Q \to \infty$, too.


\subsubsection{A formula for constant per capita growth}

The constant per capita growth model for the world population that we
are considering is
\[
\dfrac{dP}{dt} = .02\, P \qquad P(0) = 5.
\]
This differential equation has a very simple form; if $P(t)$ is a
solution, then the derivative of $P$ is just a multiple of $P$.  We
have already seen in chapter~3 that {\em exponential\/}
functions\index{function!exponential} behave this way (exercises 5--7
in section~3).  For example, if $P(t) = 2^t$, then
\[
\dfrac{dP}{dt} = .69 \cdot 2^t = .69 \, P.
\]
Of course, the multiplier that appears here is .69, not .02, so $P(t)
= 2^t$ is not a solution to our problem.

\pagebreak

However, the multiplier that appears 
%
\mar{Exponential functions satisfy $dy/dt = ky$}
%
when we differentiate an exponential function changes when we change
the base.  That is, if $P(t) = b^t$, then $P'(t) = k_b \cdot b^t$,
where $k_b$ depends on $b$.  Here is a sample of values of $k_b$ for
different bases $b$: \addtolength{\arraycolsep}{4pt}
	$$
	\begin{array}{cc}
	b & k_b \\[2pt] \hline
	\begin{array}{r}
	\rule{0pt}{14pt} .5 \\ 2 \\ 3 \\ 10 
	\end{array}
	&
	\begin{array}{r}
	\rule{0pt}{14pt} -.693147\ldots \\ .693147\ldots \\
		1.098612\ldots \\ 2.302585\ldots 
	\end{array}
	\end{array}
	$$
	\addtolength{\arraycolsep}{-4pt}%
Notice that $k_b$ gets larger as $b$ does.  Since .02\label{k_b} lies
between $-.693147$ and $+.693147$, the table suggests that the value
of $b$ we want lies somewhere between .5 and 2.

	We can say even more about the multiplier.  Since $P'(t) = k_b \cdot P(t)$ and $P(t) = b^t$, we find
	$$
	P'(0) = k_b \cdot P(0) = k_b \cdot b^0 = k_b \cdot 1
		= k_b.
	$$
In other words, {\em $k_b$ is the slope of the graph of $P(t) = b^t$ at the origin}.
\smallskip

Thus, we will be able to solve the differential equation $dP/dt = .02\,P$ if we can find an exponential function $P(t) = b^t$ 
%
\mar{The correct exponential function has slope .02\\ at the origin}
%
whose graph has slope .02 at the origin.  This is a problem that we
can solve with a computer microscope.  Pick a value of $b$ and graph
$b^t$.  Zoom in on the graph at the origin and measure the slope.  If
the slope is more than .02, choose a smaller value for $b$; if the
slope is less than .02, choose a larger value for $b$.  Repeat this
process, narrowing down the possibilities for $b$ until the slope is
as close to .02 as you wish.  Eventually, we get
	$$
	P(t) = (1.0202)^t.
	$$
You should check that $P'(0) = .02000\ldots$ ; see the exercises.  

Thus $P(t) = (1.0202)^t$ solves the differential equation $P' =
.02\,P$.  However, it does not satisfy the initial condition, because
	$$
	P(0) = (1.0202)^0 = 1 \neq 5.
	$$
This is easy to fix; $P(t) = 5 \cdot (1.0202)^t$ satisfies both
conditions.  More generally, $P(t) = C \,(1.0202)^t$ satisfies the
initial condition $P(0) = C$ as well as the differential equation
$P'(t) = .02 \,P(t)$.  To check the initial condition, we compute
	$$
	P(0) = C \, (1.0202)^0 =C \cdot 1 = C.
	$$
The differential equation is also satisfied:
	$$
	P'(t)= (C \, (1.0202)^t)' = C \cdot ((1.0202)^t)'= 
		C \cdot (.02 \, (1.0202)^t) = .02 \, P(t).
	$$
So we have verified that the solution to our problem is 
	\mar{\vspace{16pt} The formula for $P$}
	$$
	P(t) = 5 \, (1.0202)^t.
	$$

Because exponential functions are involved, constant per capita growth
is commonly called {\bf exponential growth}.\index{growth!exponential}
In the figure below we compare exponential growth $P(t)$ to
``supergrowth'' $Q(t)$.  The two graphs agree quite well when $t <
50$.  Notice that population is plotted on a logarithmic scale (a {\em
  semi--log\/} plot).  This makes the graph of $P$ a straight line!


\vspace*{2.4in}

\special{psfile=../ps/calc1/pic207.ps}

\begin{center}
	The graphs of $P(t)$ and $Q(t)$
\end{center}


\subsection{Differential Equations Involving Parameters}

\label{DEs w/ params-begin}

The $S$-$I$-$R$ model contained two parameters---the transmission 
%
\mar{How do parameters affect solutions?}
%
and recovery coefficients $a$ and $b$.\index{parameter} When we used
Euler's method to analyze $S$, $I$, and $R$, we were working
numerically.  To do the computations, we had to give the parameters
definite numerical values.  That made it more difficult to deal with
our questions about the effects of {\em changing\/} the parameters.
As a result, we took other approaches to explore those questions.  For
example, we used algebra to see that there was a threshold for the
spread of the disease: if there were fewer than $b/a$ people in the
susceptible population, the infection would fade away.

This is the situation generally.  Euler's method can be used to
produce solutions to a very broad range of initial value problems.
However, if the model includes parameters, then we usually want to
know how the solutions are affected when the parameters change.
Euler's method is a rather clumsy tool for investigating this
question.  Other methods---ones that don't require the values of the
parameters to be fixed---work better.  One possibility is to start
with a formula.
\medskip

The supergrowth problem illustrates \index{supergrowth}
%
\mar{Supergrowth parameters}
%
both how questions about parameters can arise and how useful a formula
for the solution can be to answer the questions.  One of the most
striking features of the supergrowth model is that it predicts the
population becomes infinite in 241.6\ldots years.  That prediction was
based on an initial population of 5 billion and a growth constant of
.015.  Suppose those values turn out to be incorrect, and we need to
start with different values.  Will that change the prediction?  If so,
how?

We should treat the initial population and the growth constant as
parameters---that is, as quantities that {\em can\/} vary, although
they will have fixed values in any specific situation that we
consider.  Suppose we let $A$ denote the size of the initial
population, and $k$ the growth constant.  If we incorporate these
parameters into the supergrowth model, the initial value problem takes
this form:
	$$
	\dfrac{dQ}{dt} = k\, Q^{1.2} \qquad Q(0) = A
	$$


Here is the formula for a function that solves this 
	\mar{\vspace{15pt} The supergrowth solution with parameters}
problem:\label{supergrowth soln}
	$$
	Q(t) =\left( \frac{1}{\sqrt[5]{A}} -.2kt \right)^{-5}
	$$
Notice that, when $A = 5$ and $k = .015$, this formula reduces to the one we considered earlier.

	Let's check 
	\mar{Checking the formula}
that the formula does indeed solve the initial value problem.  First, the initial condition:
	$$
	Q(0)=\left( \frac{1}{\sqrt[5]{A}} -.2k \cdot 0
		\right)^{-5} =
		\left(\frac{1}{\sqrt[5]{A}}\right)^{-5} =
		(\sqrt[5]{A})^5 = A.
	$$
Next, the differential equation.  To differentiate $Q(t)$ we introduce the chain
	$$
	Q = u^{-5} \qquad \mbox{where} \qquad 
		u = \frac{1}{\sqrt[5]{A}} -.2kt.
	$$
We see that $dQ/du = -5\,u^{-6}$.  Since $u$ is a linear function of $t$ in which the multiplier is $-.2k$, we also have $du/dt = -.2k$.  Thus, by the chain rule,
	$$
	\dfrac{dQ}{dt} = \dfrac{dQ}{du} \cdot \dfrac{du}{dt} =
		-5\,u^{-6} \cdot (-.2k) = k\,u^{-6}.
	$$
That is the left-hand side of the differential equation.  To evaluate the right-hand side, we use the fact that $Q = 
u^{-5}$.  Thus
	$$
	k\,Q^{1.2} = k\,Q^{6/5} = k(u^{-5})^{6/5} = 
		k\,u^{-5 \cdot 6/5} = k\,u^{-6}.
	$$
Since both sides equal $k\,u^{-6}$, they equal each other, proving that $Q(t)$ is a solution to the differential equation.

Next, we ask when the population becomes infinite.  Exactly as before,
this will happen when the denominator of the formula for $Q(t)$
becomes
%
\mar{\vspace{18pt} The time to infinity}
%
zero:
	$$
	\dfrac{1}{\sqrt[5]{A}} - .2\,kt = 0, \qquad \mbox{or}
		\qquad t = \dfrac{1}{.2\,k\sqrt[5]{A}}.
	$$
Here, in fact, is a {\em formula\/} that tells us how each of the parameters $A$ and $k$ affects the time it takes for the population to become infinite.

	Let's use $\tau$ (the Greek letter ``tau'') to denote the ``time to infinity.''  For example, if we double the initial population, so $A = 10$ billion people, while keeping the original growth constant $k = .015$, then the time to infinity is
	$$
	\tau = \dfrac{1}{.003 \times \sqrt[5]{10}} \approx
		 210.3 \ \mbox{years}.
	$$
By contrast, if we double the growth rate, to $k = .030$, while keeping the original $A = 5$, then the time to infinity is only
	$$
	\tau = \dfrac{1}{.006 \times \sqrt[5]{5}} \approx
		120.8 \ \mbox{years}
	$$
Conclusion: doubling the growth rate has a much greater impact than doubling the initial population.  

\label{supergrowth errors-begin}
For any specific growth rate and initial population, 
%
\mar{Uncertainty in\\ the size of $\tau$}
%
we can always calculate the time to infinity.  But we can actually do
more; the formula for $\tau$ allows us to do an {\em error
  analysis\/}\index{error analysis} along the patterns described in
chapter~3, section~4.  For example, suppose we are uncertain of our
value of the growth rate $k$; there may be an error of size $\Delta
k$.  How uncertain does that make us about the calculated value of
$\tau$?  Likewise, if the current world population $A$ is known only
with an error of $\Delta A$, how uncertain does {\em that\/} make
$\tau$?  Also, how are the {\em relative\/} errors related?  Let's do
this analysis, assuming that $k = .015$ and $A = 5$.

	Our tool is the error propagation equation---which is 
	\mar{How an error in $k$ propagates}
the microscope equation.\index{error propagation}  If we deal with $k$ first, then
	$$
	\Delta \tau \approx \dfrac{\partial \tau}{\partial k}
		\cdot \Delta k.
	$$
We have used partial derivatives because $\tau$ is a function of {\em two\/} variables, $A$ as well as $k$.  If we write
	$$
	\tau = \dfrac{1}{.2\sqrt[5]{A}} \, k^{-1},
	$$
then the differentiation rules yield
	$$
	\dfrac{\partial \tau}{\partial k} = -1 \cdot 
		\dfrac{1}{.2\sqrt[5]{A}} \, k^{-2} = 
		\dfrac{-1}{.2\sqrt[5]{5}} \times (.015)^{-2} 
		\approx -16106.
	$$
Thus $\Delta \tau \approx -16106 \cdot \Delta k$.  For example, if the uncertainty in the value of $k = .015$ is $\Delta k = \pm .001$, then the uncertainty in $\tau$ is about $\mp 16$ years.

	To determine how an error in $A$ propagates to $\tau$, 
	\mar{How an error in $A$ propagates}
we first write
	$$
	\tau = \dfrac{1}{.2\,k} A^{-1/5}.
	$$
Then
	$$
	\dfrac{\partial \tau}{\partial A} = 
		-\dfrac{1}{5} \cdot \dfrac{1}{.2\,k} A^{-6/5} = 
		\dfrac{-1}{5 \times.2 \times .015} \times 5^{-6/5} 
		\approx -9.7.
	$$
The error propagation equation is thus $\Delta \tau \approx -9.7 \cdot \Delta A$.  If the uncertainty in the world population is about 100 million persons, so $\Delta A = \pm .1$, then the uncertainty in $\tau$ is less than 1 year.

	To complete the analysis, 
	\mar{Relative errors}
let's compare relative errors.  This involves a lot of algebra.  To see how an error in $k$ propagates, we have
	$$
	\Delta \tau \approx -\dfrac{\Delta k}{.2\,k^2\sqrt[5]{A}}
		\qquad \mbox{and} \qquad 
		\tau = \dfrac{1}{.2\,k\sqrt[5]{A}}.
	$$
We can therefore compute that a given relative error in $k$ propagates as 
	$$
	\dfrac{\Delta \tau}{\tau} \approx 
		-\dfrac{\Delta k}{.2\,k^2\sqrt[5]{A}} \cdot
			\dfrac{.2\,k\sqrt[5]{A}}{1} 
		= -\dfrac{\Delta k}{k}.
	$$
Thus, a 1\% error in $k$ leads to a 1\% error in $\tau$, although the sign is reversed.

	To analyze how a given relative error in $A$ propagates, we start with
	$$
	\Delta \tau \approx -\dfrac{1}{5} \cdot
		\dfrac{\Delta A}{.2\,k A^{6/5}}.
	$$
Then
	$$
	\dfrac{\Delta \tau}{\tau} \approx -\dfrac{1}{5} \cdot
		\dfrac{\Delta A}{.2\,k A^{6/5}} \cdot
			\dfrac{.2\,k\sqrt[5]{A}}{1} =
		-\dfrac{1}{5} \cdot \dfrac{\Delta A}{A}.
	$$
This says that it takes a 5\% error in $A$ 
	\mar{How sensitive $\tau$ is to errors in $k$ and $A$}
to produce a 1\% error in $\tau$.  Consequently, the time to infinity {\em $\tau$ is 5 times more sensitive to errors in $k$ than to errors in $A$}.
\label{supergrowth errors-end}

\bigskip

	The exercises in this section will give you an opportunity to check that a particular formula is a solution to an initial value problem that arises in a variety of contexts.  Later in this chapter, we will make a modest beginning on the much harder task of {\em finding\/} solutions given by formulas for special initial value problems.  
	\mar{Special methods give formulas; general methods are numerical}
There are more sophisticated methods for finding formulas, when the formulas exist, and they provide powerful tools for some important problems, especially in physics.  However, most initial value problems we encounter cannot be solved by formulas.  This is particularly true when two or 
more variables are needed to describe the process being modelled.  The tool of widest applicability is Euler's method.  This isn't so different from the situation in algebra, where exact solutions given by formulas (e.g.\ the quadratic formula) are also relatively rare, and numerical methods play an important role.  (Chapter~5.5, presents the Newton--Raphson method\index{Newton's method} for solving algebraic equations by successive approximation.)  In most cases that will interest us, there are simply no formulas to be found---the limitation lies in the mathematics, not the mathematicians.  
\label{DEs w/ params-end}

\subsection{Exercises}


In exercises 1--4, verify that the given formula is a {\em solution\/} to the initial value problem.\index{initial value problem}

\num {\bf Powers of $y$}.

\sub  $y'=y^2$, $y(0)=5$: \quad $y(t)= 1/(\frac{1}{5} - t)$

\sub  $y'=y^3$, $y(0)=5$:   \quad $y(t) = 1/\sqrt{\frac{1}{25} - 2t}$

\sub  $y'=y^4$, $y(0)=5$:  \quad $y(t) =  1/\sqrt[3]{\frac{1}{125} - 3t}$

\sub  Write a general formula for the solution of the initial value problem $y'=y^n$, $y(0)=5$, for any integer $n > 1$.
\sub  Write a general formula for the solution of the initial value problem $y'=y^n$, $y(0)=C$, for any integer $n > 1$ and any constant $C\geq 0$.


\num  {\bf Powers of $t$}.

\sub  $y'=t^2$, $y(0)=5$:  \quad $y(t) = \frac{1}{3}t^3 + 5$

\sub  $y'=t^3$, $y(0)=5$:  \quad $y(t) = \frac{1}{4}t^4 + 5$

\sub  $y'=t^4$, $y(0)=5$:  \quad $y(t) = \frac{1}{5}t^5 + 5$

\sub  Write a general formula for the solution of the initial value problem $y'=t^n$, $y(0)=5$ for any integer $n > 1$.
\sub  Write a general formula for the solution of the initial value problem $y'=t^n$, $y(0)=C$ for any integer $n > 1$ and any constant $C$.


\num  {\bf Sines and cosines}.

\sub  $x'=-y$, $y'=x$, $x(0)=1$, $y(0)=0$: \quad $x(t)=\cos t$, $y(t) = \sin t$

\sub  $x'=-y$, $y'=x$, $x(0)=0$, $y(0)=1$: \quad $x(t) = \cos(t+\pi/2)$, $y(t) = \sin(t+\pi/2)$

\num  {\bf Exponential functions}.

\sub  $y'=2.3 \, y$, $y(0)=5$: \quad $y(t)=5 \cdot 10^t$

\sub  $y'=2.3 \, y$, $y(0)=C$: \quad $y(t)=C \cdot 10^t$

\sub  $y'=-2.3 \, y$, $y(0)=5$: \quad $y(t)=5 \cdot 10^{-t}$

\sub  $y'=4.6 \,ty$, $y(0)=5$: \quad $y(t)= 5 \cdot 10^{t^2}$

\num  {\bf Initial Conditions}.

\sub  Choose $C$ so that $y(t) = \sqrt{t+C}$ is a solution to the initial value problem
	$$
	y'=\frac{1}{2y} \quad y(3)=17.
	$$

\sub  Choose $C$ so that $  y(t) = -1/(t+C)$  is a solution to the initial value problem
	$$
	y'=y^2 \quad y(0)=-5.
	$$

\sub  Choose $C$ so that $  y(t) = -1/(t+C)$  is a solution to the initial value problem
	$$
	y'=y^2 \quad y(2)=3.
	$$

\subsubsection{World population growth with parameters}
\index{population!world}
\index{parameter}

\vspace{-10pt}
\numa  Using a graphing utility or a calculator, show that the derivative of $P(t)=(1.0202)^t$ at the origin is approximately .02: $P'(0) \approx .02$.  Since quick convergence is desirable, use
	$$
	\dfrac{\Delta P}{\Delta t} = \dfrac{P(0+h) - P(0-h)}{2h}
		= \dfrac{(1.0202)^h - (1.02020)^{-h}}{2h}
	$$

\sub By using more decimal places to get higher precision, show that $P(t) = (1.0202013)^t$ satisfies $P'(0) = .02$ even more exactly.

\numa  \label{ex:base 2} Show that the function $y = 2^{t/.69}$ satisfies the differential equation $dy/dt = y$.  Use the chain rule: $y = 2^u$, $u = t/.69$. (Recall that $k_2=.69\ldots\,.$)

\sub  Show that the function $y = 2^{kt/.69}$ satisfies the differential equation $dy/dt = k\,y$.

\sub  Show that the function $P(t) = A \cdot 2^{kt/.69}$ is a solution to the initial value problem
	$$
	\dfrac{dP}{dt} = k\,P \qquad P(0) = A.
	$$
Note that this describes a population that grows at the constant per capita rate $k$ from an initial size of $A$.


\numa  Show that the function $y = 10^{t/2.3}$ satisfies the differential equation $dy/dt = y$.  Use the chain rule: $y = 10^u$, $u = t/2.3$. (Recall that $k_{10}=2.3\ldots\,.$)

\sub  Show that the function $y = 10^{kt/2.3}$ satisfies the differential equation $dy/dt = k\,y$.

\sub  Show that the function $P(t) = A \cdot 10^{kt/2.3}$ is a solution to the initial value problem
	$$
	\dfrac{dP}{dt} = k\,P \qquad P(0) = A.
	$$
This formula provides an alternative way to describe a population that grows at the constant per capita rate $k$ from an initial size of $A$.


\numa  The formula $P(t) = 5\cdot 2^{k\,t/.69}$ describes how an initial population of 5 billion will grow at a constant per capita rate of $k$ persons per year per person.  Use this formula to determine how many years $t$ it will take for the population to double, to 10 billion persons.  

\sub  Suppose the initial population is $A$ billion, instead of 5 billion.  What is the doubling time then?

\sub  Suppose the initial population is 5 billion, and the per capita growth rate is .02, but that value is certain only with an error of $\Delta k$.  How much uncertainty is there in the doubling time that you found in part (a)?


\subsubsection{Newton's law of cooling}
\index{Newton's law of cooling}

There are formulas that describe how a body cools, or heats up, to
match the temperature of its surroundings.  See the exercises on
Newton's law of cooling in section~1.  Consider first the model
	$$
	\dfrac{dT}{dt} = -.1(T - 20) \qquad T(0) = 90,
	$$
introduced on page~\pageref{cooling coffee} to describe how a cup of
coffee cools.

\num Show that the function $y = 2^{-.1\,t/.69}$ is a solution to the
differential equation $dy/dt = -.1\,y$.  (Use the chain rule: $y =
2^u$, $u = -.1\,t/.69$.)

\numa  Show that the function 
	$$
	T = 70 \cdot 2^{-.1\,t/.69} + 20 
	$$
is a solution to the initial value problem $dT/dt = -.1(T - 20)$,
$T(0) = 90$.  This is the temperature $T$ of a cup of coffee,
initially at 90$^{\circ}$C, after $t$ minutes have passed in a room
whose temperature is 20$^{\circ}$C.

\sub Use the formula in part (a) to find the temperature of the coffee
after 20 minutes.  Compare this result with the value you found in
exercise~12~(b), page~\pageref{ex:coffee1}.

\sub Use the formula in part (a) to determine how many minutes it
takes for the coffee to cool to 30$^{\circ}$C.  In doing the
calculations you will find it helpful to know that $1/7 = 2^{-2.8}$.
Compare this result with the value you found in exercise~11~(c),
page~\pageref{ex:coffee2}.

\numa A cold drink is initially at $ Q = 36^{\circ}$F when the air
temperature is 90$^{\circ}$F.  If the temperature changes according to
the differential equation
	$$
	\dfrac{dQ}{dt} = -.2 (Q - 90) ^{\circ}
		\mbox{F per minute},
	$$
show that the function $Q(t) = 90 - 54\cdot2^{-.2\,t/.69}$ describes
the temperature after $t$ minutes.

\sub Use the formula to find the temperature of the drink after 5
minutes and after 10 minutes.  Compare your results with the values
you found in exercise~12, page~\pageref{ex:cold drink}.


\num  Find a formula for a function that solves the initial value problem
	$$
	\dfrac{dQ}{dt} = -k(Q - A) \qquad Q(0) = B.
	$$


\subsubsection{A leaking tank}
\index{leaking tank}\label{leaking tank-begin}

	The rate at which water leaks from a small hole at the bottom of a tank is proportional to the square root of the height of the water surface above the bottom of the tank.  Consider a cylindrical tank that is 10 feet tall and stands on one of its circular ends, which is 3 feet in diameter. Suppose the tank is currently half full, and is leaking at a rate of 2 cubic feet per hour.

\numa  Let $V(t)$ be the volume of water in the tank $t$ hours from now.  Explain why the leakage rate can be written as the differential equation
	$$
	V'(t) = -k\sqrt{V(t)},
	$$
for some positive constant $k$.  (The issue to deal with is this: why is it permissible to use the square root of the {\em volume\/} here, when the rate is known to depend on the square root of the {\em height\/}?)

\sub  Determine the value of $k$.  [Answer: $k \approx .3364$; you need to explain {\em why\/} this is the value.]

\numa  How much water leaks out of the tank in 12 hours; in 24 hours?  Use Euler's method, and compute successive approximations until your results stabilize.

\sub  How many hours does it take for the tank to empty?

\numa  Use the differentiation rules to show that any function of the form
	$$
	V(t) = \left\{ 
		\begin{array}{ll}
		\dfrac {{k^2}}{4} (C - t)^2 & 
			\mbox{if }\ 0 \leq t \leq C \\[.13in]
		0 &  \mbox{if }\ C < t
		\end{array}
	\right.
	$$
satisfies the differential equation.

\sub  For the situation we are considering, what is the value of $C$?  According to this solution, how long does it take for the tank to empty?  Compare this result with your answer using Euler's method.

\sub Sketch the graph of $V(t)$ for $0 \leq t \leq 2C$, taking particular care to display the value $V(0)$ in terms of $k$ and $C$.  
\label{leaking tank-end}

\subsubsection{Motion}
\index{motion}

Newton created the calculus to study the motion of the planets.  He
said that all motion obeys certain basic laws.  One law says that the
velocity of an object changes only if a force acts on the body.
Furthermore, the {\em rate\/} at which the velocity changes is
proportional to the force.  By knowing the forces that act on a body
we can construct---and then solve---a differential equation for the
velocity.
\medskip

\noindent
{\bf Falling bodies---with gravity}.\index{falling object} A body
falling through the air starts up slowly but picks up speed as it
falls.  Its velocity is thus changing, so there must be a force
acting.  We call the force that pulls objects to the earth {\bf
  gravity}.\index{gravity} At the surface of the earth, the rate
of change of velocity caused by gravity is essentially the same for
all objects.

Suppose an object is $x$ meters above the surface of the earth after
$t$ seconds have passed.  Then, by definition, its velocity is
	$$
	v = \dfrac{dx}{dt} \ \mbox{meters/second}.
	$$
According to Newton's laws of motion, the force of gravity causes the velocity to change, and we can write
	$$
	\dfrac{dv}{dt} = -g.
	$$
Here $g$ is a constant whose numerical value is about 9.8 meters/second per second.  Since $x$ and $v$ are positive when measured {\em upwards}, but gravity acts {\em downwards}, a minus sign is needed in the equation for $dv/dt$.  (The derivative of velocity is commonly called {\bf acceleration}, and $g$ is called the {\bf acceleration due to gravity}.)\index{acceleration!due to gravity}

\num Verify that $v(t) = -g\,t + v_0$\label{falling body0} is a
solution to the differential equation $dv/dt = -g$ with initial
velocity $v_0$.

\num  Since $dx/dt = v$, and since $v(t) = -g\,t + v_0$, the {\em position\/} $x$ of the body satisfies the differential equation
	$$
	\dfrac{dx}{dt} = -g\,t + v_0 \ \mbox{meters/second}.
	$$
Find a formula for $x(t)$ that solves this differential equation.  This function describes how a body moves under the force of gravity.

\num  Suppose the initial position of the body is $x_0$, so that the position $x$ is a solution to the initial value problem
	$$
	\dfrac{dx}{dt} = -g\,t + v_0 \qquad x(0) = x_0 \ 
		\mbox{meters}.
	$$
Find a formula for $x(t)$.

\numa  Suppose a body is held motionless 200 meters above the ground, and then released.  What values do $x_0$ and $v_0$ have?  What is the formula for the motion of this body as it falls to the ground?

\sub  How far has the body fallen in 1 second?  In 2 seconds?

\sub  How long does it take for the body to reach the ground?
\medskip

\label{fall w/air-begin}
\noindent
{\bf Falling bodies---with gravity and air resistance}.\index{air
  resistance} As a body falls, air pushes against it.  Air resistance
is thus another force acting on a falling body.  Since air resistance
is slight when an object moves slowly but increases as the object
speeds up, the simplest model we can make is that the force of air
resistance is proportional to the velocity: force = $-bv$ (reality
turns out to be somewhat more complicated than this).  The multiplier
$b$ is positive, and the minus sign tells us that the direction of the
force is always opposite the velocity.  The forces of gravity and air
resistance combine to change the velocity:
\[
\dfrac{dv}{dt} = -g - b\,v \ \mbox{meters/second per second}.
\]

\num\label{ex:falling body1}  Show that 
	$$
	v(t) = \dfrac{g}{b}\left(2^{-b\,t/.69} - 1\right) 
		\ \mbox{meters/second}
	$$
is a solution to this differential equation that also satisfies the initial condition $v(0) = 0$ meters/second.


\numa\label{ex:falling body2} Show that the position $x(t)$ of a body
that falls against air resistance from an initial height of $x_0$
meters is given by the formula
	$$
	x(t) = x_0 - \dfrac{g}{b} \, t - \dfrac{g}{b^2}
		\left(2^{-b\,t/.69} - 1\right) \ \mbox{meters}.
	$$

\sub\label{ex:falling body3} Suppose the coefficient of air resistance
is $b = .2$ per second.  If a body is held motionless 200 meters above
the ground, and then released, how far will it fall in 1 second?  In 2
seconds?  Compare these values with those you obtained assuming there
was no air resistance.

\sub  How long does it take for the body to reach the ground?  (Use a computer graphing package to get this answer.)  Compare this value with the one you obtained assuming these was no air resistance.  How much does air resistance add to the time?


\numa  According to the equation $dv/dt = -g - b\,v$, there is a velocity $v_T$ at which the force of air resistance exactly balances the force of gravity, and the velocity doesn't change.  What is $v_T$, expressed as a function of $g$ and $b$.  Note: $v_T$ is called the {\bf terminal velocity}\index{terminal velocity} of the body.  Once the body reaches its terminal velocity, it continues to fall at that velocity.

\sub  What is the terminal velocity of the body in the previous exercise?
\label{fall w/air-end}

\medskip

\noindent
{\bf The oscillating spring}.\index{oscillating
  spring}\label{spring-begin} Springs can smooth out life's little
irregularities (as in the suspension of a car) or amplify and measure
them (as \linebreak

\vspace{-10pt}
\noindent
\begin{minipage}[b]{3.4in}
in earthquake detection devices).  Suppose a spring that hangs from a
hook has a weight at its end.  Let the weight come to rest.  Then,
when the weight moves, let $x$ denote the position of the weight above
the rest position.  (If $x$ is negative, this means the weight is
below the rest position.)  If you pull down on the weight, the spring
pulls it back up.  If you push up on the weight, the spring (and
gravity) push it back down.  This push is the {\bf spring
  force}.\index{spring force}
\end{minipage}

\special{psfile=../ps/calc1/spring9.ps}

The simplest assumption is that the spring force is proportional to
the amount $x$ that the spring has been stretched: force = $-c^2x$.
The constant $c^2$ is customarily written as a square to emphasize
that it is positive.  The minus sign tells us the force pushes down if
$x > 0$ (so the weight is above the rest position), but it pushes up
if $x < 0$.

	If $v = dx/dt$ is the velocity of the weight, then Newton's law of motion says
	$$
	\dfrac{dv}{dt} = - c^2 x.
	$$
Suppose we move the weight to the point $x = a$ on the scale, hold it motionless momentarily, and then release it at time $t = 0$.  This determines the initial value problem
	$$
	\begin{array}{ll}
	x' = v & x(0) = a \\
	v' = -c^2 x \qquad & v(0) = 0.
	\end{array}
	$$

\numa  Show that
	$$
	x(t) = a \cos(ct) \qquad v(t) = -ac \sin(ct)
	$$
is a solution to the initial value problem.

\sub What range of values does $x$ take on; that is, how far does the
weight move from its rest position?

\numa Use a graphing utility to compare the graphs of
$y=\cos(x),\ y=\cos(2x),\ y=\cos(3x)$, and $y = \cos(.5x)$.  Based on
your observations, explain how the value of $c$ affects the nature of
the motion $x(t)=a \cos(ct)$ for a fixed value of $a$.

\sub  How long does it take the weight to complete one cycle (from $x = a$ back to $x = a$) when $c = 1$?  The motion of the weight is said to be {\bf periodic}, and the time it takes to complete one cycle is called its {\bf period}.\index{period!of an oscillation}  

\sub What is the period of the motion when $c = 2$?  When $c = 3$?
Does the period depend on the initial position $a$?

\sub Write a formula that expresses the period of the motion in terms
of the parameters $a$ and $c$.


\numa The parameter $c$ depends on two things: the mass $m$ of the
weight, and the stiffness $k$ of the spring:
	$$
	c = \sqrt{\dfrac{k}{m}}. 
	$$
Write a formula that expresses the period of the motion of the weight in terms of $m$ and $k$.

\sub Suppose you double the weight on the spring.  Does that {\em
  increase\/} or {\em decrease\/} the period of the motion?  Does your
answer agree with your intuitions?

\sub Suppose you put the first weight on a second spring that is twice
as stiff as the first (i.e., double the value of $k$).  Does that {\em
  increase\/} or {\em decrease\/} the period of the motion?  Does your
answer agree with your intuitions?

\sub When you calculate the period of the motion using your formula
form part (a), suppose you know the actual value of the mass only to
within 5\%.  How accurately do you know the period---as a percentage
of the calculated value?
\label{spring-end}

%%%%%%%%%%%%%%%%%%%%%%%%%%%%%%%%%%%%%%%%%%%%%%%%%%%%%%%%%%%%%%%%%%%%%%%%%%%
%                                                                         %
%                                Section 3                                %
%                                                                         %
%%%%%%%%%%%%%%%%%%%%%%%%%%%%%%%%%%%%%%%%%%%%%%%%%%%%%%%%%%%%%%%%%%%%%%%%%%%

\section{The Exponential Function}

\subsection{The Equation $y'= ky$}

As we have seen, initial value problems {\em define\/} functions---as
their solutions.  They therefore provide us with a vast, if somewhat
bewildering, array of new functions.  Fortunately, a few differential
equations---in fact, the very simplest---arise over and over again in
an astonishing variety of contexts.  The functions they define are
among the most important in mathematics.

One of the simplest differential equations is 
%
\mar{A simple and\\ natural model of\\ growth and decay}
%
$dy/dt =ky$, where $k$ is a constant.  It is also one of the most
useful.  We used it in chapter 1 to model the populations of Poland
and Afghanistan, as well as bacterial growth and radioactive decay.
In this chapter, it was our initial model of a rabbit population and
one of our models of the world population.  Later, we will use it to
describe how money accrues interest in a bank and how radiation
penetrates solid objects.

In this section we will look at the solutions to differential
equations of this form from two different vantage points.  On the one
hand, we already have named functions which solve such equations---the
exponential functions.  On the other hand, the fact that Euler's
method produces the same functions will allow us to prove properties
of such functions and to compute their values effectively.

In chapter 3 we established that the solutions 
%
\mar{Exponential solutions---with different bases~\ldots}
%
to $dy/dt = ky$ are exponential functions.\index{function!exponential}
Specifically, for each base $b$, the exponential function $y = b^t$
was a solution to $dy/dt = k_b \cdot b^t = k_b \cdot y$, where $k_b$
was the slope of the graph of $y = b^t$ at the origin.  In this
approach, if the constant $k$ changes, we must change the base $b$ so
that $k_b = k$.


Exercise~7 in the previous section (page~\pageref{ex:base 2}) 
%
\mar{\ldots and a\\ fixed base} 
%
opened up a new possibility: for the {\em fixed\/} base 2, the
function
	$$
	y = 2^{kt/.6931\ldots}
	$$
was a solution to the differential equation $dy/dt = kt$, no matter what value $k$ took.  There was nothing special about the base 2, of course.  In the next exercise, we saw that the functions
	$$
	y = 10^{kt/2.3025\ldots}
	$$
would serve equally well as solutions.  

	In fact, we can show that, for any base $b$, the functions
	$$
	y = b^{kt/k_b}
	$$
are also solutions to $dy/dt = ky$.  Construct the chain
	$$
	y = b^u \qquad \mbox{where} \qquad u = kt/k_b.
	$$
Then $dy/du = k_b\cdot b^u = k_b\cdot y$, while $du/dt = k/k_b$.  Thus, by the chain rule we have
	$$
	\dfrac{dy}{dt} = \dfrac{dy}{du} \cdot \dfrac{du}{dt} =
		k_b\, y \cdot \dfrac{k}{k_b} = k y.
	$$

If we express solutions to $dy/dt = ky$ by 
%
\mar{Advantages of a\\ fixed base}
%
exponential functions with a {\em fixed\/} base, it is easy to alter the solution if the growth constant $k$ changes.  We just change the value of $k$ in the exponent of $b^{kt/k_b}$.  Let's see how this works when $b = 2$ and $b = 10$:
\addtolength{\arraycolsep}{6pt}
	$$
	\begin{array}{ccc}
	\mbox{differential} & \mbox{solution} & 
		\mbox{solution} \\[-2pt]
	\mbox{equation} & \mbox{base 2} & 
		\mbox{base 10}  \\[2pt] \hline
	\dfrac{dy}{dt} = .16 \, y &
		\rule{0pt}{21pt} 2^{.231\,t} \rule{0pt}{21pt} &
		10^{.069\,t} \\ [12pt]
	\dfrac{dy}{dt} = .18 \, y &
		2^{.260\,t} & 10^{.078\,t}
	\end{array}
	$$
Notice that the growth constant $k$ gets ``swallowed up'' in the
exponent of the solution when $k$ has a specific numerical value.  The
number that appears\linebreak

\pagebreak

\noindent
 in the exponent is $k$ divided by $k_2 =
.6931\ldots$ (when the base is 2) and by $k_{10} = 2.3025\ldots$ (when
the base is 10).

The most vivid solution to $dy/dt = ky$ would use the base $b$ for
which $k_b = 1$ {\em exactly}.
%
\mar{The base $e$}
%
There is such a base, and it is always denoted $e$\index{$e$}.  (We
will determine the value of $e$ in a moment.)  Since $k_e = 1$, $k$
would stand out in the exponent:
	$$
	\begin{array}{cc}
	\mbox{differential} & \mbox{solution} \\[-2pt]
	\mbox{equation} & \mbox{base $e$}  \\[2pt] \hline
	\dfrac {dy}{dt} = .16 \, y &
		\rule{0pt}{21pt} e^{.16\,t} 
		\rule{0pt}{21pt} \\ [12pt]
	\dfrac{dy}{dt} = .18 \, y &
		e^{.18\,t} 
	\end{array}
	$$
\addtolength{\arraycolsep}{-6pt}
The simplicity and clarity of this expression have led to the universal adoption of the base $e$\index{$e$} for describing exponential growth and decay---that is, for describing solutions to $dy/dt = ky$.

\aside{The use of the symbol $e$ to denote the base dates back to a
  paper that Euler\index{Euler, L.} wrote at age 21, entitled {\em
    Meditatio in experimenta explosione tormentorum nuper instituta}
  (Meditation upon recent experiments on the firing of cannons), where
  the symbol $e$ was used sixteen times.  It is now in universal use.
  The number $e$ is, like $\pi$, one of the most important and
  ubiquitous in mathematics.}

By design, $y = e^t$ is a solution to the differential equation $dy/dt
= y$.  In particular, the slope of the graph of $y = e^t$ at the
origin is exactly~1.  As we have just seen, the function $y = e^{kt}$
is a solution to the differential equations $dy/dt = ky$ whose growth
constant is $k$.
%
\mar{\vspace{30pt}The general initial value problem for exponential functions}
%
Finally: 
\begin{center}
\begin{picture}(310,60)		
\put(0,0) {\framebox(310,60)
{
\begin{minipage}{300pt}			
\mbox{}\quad $y = C \cdot e^{kt}$ is the solution to the initial value problem
	
	\vspace{-6pt}
	$$
	\dfrac{dy}{dt} = ky \qquad y(0) = C.
	$$
\end{minipage}
}}
\end{picture}
\end{center}
We can check this quickly.  The initial condition is satisfied because $e^0 = 1$, so $y(0) = C \cdot e^{k \cdot 0} = C \cdot 1 = C$.  The differential equation is satisfied because
	$$
	 \left(C \cdot e^{kt} \right)' =
		C \cdot \left(e^{kt} \right)' = 
		C \cdot k\,e^{kt} = k\,y.
	$$
We used the differentiation rule for a constant multiple of a function, and we used the fact that the derivative of $e^{kt}$ was already established to be $k\,e^{kt}$.


\subsection{The Number $e$} 

	The number $e$ is determined by the property that $k_e = 1$.  
%
\mar{Finding $e$ with a computer microscope}
%
Since this number is the slope of the graph of $y = e^t$ at the
origin, one way to find $e$ is with a computer microscope.  Pick an
approximation $E$ for $e$ and graph $y = f^t$.  Zoom in the graph at
the origin, and measure the slope.  If the slope is more than~1,
choose a smaller approximation; if the slope is less than~1, choose a
larger value.  Repeat this process, narrowing down the value of $e$
until you know its value to as many decimal places as you wish.

	We already know $E= 2$ is too small, 
%
\mar{Under successive magnifications: $e = 2.71828\ldots$}
%
because the slope of $y = 2^t$ at the origin is .69.  Likewise, $E = 3$ is too large, because the slope of $y = 3^t$ at the origin is 1.09.  Thus $2 < e < 3$, and is closer to 3 than to 2.  
At the next stage we learn that 2.7 is too small (slope $= .9933$) but 2.8 is too large (slope $= 1.0296$).  Thus, at least we know $e = 2.7\ldots$ .  Several stages later we would learn $e = 2.71828\ldots$ .  
\medskip

	While the method just described for finding the value of $e$ works, it is somewhat ponderous. 
\mar{Finding $e$ by\\ Euler's method}
We can take a very different approach to finding the numerical value of $e$ by using the fact that $e$ is defined by an initial value problem.  Here is the idea: $e$ is the value of the function $e^t$ when $t = 1$, and $y(t) = e^t$ is the solution to the initial value problem 
	$$
	y' = y \qquad y(0) = 1.
	$$
We can then find $e = y(1)$ in the usual way by solving this initial value problem using Euler's method.  Due to some convenient algebraic simplifications, this approach yields powerful insights about the nature of $e$.

	Suppose we take $n$ steps to go from $t = 0$ to $t = 1$.  Then the step size is $\Delta t = 1/n$.  The following table shows the calculations:
	\begin{center}
	Finding $y(1)$ by Euler's method when $y' = y$ and $y(0) = 1$
	\end{center}
	\addtolength{\arraycolsep}{6pt}
	$$
	\begin{array}{cccc}
	t & y & y' = y & \Delta y = y' \cdot \Delta t 
		\\[2pt] \hline
	\rule{0pt}{14pt} 0 \rule{0pt}{14pt} & 1 & 1 & 
		1 \cdot 1/n \\
	1/n & 1 + 1/n & 1 + 1/n & (1 + 1/n) \cdot 1/n \\
	2/n & (1 + 1/n)^2 & (1 + 1/n)^2 & 
		(1 + 1/n)^2 \cdot 1/n \\
	3/n & (1 + 1/n)^3 & (1 + 1/n)^3 & 
		(1 + 1/n)^3 \cdot 1/n \\
	\vdots & \vdots & \vdots & \vdots \\
	n/n & (1 + 1/n)^n 
	\end{array}
	$$
	\addtolength{\arraycolsep}{-6pt}%
The entries in the $y$ column need to be explained.  The first two should be clear: $y(0)$ is the initial value~1, and $y(1/n) = y(0) + \Delta y = 1 + 1/n$.  To get from any entry to the next we must do the following:
	\begin{align*}
	\mbox{new $y$} &= \mbox{current $y$} + \Delta y \\
		&= \mbox{current $y$} + y' \cdot \Delta t \\
		&= \mbox{current $y$} + 
			\mbox{current $y$} \cdot \Delta t \\
		&= \mbox{current $y$} \cdot (1 + \Delta t) \\
		&= \mbox{current $y$} \cdot (1 + 1/n)
	\end{align*}
The new $y$ is the current $y$ multiplied by $(1 + 1/n)$.  Since the second $y$ is itself $(1 + 1/n)$, the third will be $(1 + 1/n)^2$, the fourth will be $(1 + 1/n)^3$, and so on.

	Euler's method with $n$ steps therefore gives us the following estimate for $e = y(1) = y(n/n)$:
	$$
	e \approx (1 + 1/n)^n
	$$
We can calculate these numbers on a computer.  In the following table we give values of $(1+1/n)^n$ for increasing values of $n$.  By the time $n=2^{40}$ (about $10^{12}$), eleven digits of $e$ have stabilized.\index{$e$}

	\addtolength{\arraycolsep}{12pt}
	$$
\begin{array}{cc}
	n & (1 + 1/n)^n \\[2pt] \hline
	\begin{array}{l}
	2^0 \rule{0pt}{14pt} \\ 2^4 \\ 2^8 \\ 2^{12} \\ 2^{16}
		\\ 2^{20} \\ 2^{24} \\ 2^{28} \\ 2^{32} \\ 2^{36}
		\\ 2^{40} 
	\end{array}
	&
	\begin{array}{l}
	2.0 \rule{0pt}{14pt} \\
	2.638 \\
	2.712\,992 \\
	2.717\,950\,081 \\
	2.718\,261\,089\,905 \\
	2.718\,280\,532\,282 \\
	2.718\,281\,747\,448 \\
	2.718\,281\,823\,396 \\
	2.718\,281\,828\,142 \\
	2.718\,281\,828\,439 \\
	2.718\,281\,828\,458 \\
	\end{array}
\end{array}
	$$
	\addtolength{\arraycolsep}{-12pt}

	The {\em true\/} value of $e$\index{$e$} is 
	\mar{Expressing $e$ as a limit}
the limit of these approximations as we take $n$ arbitrarily large:
	$$
	e = \lim_{n \rightarrow \infty} 
		\left( 1 + 1/n \right)^n = 2.71828182845904\ldots
	$$

We can generalize the preceding to get an expression for $e^T$ for any
value of $T$.  In exactly the same way as we did above, divide the
interval from 0 to $T$ into $n$ pieces, each of width $\Delta t =
T/n$.  Starting from $t=0$ and $y(0)=1$ and applying Euler's method,
we find, using the same algebraic simplifications, that after $n$
steps the value for $t$ will be $T$ and $y$ will be $(1+T/n)^n$.
Since these approximations approach the true value of the function as
$n \rightarrow \infty$, we have that
\[
e^T = \lim_{n \rightarrow \infty} 
	\left( 1 + T/n \right)^n \mbox{     for any value of $T$.}
\]


\subsection{Differential Equations Define Functions}

There is an important point underlying the operations we just
performed having to do with the question of {\bf
  computability}\index{computability}.  While it may be appalling to
think about doing it by hand, there is nothing conceptually difficult
about evaluating an expression like
%
\mar{Euler's method uses only arithmetic}
%
$(1+1/1000)^{1000}$---all we need are ordinary addition, division, and
multiplication.  In fact, for any differential equation, Euler's
method generates a solution using only ordinary arithmetic.

By contrast, think for a moment about the earlier method for
evaluating $e$ by evaluating expressions like
$(2.718^{.0001}-2.718^{-.0001})/.0002$ and seeing whether we get a
value bigger than or less than 1.  While a calculator or a computer
will readily give us a value, how does it ``know'' what
$2.718^{.0001}$ is?  The fact is, it doesn't have a built-in
exponentiator which lets it know immediately what the value of this
expression is any more than we do.  A computer---like humans---can
essentially only add, subtract, and multiply.  Any other operation has
to be reduced to these operations somehow.  Thus when we use a
computer to evaluate something like $2.718^{.0001}$, we actual trigger
a fairly elaborate program having little directly to do with raising
numbers to powers which produces an approximation to the desired
number.  It turns out that if you use the \fbox{$x^y$} key on your
calculator to evaluate $2^5$ it doesn't come up with the answer by
multiplying 2 by itself 5 times, but uses this more complicated
program.

There is often a large gap between naming and defining a function, and
being able to compute
%
\mar{Defining a function is not the same as being able to evaluate it}
%
values for it to four or five decimals.  Think about the trigonometric
functions for a moment.  You have probably seen several definitions of
the cosine function by now---as the ratio of the adjacent side over
the hypotenuse of a right triangle, or as the $x$-coordinate of a
point moving around a circle of radius 1.  Yet neither of these
definitions would help you calculate $\cos(2)$ to five decimals.  It
turns out that most methods for evaluating functions are based on the
way the derivatives of the functions behave.  While we will have much
more to say about this in chapter 10, Euler's method is a good first
example of this.

Returning to exponents, think what would be involved in evaluating
$2^{\sqrt{3}}$ using the pre-calculus concept of exponents.  We might
first get a series of rational approximations to
$\sqrt{3}=1.73205081\ldots$: $17/10=1.7$, $173/100=1.73$,
$433/250=1.732$, and so on.  We would then calculate
\begin{align*}
2^{17/10}&=(\sqrt[10]{2})^{17} \\
2^{173/100}&=(\sqrt[100]{2})^{173} \\
2^{433/250}&=(\sqrt[250]{2})^{433} \\
           &\ \, \vdots
\end{align*}
Even evaluating the first of these approximations would involve
finding the 10th root of 2 and raising it to the 17th power, which
would not be easy.  We would continue with these approximations until
the desired number of digits remained fixed.

By contrast, evaluating $e^{\sqrt{3}}$ by Euler's method is very
straightforward.  As we saw above, it reduces to evaluating
$(1+1.73205\ldots/n)^n$ for increasing values of $n$ until the desired
number of digits remains fixed.  Moreover, this same process works
just as well for any kind of exponent---positive or negative, rational
or irrational.

In fact, {\em all} the properties of the exponential function follow
from the fact that it is the solution to its initial value problem, so
we could have made this the definition in the first place.  This would
have given us the benefit of coherence (not having to distinguish
among different kinds of exponents) and direct computability.  It
would also directly reflect the primary reason the exponential
function is important, namely that its rate of change is proportional
to its value.  Since the process of deducing the properties of a
function from its defining equation will be important later on, and
since it is a good exercise in some of the theoretical ideas we've
been developing, let's see how this works.

We will assume {\em nothing\/} about the function $y = e^t$.  
%
\mar{Defining $\exp(t)$}
%
Instead, we begin simply with the observation that each initial value
problem defines a function---its solution.  Therefore, the specific
problem
	$$
	y' = y \qquad y(0) = 1
	$$
defines a function;\index{$\exp$} 
we call it $y = \exp(t)$.  At the outset, all we know about the function $\exp(t)$ is that
	$$
	\exp'(t) = \exp(t) \qquad \exp(0) = 1.
	$$
As before, we can use Euler's method to evaluate $\exp(1)$, which we
will call $e$.  From these facts alone we want to {\em deduce\/} that
$\exp(t)=e^t$ for all values of~$t$.  We will actually show this only
for all {\em rational} values of $t$, since there is, as we've seen, a
bit of hand-waving about what it means to raise a number to an
irrational power. The following theorem is the key to establishing
this result.

\medskip
\noindent
{\bf Theorem 1}.  {\em For any real numbers\/ $r$ and\/ $s$},
	$$
	\exp(r+s) = \exp(r) \cdot \exp(s).
	$$

We will prove this result shortly, but let's see what we can deduce
from it.  First off, note that
$$
	\exp(2) = \exp(1+1) = \exp(1) \cdot \exp(1) = 
		\left(\exp(1)\right)^2=e^2.
	$$
Notice that we invoked Theorem 1 to equate $\exp(1+1)$ with $\exp(1) \cdot \exp(1)$.  In a similar way,
	$$
	\exp(3) = \exp(2+1) = \exp(2) \cdot \exp(1) = 
		\left(\exp(1)\right)^2 \cdot \exp(1) = 
		\left(\exp(1)\right)^3=e^3.
	$$
Repeating this argument for any positive integer $m$, we get

\noindent
{\bf Corollary 1}.  {\em For any positive integer $m$},	$$
	\exp(m) = \left(\exp(1)\right)^m = e^m.
	$$

We can also express $\exp(t)$ in terms of $e $ when $t$ is a {\em
  negative\/} integer.
%
\mar{Negative integers}
%
We begin with another consequence of Theorem~1:
	$$
	1 = \exp(0) = \exp(-1+1) = \exp(-1) \cdot \exp(1).
	$$
This says $e = \exp(1)$ is the reciprocal of $\exp(-1)$:
	$$
	\exp(-1) = \left(\exp(1)\right)^{-1} = e^{-1}.
	$$
Since $-2 = -1 - 1$, $-3 = -2 - 1$, and so forth, we can eventually show that

\pagebreak

\noindent
{\bf Corollary 2}.  {\em For any negative integer $-m$},	$$
	\exp(-m) = \exp(-1-1-\ldots -1)=\left(\exp(-1)\right)^{m}=\left(\exp(1)\right)^{-m} = e^{-m}.
	$$

	We can even do the same thing with fractions.  
	\mar{Rational numbers}
Here's how to deal with $\exp(1/3)$, for example:
	\begin{align*}
\exp(1) &= \exp\left(\tfrac{1}{3} 
		+ \tfrac{1}{3} + \tfrac{1}{3}\right)  \\
	&= \exp(1/3) \cdot \exp(1/3) \cdot \exp(1/3) \\
	&= \left( \exp(1/3) \right)^3,
	\end{align*}
so $\exp(1/3)$ is the cube root of $\exp(1)$:
	$$
	\exp(1/3) = \left( \exp(1) \right)^{1/3} = e^{1/3}.
	$$
A similar argument will show that 

\medskip

\noindent
{\bf Corollary 3}.  {\em For any positive integer $n$}, $\exp(1/n) =
e^{1/n}$.

\medskip

Finally, we can deal with any rational number $m/n$:
	$$
	\exp(m/n) = \left( \exp(1/n) \right)^m = 
		\left( e^{1/n} \right)^m = e^{m/n}.
	$$
This leads to
\medskip

\noindent
{\bf Theorem 2}.  {\em For any rational number\/ $r$}, $\exp(r) = \left( \exp(1) \right)^{\displaystyle r} = e^{\displaystyle r}$.
\medskip

	In other words, Theorem~1 implies that the function $\exp(t)$ is the same function as the exponential function $e^t$---at least when $t$ is a rational number $m/n$, as claimed.  We could now prove that $\exp(t)=e^t$ when $t$ is an irrational number, which would require being clearer about what it means to raise a number to an irrational power than most high school texts are.  

We adopt a more attractive option.  Since $\exp(t)$ and $e^t$ agree at rational values of $t$, and since $\exp(t)$ is well-defined for {\em all} values of $t$---including irrational numbers---we {\em define} $e^t$ for irrational values of $t$ by setting it equal to $\exp(t)$.

\smallskip
\noindent
{\bf Proof of Theorem 1}

\noindent
The proof uses the Existence and Uniqueness Principle for differential equations we articulated  	\mar{The proof of Theorem~1}
earlier:  if two functions satisfy the same differential equation and satisfy the same initial conditions then they have to be the same function.

	Theorem 1 involves two fixed real numbers, $r$ and $s$.  
We fix one of them, say $r$, to define two new functions of $t$:
	$$
	P(t) = \exp(r + t) \qquad Q(t) = \exp(r) \cdot \exp(t).
	$$
We shall show that {\em both\/} of these functions are solutions to the same initial value problem:
	$$
	\dfrac{dy}{dt} = y \qquad y(0) = \exp(r).
	$$
(Remember, $\exp(r)$ is a constant, because $r$ is fixed.)  

If we show this, it will then follow that $P(t)$ and $Q(t)$ must be
the same function. Since
\begin{align*}
P(0) &= \exp(r + 0) = \exp(r) \\
Q(0) &= \exp(r) \cdot \exp(0) = \exp(r) \cdot 1 = \exp(r),
\end{align*}
$P(t)$ and $Q(t)$ both satisfy the initial condition $y(0) = \exp(r)$.
Next we show that they both satisfy the differential equation $y' = y$:
	$$
	Q'(t) = \left(\exp(r) \cdot \exp(t) \right)' =
		\exp(r) \cdot \left(\exp(t) \right)'= 
		\exp(r) \cdot \exp(t) = Q(t),
	$$
so $Q(t)$ is a solution.  To differentiate $P(t)$ we construct a chain:
	$$
	P = \exp(u) \qquad \mbox{where} \qquad u = r+ t.
	$$
Then $dP/du = \exp(u)$ and $du/dt = 1$, so
	$$
	P'(t) = \dfrac{dP}{du} \cdot \dfrac{du}{dt} = 
		\exp(u) \cdot 1 = P(t) \cdot 1 = P(t),
	$$ 
so $P(t)$ is also a solution.  Therefore $P(t)$ and $Q(t)$ must be the
same function.  It follows then that $P(t)=Q(t)$ for all values of
$t$, in particular for $t=s$.  But this means that
	$$
	\exp(r + s) = P(s) = Q(s) = \exp(r) \cdot \exp(s),
	$$
which is exactly the statement of Theorem 1, and so completes the proof.

\medskip 

Now that we have established $\exp(x) = e^x$, we shall call $\exp(x)$
{\bf \emph{the} exponential function} and we shall use the forms $e^x$ and
$\exp(x)$ interchangeably.  The following theorem summarizes several
more properties of the exponential
function.\index{function!exponential}
\medskip

\noindent
{\bf Theorem 3}.  {\em For any real numbers\/ $r$ and\/ $s$}, 
	\begin{align*}
	\exp(s) &> 0 \\
	\exp(-s) &= \dfrac{1}{\exp(s)} \\[3pt]
	\exp(r - s) &= \dfrac{\exp(r)}{\exp(s)} \\[3pt]
	\exp(rs) &= \left( \exp(r) \right)^{\displaystyle s} 
		= \left( \exp(s) \right)^{\displaystyle r}.
	\end{align*}
To make the statements in this theorem seem more natural, you should
stop and translate them from $\exp(x)$ to $e^x$.  Proofs will be
covered in the exercises.


\subsection{Exponential Growth}
\index{growth!exponential}

The function $\exp(x) = e^x$, like polynomials and the sine and cosine
functions, is defined for all real numbers.  Nevertheless, it behaves
in a way that is quite different from any of those functions.

One difference occurs when $x$ is large, either positive or negative.
The sine function and the cosine function stay bounded between $+1$
and $-1$ over their entire domain.  By contrast, every polynomial
``blows up'' as $x \to \pm \infty$.  In this regard, the exponential
function is a hybrid.  As $x \to -\infty$, $\exp(x) \to 0$.  As $x \to
+ \infty$, however, $\exp(x) \to +\infty$.

Let's look more closely at what happens 
%
\mar{How fast do $x^n$ and $e^x$ become infinite?}
%
to power functions $x^n$ and the exponential function $e^x$ as $x \to
\infty$.  Both kinds of functions ``blow up'' but they do so at quite
different rates, as we shall see.  Before we compare power and
exponential functions directly, let's compare one power of $x$ with
another---say $x^2$ with $x^5$.  As $x \to \infty$, both $x^2$ and
$x^5$ get very large.  However, $x^2$ is only a small fraction of the
size of $x^5$, and that fraction gets smaller, the larger $x$ is.  The
following table demonstrates this.  Even though $x^2$ grows enormous,
we interpret the fact that $x^2/x^5 \to 0$ to mean that $x^2$ {\em
  grows more slowly than\/} $x^5$.  \addtolength{\arraycolsep}{14pt}
	$$
	\begin{array}{cccc}
	x & x^2 & x^5 & x^2/x^5 \\[2pt] \hline
	\rule{0pt}{14pt} \hphantom{00}10 \rule{0pt}{14pt} &
		10^2 & 10^{5\hphantom{0}} & 10^{-3} \\
	\hphantom{0}100 & 10^4 & 10^{10} & 10^{-6} \\
	1000 & 10^6 & 10^{15} & 10^{-9} \\
	\downarrow & \downarrow & \downarrow & \downarrow \\
	\infty & \infty & \infty & 0 
	\end{array}
	$$
	\addtolength{\arraycolsep}{-14pt}

It should be clear to you that we can compare 
%
\mar{$x^p$ grows more slowly than $x^q$ if $p < q$}
%
{\em any\/} two powers of $x$ this way.  We will find that $x^p$ grows
more slowly than $x^q$ if, and only if, $p < q$.  To prove this, we
must see what happens to the ratio $x^p/x^q$, as $x \to +\infty$.  We
can write $x^p/x^q = 1/x^{q - p}$, and the exponent $q - p$ that
appears here is positive, because $q > p$.  Consequently, as $x \to
\infty$, $x^{q-p} \to \infty$ as well, and therefore $1/x^{q-p} \to
0$.  This completes the proof.

\medskip

How does $e^x$ compare to $x^p$?  To make it tough on $e^x$, let's
compare it to $x^{50}$.  We know already that $x^{50}$ grows faster
than any lower power of $x$.  The table below compares $x^{50}$ to
$e^x$.  However, the numbers involved are so large that the table
shows only their {\em order of magnitude\/}---that is, the number of
digits they contain.  At the start, $x^{50}$ is {\em much\/} larger
than $e^x$.  However, by the time $x = 500$, the ratio $x^{50}/e^x$ is
so small its first 82 decimal places are zero!
\addtolength{\arraycolsep}{2pt}
	$$
\begin{array}{cccc}
	x & x^{50} & e^x & x^{50}/e^x \\[2pt] \hline
	\begin{array}{r}
	\rule{0pt}{16pt} 100 \\ 200 \\ 300 \\ 400 \\ 500
	\end{array}
	&
	\begin{array}{r}
	\rule{0pt}{16pt} \sim 10^{100} \\ 10^{115} \\ 10^{123}
		\\ 10^{130} \\ 10^{134} 
	\end{array}
	&
	\begin{array}{r}
		\rule{0pt}{16pt} \sim 10^{43\hphantom{0}} \\ 
		10^{86\hphantom{0}}  \\
		10^{130}  \\
		10^{173}  \\
		10^{217}  
	\end{array}
	&
	\begin{array}{r}
		\rule{0pt}{16pt} \sim 10^{56\hphantom{-}} \\ 
		10^{28\hphantom{-}}  \\
		10^{-7\hphantom{0}}  \\
		10^{-44}  \\
		10^{-83}  
	\end{array}
	\\
	\downarrow & \downarrow & \downarrow & \downarrow \\
	\infty & \infty & \infty & 0 
\end{array}
	$$
	\addtolength{\arraycolsep}{-2pt}
	

So $x^{50}$ grows more slowly than $e^x$, 
%
\mar{$e^x$ grows more rapidly than any power of $x$}
%
and so does any lower power of $x$.  Perhaps a higher power of $x$
would do better.  It does, but ultimately the ratio $x^p/e^x \to 0$,
no matter how large the power $p$ is.  We don't yet have all the tools
needed to prove this, but we will after we introduce the logarithm
function in the next section.

The speed of exponential growth has had an impact in computer science.
In many cases, the number of operations needed to calculate a
particular quantity is a power of the number of digits of precision
required in the answer.  Sometimes, though, the number of operations
is an {\em exponential\/} function of the number of digits.  When that
happens, the number of operations can quickly exceed the capacity of
the computer.  In this way, some problems that can be solved by an
algorithm that is straightforward in theoretical terms are intractable
in practical terms.


\subsection{Exercises}

\subsubsection{The exponential functions $b^t$}
\index{function!exponential}

\vspace{-10pt}
\num  Use a graphing utility or a calculator to approximate the slopes of the following functions at the origin and show:

\sub  If  $f(t) = (2.71)^t$, then $f'(0)<1$.

\sub  If  $g(t) = (2.72)^t$, then $g'(0)>1$.

\sub  Use parts (a) and (b) to explain why $2.71 < e < 2.72$.


\numa  In the same way find the value of the parameter $k_b$ for the bases $b =$ .5, .75, and .9 accurate to 3 decimal places.

\sub  What is the shape of the graph of $y=b^t$ when $0 < b < 1$?  What does that imply about the {\em sign\/} of $k_b$ for $0 < b < 1$?  Explain your reasoning.


\subsubsection{Differentiating exponential functions}
\index{differentiation}

\vspace{-10pt}
\num  Differentiate the following functions.

\sub  $7e^{\displaystyle 3x}$

\sub  $Ce^{\displaystyle kx}$, where $C$ and $k$ are constants.

\sub  $1.5e^{\displaystyle t}$

\sub $1.5e^{\displaystyle 2t}$

\sub  $2e^{\displaystyle 3x} - 3e^{\displaystyle 2x}$

\sub  $e^{\displaystyle \cos t}$


\num  Find partial derivatives of the following functions. 

\sub  $e^{\displaystyle xy}$

\sub  $3x^{\displaystyle 2}e^{\displaystyle 2y}$

\sub  $e^{\displaystyle u} \, \sin v$

\sub  $e^{\displaystyle u \sin(v)}$


\subsubsection{Powers of $e$}

\vspace{-10pt}
\num  Simplify the following and rewrite as powers of $e$.  For each, explain your work, citing any theorems you use.

\sub  $\exp(2x+3)$

\sub  $(\exp(x))^2$

\sub  $\exp(17x)/\exp(5x)$ 


\num  Use the second property in Theorem 3  to explain why
	$$
	\lim_{t \rightarrow - \infty} \exp(t) = 0.
	$$

\num  This purpose of this exercise is to prove the fourth property listed in Theorem~3: $\exp(rs) = (\exp(s))^r$, for all real numbers $r$ and $s$.  The idea of the proof is the same as for Theorem~1: show that two different-looking functions solve the same initial value problem, thus demonstrating that the functions must be the same.  The initial value problem is
	$$
	y' = ry \qquad y(0) = 1.
	$$

\sub  Show that $P(t) = \exp(rt)$ solves the initial value problem.  (You need to use the chain rule.)

\sub  Show that $Q(t) = (\exp(t))^r$ solves the initial value problem.  (Here use the chain $Q = u^r$, where $u = \exp(t)$.  There is a bit of algebra involved.)

\sub  From parts (a) and (b), and the fact that an initial value problem has a unique solution, it follows that $P(t) = Q(t)$, for every $t$.  Explain how this establishes the result.


\subsubsection{Solving $y'= ky$ using $e^t$}

\vspace{-10pt}
\num  {\bf Poland and Afghanistan}.\index{Poland}\index{Afghanistan}  Refer to problem 25 in chapter~1.2. 

\sub  Write out the initial value problems that summarize the information about the populations $P$ and $A$ given in parts (a) and (b) of problem 21.

\sub  Write formulas for the solutions $P$ and $A$ of these initial value problems.

\sub  Use your formulas in part (b) (and a calculator) to find the population of each country in the year 2005.  What were the populations in 1965?


\num {\bf Bacterial growth}.\index{growth!bacterial} Refer to problem
26 in chapter~1.2.

\sub  Assuming that we begin with the colony of bacteria weighing 32 grams, write out the initial value problem that summarizes the information about the weight $P$ of the colony.

\sub  Write a formula for the solution $P$ of this initial value problem.

\sub  How much does the colony weigh after 30 minutes?  after 2 hours?


\num  {\bf Radioactivity}.\index{radioactivity}  Refer to problem~27 in chapter~1.2.

\sub  Assuming that when we begin the sample of radium weighs 1~gram,  write out the initial value problem that summarizes the information about the weight $R$ of the sample.

\sub  How much did the sample weigh 20 years ago?  How much will it weigh 200 years hence?


\num {\bf Intensity of radiation}.\index{radiation} As gamma rays
travel through an object, their intensity $I$ decreases with the
distance $x$ that they have travelled.  This is called {\bf
  absorption}.  The absorption rate $dI/dx$ is proportional to the
intensity.  For some materials the multiplier in this proportion is
large; they are used as radiation shields.

\sub Write down a differential equation which models the intensity of
gamma rays $I(x)$ as a function of distance $x$.

\sub Some materials, such as lead, are better shields than others,
such as air.  How would this difference be expressed in your
differential equation?

\sub Assume the unshielded intensity of the gamma rays is $I_0$.
Write a formula for the intensity $I$ in terms of the distance $x$ and
verify that it gives a solution of the initial value problem.


\num  In this problem you will find a solution for the initial value problem $y'=ky$ and $y(t_0)=C$.  (Notice that this isn't the original initial value problem, because $t_0$ was 0 originally.)

\sub Explain why you may assume $y=Ae^{kt}$ for some constant $A$.

\sub  Find $A$ in terms of $k, C$ and $t_0$.


\subsubsection{Solving other differential equations}

\vspace{-10pt}
\numa  {\bf Newton's law of cooling}.\index{Newton's law of cooling}  Verify that
	$$
	Q(t)=70e^{-.1\,t} + 20
	$$
is a solution to the initial value problem $Q'(t) = -.1(Q-20)$, $Q(0) = 90$.  What is the relationship between this formula and the one found in problem 11 in section~2?

\sub Verify that
$$
Q(t)=(Q_0-A)e^{-kt} + A
$$
is a solution to the the initial value problem $Q'(t)=-k(Q-A)$,
\mbox{$Q(0)=B$}.  What is the relationship between this formula and the one
found in problem~13 in section~2?


\num In {\em An Essay on the Principle of Population}, written in
1798, the British economist Thomas Robert Malthus
(1766--1834)\index{Malthus, T.} argued that food supplies grow at a
constant rate, while human populations naturally grow at a constant
{\em per capita} rate.  He therefore predicted that human populations
would inevitably run out of food (unless population growth was
suppressed by unnatural means).


\sub Write differential equations for the size $P$ of a human
population and the size $F$ of the food supply that reflect Malthus'
assumptions about growth rates.

\sub Keep track of the population in millions, and measure the food
supply in millions of units, where one unit of food feeds one person
for one year.  Malthus' data suggested to him that the food supply in
Great Britain was growing at about .28 million units per year and the
per capita growth rate of the population was 2.8\% per year.  Let
$t=0$ be the year 1798, when Malthus estimated the population of the
British Isles was $P=$ 7 million people.  He assumed his countrymen
were on average adequately nourished, so he estimated that the food
supply was $F=$ 7 million units of food.  Using these values, write
formulas for the solutions $P=P(t)$ and $F=F(t)$ of the differential
equations in (a).
 
\sub Use the formulas in (b) to calculate the amount of food and the
population at 25 year intervals for 100 years.  Use these values to
help you sketch graphs of $P=P(t)$ and $F=F(t)$ on the same axes.

\sub The per capita food supply in any year equals the ratio
$F(t)/P(t)$.  What happens to this ratio as $t$ grows larger and
larger?  (Use your graphs in (c) to assist your explanation.)  Do your
results support Malthus's prediction?  Explain.


\numa {\bf Falling bodies}.\index{falling object} Using the base $e$
instead of the base 2, modify the solution $v(t)$ to the initial value
problem
\[
\dfrac{dv}{dt} = -g - bv \qquad v(0) = 0
\]
that appears in exercise~21 on page~\pageref{ex:falling body1}.  Show
that the modified expression is still a solution.

\sub If an object that falls against air resistance is $x(t)$ meters
above the ground after $t$ seconds, and it started $x_0$ meters above
the ground, then it is the solution of the initial value problem
\[
\dfrac{dx}{dt} = v(t) \qquad x(0) = x_0,
\]
where $v(t)$ is the velocity function from the previous exercise.
Find a formula for $x(t)$ using the exponential function with base
$e$.  (Compare this formula with the one in exercise~22~(a),
page~\pageref{ex:falling body2}.)

\sub Suppose the coefficient of air resistance is .2 per second.  If a
body is held motionless 200 meters above the ground, and then
released, how far will it fall in 1 second?  In 2 seconds?  Use your
formula from part (b).  Compare these values with those you obtained
in exercise~22~(b), page~\pageref{ex:falling body3}.

\subsubsection{Interest rates} 
\index{interest rates}

Bank advertisements sometimes look like this:
\vskip .125in
\centerline{\bf Civic Bank and Trust}
\begin{itemize}
\item Annual rate of interest 6\%.
\item Compounded monthly.
\item Effective rate of interest 6.17\%.
\end{itemize}
The first item seems very straightforward.  The bank pays 6\% interest
per year.  Thus if you deposit \$100.00 for one year then at the end
of the year you would expect to have \$106.00. Mathematically this is
the simplest way to compute interest; each year add 6\% to the
account.  The biggest problem with this is that people often make
deposits for odd fractions of a year, so if interest were paid only
once each year then a depositor who withdrew her money after 11 months
would receive no interest.  To avoid this problem banks usually
compute and pay interest more frequently.  The Civic Bank and Trust
advertises interest {\bf compounded} monthly.  This means that the
bank computes interest each month and credits it (that is, adds it) to
the account.
\begin{center}
{\footnotesize
%\vskip .25in
\begin{tabular}{||r|r|r|r||} \hline
Month & Start & Interest & End \\ \hline
 1 & \$100.0000 &    .5000 & \$100.5000 \\
 2 & \$100.5000 &    .5025 & \$101.0025 \\
 3 & \$101.0025 &    .5050 & \$101.5075 \\
 4 & \$101.5075 &    .5075 & \$102.0151 \\
 5 & \$102.0151 &    .5101 & \$102.5251 \\
 6 & \$102.5251 &    .5126 & \$103.0378 \\
 7 & \$103.0378 &    .5152 & \$103.5529 \\
 8 & \$103.5529 &    .5178 & \$104.0707 \\
 9 & \$104.0707 &    .5204 & \$104.5911 \\
10 & \$104.5911 &    .5230 & \$105.1140 \\
11 & \$105.1140 &    .5256 & \$105.6396 \\
12 & \$105.6396 &    .5282 & \$106.1678 \\ \hline
\end{tabular}
%\vskip .25in
}
\end{center}

Since this particular account pays interest at the rate of 6\% per
year and there are 12 months in a year the interest rate is 6\%/12 =
0.5\% per month.  The following table shows the interest computations
for one year for a bank account earning 6\% annual interest compounded
monthly.

Notice that at the end of the year the account contains \$106.17.  It
has effectively earned 6.17\% interest.  This is the meaning of the
advertised {\em effective\/} rate of interest.  The reason that the
effective rate of interest is higher than the original rate of
interest is that the interest earned each month itself earns interest
in each succeeding month.  (We first encountered this phenomenon when
we were trying to follow the values of $S$, $I$, and $R$ into the
future.)  The difference between the original rate of interest and the
effective rate can be very significant.  Banks routinely advertise the
effective rate to attract depositors.  Of course, banks do the same
computations for loans.  They rarely advertise the effective rate of
interest for loans because customers might be repelled by the true
cost of borrowing.

The effective rate of interest can be computed much more quickly than
we did in the previous table. Let $R$ denote the annual interest rate
as a decimal.  For example, if the interest rate is 6\% then $R =
0.06$. If interest is compounded $n$ times per year then each time it
is compounded the interest rate is $R/n$.  Thus each time you compound
the interest you compute
$$
V + \left(\dfrac{R}{n}\right)V = \left({1 + \dfrac{R}{n}}\right)V
$$
where $V$ is the value of the current deposit.  This computation is
done $n$ times during the course of a year.  So, if the original
deposit has value $V$, after one year it will be worth
$$
\left({1 + \dfrac{R}{n}}\right)^nV.
$$
For our example above this works out to
$$
\left(1 + \dfrac{0.06}{12}\right)^{12}V = 1.061678 \, V 
$$
and the effective interest rate is 6.1678\%.

Many banks now compound interest daily.  Some even compound interest
{\em continuously}.  The value of a deposit in an account with
interest compounded continuously at the rate of 6\% per year, for
example, grows according to the differential equation
$$
V' = 0.06V.
$$


\num Many credit cards charge interest at an annual rate of 18\%.  If
this rate were compounded monthly what would the effective annual rate
be?

\num In fact many credit cards compound interest daily.  What is the
effective rate of interest for 18\% interest compounded daily?  Assume
that there are 365 days in a year.

\num The assumption that a year has 365 days is, in fact, {\em not\/}
made by banks.  They figure every one of the 12 months has 30 days, so
their year is 360 days long.  This practice stem from the time when
interest computations were done by hand or by tables, so simplicity
won out over precision.  Therefore when banks compute interest they
find the daily rate of interest by dividing the annual rate of
interest by 360.  For example, if the annual rate of interest is 18\%
then the daily rate of interest is 0.05\%.  Find the effective rate of
interest for 18\% compounded 360 times per year.

\num In fact, once they've obtained the daily rate as 1/360-th of the
annual rate, banks then compute the interest {\em every\/} day of the
year.  They compound the interest 365 times.  Find the effective rate
of interest if the annual rate of interest is 18\% and the
computations are done by banks.  First, compute the daily rate by
dividing the annual rate by 360 and then compute interest using this
daily rate 365 times.

\num Consider the following advertisement.
\vskip .125in
\centerline{\bf Civic Bank and Trust}
\begin{itemize}
\item Annual rate of interest 6\%.
\item Compounded daily.
\item Effective rate of interest 6.2716\%.
\end{itemize}
Find the effective rate of interest for an annual rate of 6\%
compounded daily in the straightforward way---using 1/365-th of the
annual rate 365 times.  Then do the computations the way they are done
in a bank.  Compare your two answers.

\num There are two advertisements in the newspaper for savings
accounts in two different banks.  The first offers 6\% interest
compounded quarterly (that is, four times per year).  The second
offers 5.5\% interest compounded continuously.  Which account is
better? Explain.


%%%%%%%%%%%%%%%%%%%%%%%%%%%%%%%%%%%%%%%%%%%%%%%%%%%%%%%%%%%%%%%%%%%%%%%%%%%
%                                                                         %
%                                Section 4                                %
%                                                                         %
%%%%%%%%%%%%%%%%%%%%%%%%%%%%%%%%%%%%%%%%%%%%%%%%%%%%%%%%%%%%%%%%%%%%%%%%%%%

\section{The Logarithm Function}

Suppose a population is growing at the net rate of 3 births per
thousand persons per year.  If there are 100,000 persons now, how many
will there be 37 years from now?  How long will it take the population
to double?

Translating into mathematics, we want to find the function $P(t)$ that
solves the initial value problem
	$$
	P'(t) = .003P(t) \qquad \hbox{and} \qquad 
		P(0) = 100000.
     $$ 
Using the results of section~3 we know that the solution is the exponential
function
	$$
	P(t) = 100000\,e^{.003\,t}.
	$$ 
The size of the population 37 years from now will therefore be 
\begin{align*}
	P(37) & =  100000\,e^{.111} \\
	      & =  100000 \times  1.117395 \\
	      & \approx  111740 \hbox{ people}
\end{align*}
To find out how long it will take 
%
\mar{The doubling time\\ of a population}
%
the population to double,\index{doubling time} we want to find a value
for $t$ so that $P(t) = 200000$.  In other words, we need to solve for
$t$ in the equation
	$$
	100000\,e^{.003\,t} = 200000.
	$$
Dividing both sides by 100,000, we have
	$$
	e^{.003\,t} = 2.
	$$

We can't proceed because one side is expressed in exponential form
while the other isn't.  One remedy is to express 2 in exponential
form.  In fact, $2 = e^{.693147}$, as you should verify with a
calculator.  Then
	$$
e^{.003\,t} = 2 = e^{.693147}
	 \qquad \mbox{implies} \qquad .003\,t = .693147,
	$$
so $t = .693147/.003 = 231.049$.  Thus it will take about 231 years
for the population to double.

\medskip

To determine the doubling time of the population we had to know the
number $b$ for which
	$$
	\exp(b) = e^{\displaystyle b} = 2.
	$$
This is an aspect 
%
\mar{Solving an\\ exponential equation}
%
of a very general question: given a positive number $a$, find a number
$b$ for which
	$$
	e^{\displaystyle b} = a.
	$$
A glance at the graph of the exponential function below shows that, by
working backwards from any point $a>0$ on the vertical axis, we can
indeed find a unique point $b$ on the horizontal axis which gives us
$\exp(b) = a$.

\pagebreak

\mbox{}

\vspace*{2.2in}
\special{psfile=../ps/calc1/log1.ps}

This process of obtaining the number $b$ that satisfies $\exp(b) = a$
for any given positive number $a$ is a clear and unambiguous rule.
%
\mar{The natural logarithm function}
%
Thus, it defines a function.\index{function!logarithm}\index{logarithm
  function}\index{logarithm}\index{natural logarithm} This function is
called {\bf the natural logarithm},\index{natural logarithm} and it is
denoted $\ln(a)$, or sometimes $\log(a)$.  That is,
	$$
  \ln(a) = \log(a) = 
	\left\{ \mbox{the number $b$ for which $\exp(b) = a$}
	\right\}.
	$$
In other words, the two statements
	$$
	\ln(a) = b \quad \mbox{and} \quad \exp(b) = a
	$$
express exactly the same relation between the quantities $a$ and $b$.  

The question that led to the introduction of the logarithm function
was: what number gives the exponent to which $e$ must be raised in
order to produce the value 2?  This number is $\ln(2)$, and we
verified that $\ln(2) = .693147$.  Quite generally we can say that the
number $\ln(x)$ gives the exponent to which $e$ must be raised in
order to produce the value $x$:
	$$
	e^{\ln(x)} = x.
	$$

\noindent
If we set $y = \ln(x)$, then $x = e^{\displaystyle y}$ and we can
restate the last equation as a pair of companion equations:
	$$
	e^{\ln(x)} = x \quad \mbox{and} \quad
		\ln(e^{y}) = y.
	$$
The first equations says 
%
\mar{The logarithm and exponential functions are inverses}
%
{\em the exponential function ``undoes'' the effect of the logarithm
  function\/} and the second one says {\em the logarithm function
  ``undoes'' the effect of the exponential function\/}.  For this
reason the exponential and logarithm functions are said to be {\bf
  inverses}\index{inverse} of each other.

\pagebreak

Many of the other pairs of functions---sine and arscsine, squareroot
and squaring---that share a key on a calculator have this property.
There are even functions (at least one can be found on any calculator)
that are their own inverses---apply such a function to any number,
then apply this same function to the result, and you're back at the
original number.  What functions do this?  We will say more about
inverse functions later in this section.


\subsection{Properties of the Logarithm Function}

The inverse relationship allows us to translate each of the properties
of the exponential function into a corresponding statement about the
logarithm function. We list the major pairs of properties below.
\begin{center}
\begin{tabular}{ll}
{\bf exponential version} & {\bf logarithmic version}\\[5pt]
  $e^{0}=1$          & $\ln(1)=0$\\
  $e^{\displaystyle a+b}=e^{\displaystyle a}\cdot e^{\displaystyle b}$            & $\ln(m\cdot n)=\ln(m) + \ln(n)$\\
  $e^{\displaystyle a-b} = e^{\displaystyle a}/e^{\displaystyle b}$               & $\ln(m/n) = \ln(m) - \ln(n)$ \\
  $(e^{\displaystyle a})^{\displaystyle s} = e^{\displaystyle as}$                & $\ln(m^s) = s \cdot \ln(m)$ \\
 range of $e^{\displaystyle x}$ is all positive reals & domain of $\ln(x)$ is all positive reals\\
domain of $e^{\displaystyle x}$ is all real numbers & range of $\ln(x)$ is all real numbers \\
$e^{\displaystyle x} \to 0$ as $x \to -\infty$ &
	$\ln(x) \to -\infty$ as $x \to 0$ \\
$e^{\displaystyle x}$ grows faster & $\ln(x)$ goes to infinity slower\\
\hspace{.25in} than $x^n$, any $n> 0$ & \hspace{.25in} than $x^{1/n}$, any  $n > 0$
\end{tabular}
\end{center}
For each pair, we can use the exponential property and the inverse
relationship between exp and ln to establish the logarithmic property.
As an example, we will establish the second property.  You should be
able to demonstrate the others.

\smallskip

\noindent {\em Proof of the second property}. Remember that to show
$\ln(a)=b$, we need to show $e^b=a$.  In our case $a$ and $b$ are more
complicated.  We have
\[
a =  m\cdot n,  \qquad\quad
b =  \ln(m)+\ln(n);
\]
thus we need to show
$$
e^{\displaystyle \ln(m)+\ln(n)} = m \cdot n \, .
$$
But, by the exponential version of property 2, 
	$$
	e^{\displaystyle \ln(m)+\ln(n)} = 
		e^{\displaystyle \ln(m)} \cdot
		e^{\displaystyle \ln(n)} = m \cdot n,
$$
and our proof is complete.

\pagebreak

\subsection{The Derivative of the Logarithm Function}
\index{derivative}

Since the natural logarithm is a function in its own right, it is
reasonable to ask: what is the derivative of this function?
%
\mar{The graph of $\ln$}
%
Since the derivative describes the slope of the graph, let us begin by
examining the graph of ln.  Can we take advantage of the relationship
between ln and exp---a function whose graph we know well---as we do
this?  Indeed we can, by making the following observations.

\begin{itemize}
\item We know the point $(a, b)$ is on the graph of $y=\ln(x)$ if and
  only if $b = \ln(a)$.

\item  We know $b =  \ln(a)$ says the same thing as $a=e^{\displaystyle b}$.

\item Finally, we know $a=e^{\displaystyle b}$ is true if and only if
  the point $(b, a)$ is on the graph of $y=e^{\displaystyle x}$.

\end{itemize}

\noindent
Putting our observations together, we have
\begin{center}
\begin{tabular}{c}
$(a, b)$ is on the graph of $y=\ln(x)$  \\[3pt]
if and only if  \\[3pt]
$(b, a)$ is on the graph of $y=e^{\displaystyle x}$.
\end{tabular}
\end{center}
The picture below 
%
\mar{Reflection across\\ the $45^{\circ}$~line}
%%
demonstrates that the point $(a, b)$ and the point $(b, a)$ are
reflections\index{reflection} of each other about the $45^{\circ}$
line.  (Remember that points on the $45^{\circ}$ line have the same
$x$ and $y$ coordinates.)  This is because these two points are the
endpoints of the diagonal of a square whose other diagonal is the line
$y=x$.

\vspace*{120pt} 

\special{psfile=../ps/calc1/log2.ps}

Since we have just seen that every point $(a, b)$ on the graph of
$y=\ln(x)$ corresponds to a point $(b, a)$ on the graph of $y=e^x$, we
see that {\em the graphs of $y=\ln(x)$ and $y=e^x$ are the reflections
  of each other about the line $y=x$.}

\vspace*{190pt}

\special{psfile=../ps/calc1/log3.ps}

Finally, since the two graphs are reflections of one another,
\index{microscope}
%
\mar{Microscopic views at mirror image points}
%
a microscopic view of $\ln(x)$ at any point $(b, a)$ will be the
mirror image of the microscopic view of of $e^x$ at the point $(a,
b)$.  Any change in the $y$-value on one of these lines will
correspond to an equal change in the $x$-value in its mirror image,
and vice versa.  The figure below shows what microscopic views of a
pair of corresponding points look like, showing how a vertical change
in one line equals the horizontal change in the other, and conversely.

\vspace{210pt}

\special{psfile=../ps/calc1/log4.ps}

\newpage

It follows that the slopes of the two lines 
%
\mar{The slopes\\ are reciprocals}
%
must be reciprocals of each other.  This says that the rate of change
of $\ln(x)$ at $x = b$ is just the reciprocal of the rate of change of
$e^x$ at $x = a$, where $a = \ln(b)$.  But the rate of change of $e^x$
at $x = \ln(b)$ is just $e^{\ln(b)} = b$.  Therefore the rate of
change of $\ln(x)$ at $x = b$ is the reciprocal of this value, namely
$1/b$.  We have thus proved the following result:

\bigskip 
\noindent
{\bf Theorem 1}.   $(\ln(x))' = 1/x$.
\bigskip

Note that one interpretation of this theorem is that the function
$\ln(x)$ is the solution to a certain initial value problem, namely
	$$
	\dfrac{dy}{dx} = \dfrac{1}{x} \qquad y(1) = 0.
	$$
As was the case with the exponential function, we can now apply
Euler's method to this differential equation as an effective way to
compute values of $\ln(x)$.  Applications of this idea can be found in
the exercises.


\subsection{Exponential Growth}
\index{growth!exponential}

The logarithm gives us a useful tool for comparing 
%
\mar{Comparing rates\\ of growth}
%
the growth rates of exponential and power functions.  In the last section we claimed that $e^x$ grows faster than any power $x^p$ of $x$, as $x \to +\infty$.  We interpreted that to mean
	$$
	\lim_{x \to +\infty} \dfrac{x^p}{e^x} = 0,
	$$
for any number $p$.  Using the natural logarithm, we can now show why it is true.

	To analyze the quotient $Q = x^p/e^x$, we first replace it by its logarithm
	$$
	\ln{Q} = \ln\left(x^p/e^x\right) 
		= \ln\left(x^p\right) - \ln\left(e^x\right) 
		= p \ln{x} - x.
	$$
Several properties of the logarithm function were invoked here to reduce $\ln{Q}$ to $p \ln{x} - x$.  By another property of the logarithm function, if we can show $\ln{Q} \to -\infty$ we will have established our original claim that $Q \to 0$.  

	Let $y = \ln{Q} = p \ln{x} - x$.  We know $y$ is increasing when $dy/dx > 0$ and decreasing where $dy/dx < 0$.  Using the rules of differentiation, we find
	$$
	\dfrac{dy}{dx} = \dfrac{p}{x} - 1.
	$$

\pagebreak

\noindent
The expression $p/x - 1$ is positive when $x$ is 
%
\mar{$y = \ln{Q}$ decreases when $x > p$ \ldots}
%
less than $p$ and negative when $x$ is greater than $p$.  For $x$ near
$p$, the graph of $y$ must therefore look like this:
\begin{center}
\begin{picture}(190,90)(-30,-40)
	\put(-30,0){\vector(1,0){190}}
	\put(0,-40){\vector(0,1){90}}
	\put(160,-3){\makebox(0,0)[t]{\footnotesize $x$}}
	\put(-3,50){\makebox(0,0)[r]{\footnotesize $y$}}

	\begin{thicklines}
	\bezier{75}(15,30)(41.7,41.7)(75,25)
	\end{thicklines}

	\put(37.5,0){\line(0,1){1.5}}
	\put(37.5,-3){\makebox(0,0)[t]{\footnotesize $p$}}

\end{picture}
\end{center}

	Since $dy/dx$ remains negative as $x$ gets large, $y$ will continue to decrease.  
	\mar{\ldots but $y$ may still\\ ``level off''}
This does not, in itself, imply that $y \to -\infty$, however.  It's conceivable that $y$ might ``level off'' even as it continues to decrease---as it does in the next graph.
\begin{center}
\begin{picture}(190,90)(-30,-40)
	\put(-30,0){\vector(1,0){190}}
	\put(0,-40){\vector(0,1){90}}
	\put(160,-3){\makebox(0,0)[t]{\footnotesize $x$}}
	\put(-3,50){\makebox(0,0)[r]{\footnotesize $y$}}

	\begin{thicklines}
	\bezier{75}(15,30)(41.7,41.7)(75,25)
	\bezier{100}(75,25)(115,5)(155,5)
	\end{thicklines}

	\put(37.5,0){\line(0,1){1.5}}
	\put(37.5,-3){\makebox(0,0)[t]{\footnotesize $p$}}

\end{picture}
\end{center}

\enlargethispage*{5pt}

	However, we can show that $y$ does {\em not\/} ``level off'' in this way; it continues to plunge down to $-\infty$.  We start by assuming that $x$ has already become larger than $2p$: $x > 2p$.  Then $1/x < 1/2p$ (the bigger number has the smaller reciprocal), and thus $p/x < p/2p = 1/2$.  Thus, when $x > 2p$, 
	$$
	\dfrac{dy}{dx} = \dfrac{p}{x} - 1 < \dfrac{1}{2} - 1 = 
		- \, \dfrac{1}{2}.
	$$
In other words, 
	\mar{$y$ lies below a line that slopes down 
to~$-\infty$}
the slope of the graph of $y$ is more negative than $-1/2$.  The graph of $y$  must therefore lie {\em below\/} the straight line with slope $-1/2$ that we see below:
\begin{center}
\begin{picture}(190,90)(-30,-40)
	\put(-30,0){\vector(1,0){190}}
	\put(0,-40){\vector(0,1){90}}
	\put(160,-3){\makebox(0,0)[t]{\footnotesize $x$}}
	\put(-3,50){\makebox(0,0)[r]{\footnotesize $y$}}

	\begin{thicklines}
	\bezier{75}(15,30)(41.7,41.7)(75,25)
	\bezier{20}(75,25)(115,5)(155,-15)
	\end{thicklines}

	\put(37.5,0){\line(0,1){1.5}}
	\put(75,0){\line(0,1){1.5}}
	\put(37.5,-3){\makebox(0,0)[t]{\footnotesize $p$}}
	\put(75,-3){\makebox(0,0)[t]{\footnotesize $2p$}}

\end{picture}
\end{center}
This guarantees that $y = \ln{Q} \to -\infty$ as $x \to \infty$.  Hence $Q \to 0$, and since $Q = x^p/e^x$, we have shown that $e^x$ grows faster than any power of $x$.

\pagebreak

\subsection{The Exponential Functions $b^{\displaystyle x}$}
\index{function!exponential}

	We have come to adopt the exponential function $\exp(x) = e^x$ as the natural one for calculus, 
	\mar{One exponential function is enough}
and especially for dealing with differential equations of the form $dy/dx = ky$.  Initially, though, all exponential functions $b^x$ were on an equal footing.  With the natural logarithm function, however, a {\em single\/} exponential function will meet our needs.  Let's see why.

If $b$ is any positive real number, then $b = e^{\ln{b}}$.  Consequently,
	$$
	b^{x} = 
		\left(e^{\ln{b}}\right)^{x} =
		e^{\ln{b} \, \cdot \, x}.
	$$
In other words, $b^{\displaystyle x} = e^{\displaystyle cx}$, where $c
= \ln{b}$.
%
\mar{$b^{x} = e^{\ln{b} \cdot x}$}
%
Thus, every exponential function can be expressed in terms of $\exp$
in a simple way.  This is, in fact, the way computers evaluate
exponents, since the computer can raise any number to any power so
long as it has a way to evaluate the functions ln and exp.  For
instance, when you ask a computer or calculator to evaluate
%
\mar{How to calculate $2^5$} 
%
\verb+2^5+ (2 to the 5th power in most computer languages), it will
first calculate $\ln{2}$, then multiply this number by 5, then apply
exp to the result.  That is, it evaluates $2^5$ by thinking
$e^{5\ln{2}}$!  While this may seem a roundabout way to come up with
32, its virtue is that the computer needs only one algorithm to
calculate any base to any power, without having to consider different
cases.

	This expression gives us a new way to find the derivative of $b^x$.  We already know that
	$$
	( e^{\displaystyle cx})' = 
		c \cdot e^{\displaystyle cx},
	$$
for any constant $c$.  This follows from the chain rule.  When $c = \ln{b}$, we get
	$$
	(b^{\displaystyle x})' = 
		(e^{\displaystyle \ln(b) \cdot x})' = 
		\ln(b) \cdot e^{\displaystyle \ln(b) \cdot x} = 
		\ln(b) \cdot b^{\displaystyle x}.
	$$
Thus, $y = b^{\displaystyle x}$ is a solution to the differential equation
	$$
	\dfrac{dy}{dx} = \ln(b) \cdot y.
	$$
In chapter 3, we wrote this differential equation as
	$$
	\dfrac{dy}{dx} = k_b \cdot y.
	$$
We see now that $k_b = \ln(b)$.  
	\mar{$k_b = \ln(b)$}


We can use the connection between $k_b$ and the natural logarithm, and
between the natural logarithm and the exponential function, to gain
new insights.  For example, on page~\pageref{k_b} we argued that there
must be a value of $b$ for which $k_b = .02$.  This simply means
\[
\ln{b} = .02 \qquad \mbox{or} \qquad b = e^{.02}.
\]
In other words, we now have an explicit formula that tells us the
value of $b$ for which $k_b = .02$:
\[
b = e^{.02} = 1.02020134\ldots .
\]


\subsection{Inverse Functions}

	Most of what we have said about the exponential and logarithm functions carries over directly to {\em any\/} pair of inverse functions.  We begin by saying precisely what it means for two functions $f$ and $g$ to be inverses of each other.\index{inverse function}\index{function!inverse}

\bigskip
\noindent
{\bf Definition}.  Two functions $f$ and $g$ are {\bf inverses} if
\begin{align*}
	f(g(a)) &=  a  \\
	\mbox{and} \ g(f(b)) &=  b
\end{align*}
for every $a$ in the domain\index{domain} of $g$ and every $b$ in the domain of $f$.
\bigskip

Observe that if $f$ and $g$ are inverses of each other, then each one ``undoes'' the effect of the other by sending any value back to the number it came from via the other function.  One implication 
of this is that neither function can have two different input values going to the same output value.  For instance, suppose $b_1$ and $b_2$ get sent to the same value by $f$: $f(b_1)=f(b_2)$.  Applying $g$ to both sides of this equation we would get $b_1=g(f(b_1))=g(f(b_2))=b_2$, so $b_1$ and $b_2$ were actually the same number.  This is an important enough property that there is a name for it:

\medskip
\noindent
{\bf Definition.}  We say that a function $f$ is {\bf one-to-one}\index{one-to-one}, usually written as {\bf 1--1}, if it is true that whenever $x_1\neq x_2$ then it is also the case that  $f(x_1)\neq f(x_2)$.  Equivalently, $f$ is 1--1 if whenever $f(x_1)=f(x_2)$, then it must be true that $x_1=x_2$.

\medskip
\noindent
We have thus seen that only functions which are one-to-one can have
inverses.  This means that to establish inverses for some functions,
we will need to restrict their domains to regions where they are
one-to-one.  Let's re-examine the examples we mentioned earlier to see
how they fit this definition.

\medskip
\noindent
{\bf Example 1}.  Suppose $f(x)=\exp(x)$ and $g(x)=\ln(x)$.  Then the
equations
\begin{align*}
f(g(a)) & =  \exp(\ln(a)) = e^{\displaystyle \ln(a)} 
        = a \;\; \mbox{for}\ a>0 \\
	\mbox{and} \ g(f(b)) & =  \ln(\exp(b)) = 
	\ln(e^{\displaystyle b}) = b 
\end{align*}
hold for all real numbers $b$ and for all {\em positive\/} real
numbers $a$.  The domain of the exponential function is all real
numbers and the domain of the natural logarithm function is all
positive real numbers.

\medskip

\noindent
{\bf Example 2}.  Suppose $f(x)= x^2$ and $g(x) = \sqrt{x}$.
	\mar{$x^2$ is invertible\\ on $x \geq 0$}
The squaring function is not invertible on its natural domain because it is not one-to-one.  Since a number and its negative have the same square, we wouldn't know which one to send the square back to when we took the square root.  We can't avoid the problem by saying that $g(4)=\pm 2$, since a function has to have only one output for each input. The squaring function is invertible, though, if we restrict it to non-negative real numbers.  Then
\begin{align*}
f(g(a)) & =  (\sqrt{a})^2 = a \;\; (\mbox{for} \;\; a \geq 0) \\
\mbox{and} \ g(f(b)) & =  \sqrt{b^2} = b \;\; (\mbox{for} \;\; b \geq 0).
\end{align*}
The domain of the square root function is all $b \geq 0$.

Note that we could have restricted the domain of $f$ in another way to
make it one-to-one by considering only non-positive real numbers.  Now
$g$ is no longer the inverse of this restricted $f$.  For instance,
$g(f(-3))=g(9) = 3\neq -3$.  What would the inverse of $f$ be in this
case?

\medskip

\noindent
{\bf Example 3.}  Suppose $f(x) = \sin(x)$ and $g(x)=\arcsin(x)$.  
	\mar{$\sin{x}$ is invertible\\ on
		\mbox{$-\pi/2 \leq x \leq \pi/2$}}
Since $f$ is not one-to-one on its natural domain, we again need to restrict it in order for it to have an inverse.  By convention, the domain of $\sin(x)$ is taken to be $-\pi/2 \leq x \leq \pi/2$.  
\begin{align*}
f(g(a)) & =  \sin(\arcsin(a)) = a \;\; (\mbox{for} \;\; -1 \leq a \leq 1) \;\; \mbox{and}  \\
g(f(b)) & =  \arcsin(\sin(b)) = b \;\; (\mbox{for} \;\; -\pi/2 \leq b \leq \pi/2).
\end{align*}

	Each pair of inverse functions share corresponding properties, just as the logarithm and exponential functions do---the particular properties depending on the particular functions.  But two they all share are
\begin{itemize}
\item  The range of $f$ is the domain of $g$.
\item  The domain of $f$ is the range of $g$.
\end{itemize}
The exercises check this for examples 2 and 3.

\pagebreak

	Finally, the graphs---and therefore the derivatives---of a function and of its inverse are mirror images, exactly like those of the exponential and logarithm functions.  We begin with the same list of observations.
\begin{itemize}
\item  We know the point $(a, b)$ is on the graph of $y=g(x)$ if and only if
  $b = g(a)$.
\item  We know $b = g(a)$ says the same thing as $a = f(b)$.
\item  Finally, we know $a = f(b)$ is true if and only if the point $(b, a)$ is on the graph of $y = f(x)$.
\end{itemize}

\noindent
As before, putting our observations together, we have
\begin{center}
\begin{tabular}{c}
$(a, b)$ is on the graph of $y=g(x)$ \\
if and only if \\
$(b, a)$ is on the graph of $y=f(x)$.
\end{tabular}
\end{center}
\noindent Exactly as before, we have that the point $(a, b)$ and the point $(b, a)$ are reflections of each other about the line $y=x$.  
	\mar{The graphs of inverse functions are mirror images \ldots}
Since we have just seen that every point  $(a, b)$ on the graph of $y=g(x)$ corresponds to a point $(b, a)$ on the graph of $y=f(x)$, we again see that {\em the 
graphs of $y=g(x)$ and $y=f(x)$ are the reflections of each other     about the line $y=x$.}


	Finally, since the two graphs are reflections of one another, the local linear 
approximation of $g(x)$  at any point $(a, b)$ will be the mirror image of 
the local linear approximation of $f(x)$ at the point $(b, a)$. 
Any change in the $y$-value on one of these local lines will correspond to an equal change 
in the $x$-value in its mirror image, and vice versa.  
Just as before, it follows that the slopes of the two lines must be reciprocals of each other.  This says that the rate of change of $g(x)$ at $x = b$ is the reciprocal of 
the rate of change of $f(x)$ at $x = a$, where $a = g(b)$.  
	\mar{\ldots and their derivatives are reciprocals}
But the rate of change of $f(x)$ at $x = g(b)$ is just $f'(g(b))$.  Therefore the rate of change of $g(x)$ at $x = b$ is the reciprocal of this value, namely $1/f'(g(b))$.  We have thus proved the 
following result:  

\bigskip \noindent
{\bf Theorem 2}.   {\it If the functions} $f$ {\it and} $g$ {\it are inverses, then} $g$ {\it is locally linear at} $(b,a)$ {\it if and only if} $f$ {\it is locally linear at} $(a,b)$.  {\it When local linearity holds,}
$$
g'(b) = \frac{1}{f'(a)}.
$$

\subsection{Exercises}

\vspace{-10pt}
\num  Determine the numerical value of each of the following.
\smallskip

\noindent
\begin{tabular}{p{3cm}p{3cm}p{3cm}p{3cm}}
a) $\ln(2e)$ & b) $\ln(e^3)$ & c) $e^{-1}$ & 
	d) $\ln(\sqrt{e})$ 	\\ [2pt]
e) $e^{\ln 2}$ & f) $e^{3\ln 2}$ & g) $(e^{\ln 2})^3$ & 
	h) $e^{2\ln 3}$  \\[2pt]
i) $\ln 10$ & j) $\ln 10^3$  & k) $e^{\ln 10}$ & 
	l) $e^{\ln 1000}$  \\[2pt]
m) $\ln(1/e)$ & n) $\ln(1/2)$ & o) $e^{-\ln 2}$ & 
	p) $e^{-3\ln 2}$  
\end{tabular}

\numa In the text we noted that the function $\ln{x}$ is the solution to the initial-value problem
	$$
	\dfrac{dy}{dx} = \dfrac{1}{x} \qquad y(1) = 0,
	$$
so that we can use Euler's method to compute values for $\ln{x}$.  Use this method to evaluate $\ln{2}$ to 3 decimal places.  What value of $\Delta x$ gives the desired accuracy?
\sub If you now wanted to calculate $\ln{6}$ to 3 decimals, can you think of a better way to do it than simply starting at $x=1$ and running Euler's method out to $x=6$?  Remember the basic properties of logarithms, and figure out a way to use the results of part (a).
\sub Suppose you had figured out that $\ln{2}=0.693147\ldots$\,.  How would you use Euler's method to calculate $\ln{1300}$ quickly?  You might find the fact that $2^{10}=1024$ helpful.

\num  The rate of growth of the population of a particular country is proportional to
 the
population.  The last two censuses determined that the population in 1980 was
 40,000,000, and in 1985
it was 45,000,000.  What will the population be in 1995?

\num  Find the derivatives of the following functions.

\vspace{-10pt}
\noindent
\begin{minipage}[t]{2.5in}
\sub  $\ln(3x)$ 

\sub  $17\ln(x)$

\sub  $ \ln(e^{\displaystyle w})$ 
\end{minipage}
\noindent
\begin{minipage}[t]{2.5in}
\sub  $\ln(2^{\displaystyle t})$

\sub  $\pi\ln(3e^{\displaystyle 4s})$
\end{minipage}

\num  Suppose a bacterial population\index{growth!bacterial} grows so that its mass is 
	$$
	P(t) = 200e^{\displaystyle .12t} \quad \mbox{grams}
	$$
after $t$ hours.  Its initial mass is $P(0) = 200$ grams.  
When will its mass double, to 400 grams?  How much longer 
will it take to double again, to 800 grams?  After the 
population reaches 800 grams, how long will it take for yet 
another doubling to happen?  What is the {\em doubling time} 
of this population?


\num  Suppose a beam of X-rays whose intensity is $A$ rads (the ``rad'' is a unit of radiation)\index{radiation} falls perpendicularly on a heavy concrete wall.  After the rays have penetrated $s$ feet of the wall, the radiation intensity has fallen to 
	$$
	R(s) = Ae^{\displaystyle -.35s} \quad \mbox{rads}.
	$$
What is the radiation intensity 3 inches inside the wall; 18 inches?  (Your answers will be expressed in terms of $A$.) How far into the wall must the rays travel before their intensity is cut in half, to $A/2$?  How much further before the intensity is $A/4$?  

\num Virtually all living things take up carbon as they grow.  This
carbon comes in two principal forms: normal, stable
carbon---C$^{12}$---and radioactive carbon---C$^{14}$.  C$^{14}$
decays into C$^{12}$ at a rate proportional to the amount of C$^{14}$
remaining.  While the organism is alive, this lost C$^{14}$ is
continually replenished.  After the organism dies, though, the
C$^{14}$ is no longer replaced, so the percentage of C$^{14}$
decreases exponentially over time.  It is found that after 5730 years,
half the original C$^{14}$ remains.  If an archaeologist finds a bone
with only 20\% of the original C$^{14}$ present, how old is it?

\num The human population of the world appears to be growing
exponentially.  If there were 2.5 billion people in 1960, and 3.5
billion in 1980, how many will there be in 2010?

\num If bacteria increase at a rate proportional to the current
number, how long will it take 1000 bacteria to increase to 10,000 if
it takes them 17 minutes to increase to 2000?

\num Suppose sugar in water dissolves at a rate proportional to the
amount left undissolved.  If 40 lb. of sugar reduces to 12 lb. in 4
hours, how long will you have to wait until 99\% of the sugar is
dissolved?

\num Atmospheric pressure is a function of altitude.  Assume that at
any given altitude the rate of change of pressure with altitude is
proportional to the pressure there.  If the barometer reads 30 psi
(pounds per square inch) at sea level and 24 psi at 6000 feet above
sea level, how high are you when the barometer reads 20 psi?

\numa An important concept in many economic analyses is the idea of
      {\bf present value}\index{present value}.  It is used to compare
      the values of different possible payments made at different
      times.  As a simple example, suppose you had a small wood lot
      and had the choice of selling the timber on it now for \$5,000
      or waiting 10 years for the trees to get larger, at which point
      you estimate the timber could be sold for \$8,000.  To compare
      these two options, you need to convert the prospect of \$8,000
      ten years from now into an equivalent amount of money now---its
      present value.  This is the amount of money you would need to
      invest now to have \$8,000 in 10 years.  Suppose you thought you
      could invest money at an annual interest rate of 4\% compounded
      continuously.  If you invested \$5,000 now at this rate, then in
      10 years you would have $5000\,e^{.4}=\$7,459.12$.  That is,
      \$5,000 now is worth \$7,459.12 in 10 years---both amounts have
      the same present value.  Clearly \$8,000 in 10 years must have a
      slightly greater present value under the assumption of a 4\%
      annual interest rate.  What is it?  \sub On the other hand, if
      you can get a higher interest rate than 4\%, the present value
      of the \$8,000 will be much less.  What is the present value of
      a payment of \$8,000 ten years from now if the annual interest
      rate is 8\%?  \sub At what interest rate do \$5,000 now and
      \$8,000 in ten years have the same present value?

\num  Use properties of exp to prove the following properties of the logarithm.
  (Remember that $\ln{a} = b$  means $a = \exp{b}$.)

\sub  $\ln(1)=0$.

\sub  $\ln(m/n) = \ln(m) - \ln(n)$.

\sub  $\ln(m^n) = n \ln(m)$. 

\numa Use a graphing program to find a good numerical approximation to
$(\ln{x})'$ at $x=2$.  Make a short table, for decreasing interval
sizes $\Delta x$, of the quantity $\Delta (\ln{x})/\Delta x$.

\sub Use a graphing program to find a good numerical approximation to
$(e^x)'$ at $x=\ln(2)=0.6931\ldots$.  Make a short table for
decreasing interval sizes $\Delta x$, of the quantity $\Delta
(e^x)/\Delta x$.

\sub What is the relationship between the values you got in parts (a)
and (b)?

\num  Find a solution (using $\ln{x}$) to the differential equation
	$$
	f'(x) = 3/x \quad \mbox{satisfying} \quad f(1)=2.
	$$


\numa Find a formula using the natural logarithm function giving the
solution of $y'= a/x$ with $y(1)=b$.

\sub  Solve $P' = 2/t$ with $P(1)=5$.


\num Find the domain and range of each of the following pairs of
inverse functions.

\sub  $f(x) = x^2$ (restricted to $x \geq 0$) and $g(x)=\sqrt{x}$.

\sub $f(x) = \sin(x)$ (restricted to $-\pi/2 \leq x \leq \pi/2\/$) and
$g(x) = \arcsin(x)$.

\num Show that $f(x) = 1/x$ equals its own inverse.  What are the
domain and range of $f$?

\num Let $n$ be a positive integer. and let $f(x)=x^n$.  What is an
inverse of $f$?  How do we need to restrict the domain of $f$ for it
to have an inverse?  Caution: the answer depends on $n$.

\numa What is the inverse $g$ of the function $f(x)=1-3x$?  \sub Do
$f$ and $g$ satisfy Theorem 2?

\num  What is an inverse of $f(x) = x^2-4$?

\num Use the relationship between the derivatives of a function and
its inverse to find the indicated derivatives.

\sub  $g'(100)$ for $g(x)=\sqrt{x}$.

\sub  $g'(\sqrt{2}/2)$ for $g(x)=\arcsin(x)$.

\sub  $g'(1/2)$ for $g(x)=1/x$.  

\numa Use Theorem 2 and the fact that $(x^2)'=2x$ to derive the
formula for the derivative of $\sqrt{x}$.

\sub Use Theorem 2 and the fact that $(x^n)'=nx^{n-1}$ to derive the
formula for the derivative of $\sqrt[n]{x}$


\num Compare the rates of growth of $e^x$ and $b^x$ for both $e < b$
and $1 < b < e$.

%%%%%%%%%%%%%%%%%%%%%%%%%%%%%%%%%%%%%%%%%%%%%%%%%%%%%%%%%%%%%%%%%%%%%%%%%%%
%                                                                         %
%                                Section 5                                %
%                                                                         %
%%%%%%%%%%%%%%%%%%%%%%%%%%%%%%%%%%%%%%%%%%%%%%%%%%%%%%%%%%%%%%%%%%%%%%%%%%%

\section{The Equation $y' = f(t)$}

Most differential equations we have encountered express the rate of
growth of a quantity {\em in terms of the quantity itself}.  The
simplest models for biological growth had this form: $y' = ky$ and $y'
= ky^p$.  Even when several variables were present---as in the
$S$-$I$-$R$ model and the predator-prey models---it was most natural
to express the rates at which those variables change in terms of the
variables themselves.  Even the motion of a spring (pages
\pageref{spring-begin}--\pageref{spring-end}) was described that way:
the rate of change of position equalled the velocity, and the rate of
change of velocity was proportional to the position.

Sometimes, though, a differential equation \index{falling body}
\mar{The motion of\\ a falling body \ldots} will express the rate of
change of a variable {\em directly in terms of the input variable}.
For example, on page~\pageref{falling body0} we saw that the velocity
$dx/dt$ of a body falling under the sole influence of gravity is a
linear function of the time:
        $$
        \dfrac{dx}{dt} = -gt + v_0.
        $$
Here $x$ is the height of the body above the ground, $g$ is the acceleration due to gravity, and $v_0$ is the velocity at time $t = 0$.  This equation has the general form
        $$
        \dfrac{dx}{dt} = f(t),
        $$
where $f(t)$ is a given function of $t$.  We will now consider special methods that can be used to study differential equations of this special form.  

\subsection{Antiderivatives}

        To solve the equation of motion of a body 
        \mar{\ldots and its solution}
falling under gravity, we must find a function $x(t)$ whose derivative is given as
        $$
        x' = -gt + v_0.
        $$
We can call upon our knowledge of the rules of differentiation to find $x$.  Consider $-gt$ first.  
        \mar{$-gt$ is the derivative\\ of $-gt^2/2$}
What function has $-gt$ as its derivative?  We can start with $t^2$, whose derivative is $2t$.  Since we want the derivative to turn out to be $-gt$, we can reason this way:
        $$
        -gt = -\dfrac{g}{2} \cdot 2t = -\dfrac{g}{2} \times 
                \mbox{the derivative of $t^2$}.
        $$
This leads us to identify $-gt^2/2$ as a function whose derivative is $-gt$.  Check for yourself that this is correct by differentiating $-gt^2/2$.

Now consider $v_0$, the other part of $dx/dt$.  What function has the
constant $v_0$ as its derivative?
%       
\mar{$v_0$ is the derivative\\ of $v_0t$ and of $v_0t + b$} 
%
A derivative is a rate of growth, and we know that the linear functions
are precisely the ones that have constant growth rates.  Furthermore,
the rate is the multiplier for a linear function, so we conclude that
any linear function of the form $v_0 t + b$ has derivative $v_0$.

If we put the two pieces together, we find that 
\[
x(t) = -\dfrac{g}{2} \, t^2 + v_0 t + b
\]
is a solution to the differential equation, for any value of $b$.
(Recall from section~2 that a differential equation can have many
solutions.)  We constructed this formula for $x(t)$ by ``undoing''
%
\mar{The antiderivative\\ of a function}
%
the process of differentiation, a process sometimes called {\bf
  antidifferentiation}.\index{antidifferentiation} The function
produced is called an {\bf antiderivative}.\index{antiderivative}
Thus:
\begin{gather*}
-\dfrac{g}{2} t^2 + v_0 t + b 
              \text{ is an \emph{antiderivative} of} -gt + v_0 \\
\text{because }
-gt + v_0 \text{ is the \emph{derivative} of } -\dfrac{g}{2} t^2 + v_0 t + b.
\end{gather*}
Note that a function 
        \mar{All the functions $F(x) +C$ are antiderivatives of $F'(x)$}
has only one derivative, but it has many antiderivatives.  If $F(t)$ is an antiderivative of $f(t)$, then so is $F(t) + C$, where $C$ is any constant.

\enlargethispage*{22pt}

The list of functions and their derivatives that we compiled in
chapter~3 (see page~\pageref{basic derivs}) can be ``turned around''
to become a list of functions and their antiderivatives.  Note that
the antiderivative column should really be labelled ``an
antiderivative'' since we could add a constant to any of the listed
functions and still have an antiderivative for the function in the
first column.  \addtolength{\arraycolsep}{8pt}
        $$
        \begin{array}{cc}
        \mbox{function} & \mbox{antiderivative} \\[2pt] \hline
        \rule{0pt}{14pt} 0 \rule{0pt}{14pt} & c \\[2pt]
        x^p & \hspace*{.6in} \frac{1}{p+1} \, x^{p+1} 
                         \mbox{ (if $p\neq -1$)}\\[2pt]
        x^{-1} & \ln{x} \\[2pt]
        \sin{x} & -\cos{x} \\[2pt]
        \cos{x} & \sin{x} \\[2pt]
        \exp{x} = e^x & \exp{x} = e^x \\[2pt]
        b^x & \frac{1}{\ln{b}}\, b^x 
        \end{array}
        $$
        \addtolength{\arraycolsep}{-8pt}%
Notice the formula for the antiderivative of $x^p$ 
%
\mar{Every power of $x$ has an antiderivative}
%
requires $p+1 \neq 0$, that is, $p \neq -1$.  This leaves out
$x^{-1}$.  However, the antiderivative of $x^{-1}$ is $\ln{x}$, so no
power of $x$ is excluded from the table.

\pagebreak

We also had differentiation rules that told us how to deal with
different {\em combinations\/} of functions.  Each of these rules has
an analogue in antidifferentiation.  The simplest combinations are a
sum and a constant multiple.
        \addtolength{\arraycolsep}{15pt}
        $$
        \begin{array}{cc}
        \mbox{function} & \mbox{antiderivative} \\[2pt] \hline
        \rule{0pt}{14pt} f(x) \rule{0pt}{14pt} & F(x) \\[2pt]
        g(x) & G(x) \\[2pt]
        c \cdot f(x) & c \cdot F(x) \\[2pt]
        f(x) + g(x) & F(x) + G(x) 
        \end{array}
        $$
        \addtolength{\arraycolsep}{-15pt}
We defer a discussion of the analogue of the chain rule to chapter~11.

With just these rules we can find the antiderivative of any
polynomial, for instance.  (Recall that a polynomial is a sum of
constant multiples of powers of the input variable.)  Here is a
collection of sample antiderivatives that illustrate the various
rules.  To emphasize the fact that antiderivatives are determined only
up to an additive constant, various constants have been tacked
on---any other constant would work just as well.  You should compare
this table with the one on page~\pageref{deriv table}.\label{antideriv
  table} \addtolength{\arraycolsep}{8pt}
        $$
        \begin{array}{cc}
        \mbox{function} & \mbox{antiderivative} \\[2pt] \hline
        \rule{0pt}{16pt} 5x^4 - 2x^3\rule{0pt}{16pt} & 
                x^5 - \frac{1}{2}x^4 +17 \\[6pt]
        5x^4 - 2x^3 + 17x & 
                x^5 - \frac{1}{2}x^4 + \frac{17}{2}x^2 -243.77 \\[6pt]
        6 \cdot 10^z + 17/z^7 & 
                 6 \cdot 10^z/\ln{10} -  17/6z^6 +.002 \\[6pt]
        3\sin{t} - 2t^3 & -3\cos{t} - \frac{1}{2}t^4 +5\ln{7} \\[6pt]
        \pi\cos{x} + \pi^2 & \pi\sin{x} + \pi^2 x  - 12e^{7.21}
        \end{array}
        $$
        \addtolength{\arraycolsep}{-8pt}


\subsection{Euler's Method Revisited}
\index{Euler's method}

If we know the formula for an antiderivative of $f(t)$, then we can
write down a solution to the differential equation $dy/dt = f(t)$.
For example, the general solution to
        $$
        \dfrac{dy}{dt} = 12t^2 + \sin{t}
        $$
is $y = 4t^3 - \cos{t} + C$.  In such a case we have a shortcut to solving the differential equation without needing to use Euler's method.  Often, though, there is no formula for an antiderivative of $f(t)$---even when $f(t)$ itself has a simple formula.  There is no formula for the antiderivative of $\cos(t^2)$, or $\sin{t}/t$, or $\sqrt{1 + t^3}$, for instance.  In other cases, $f(t)$ may not even be given by a formula.  It may be a data function, given as a graph made by a pen tracing on a moving sheet of graph paper.

        Whether we can find a formula for an antiderivative of $f(t)$ or not, we can still solve the differential equation $dy/dt = f(t)$ by Euler's method. It turns out that Euler's method takes on a relatively simple form in such cases.  Let's investigate this in the following context.  

        Let $V$ be the volume of water in a reservoir \index{reservoir}
        \mar{The volume\\ of a reservoir\\ varies over time}
serving a small town, measured in millions of gallons.  Then $V$ is a function of the time $t$, measured in days.  Rainfall adds water to the reservoir, while evaporation and consumption by the townspeople take it away.  Let $f$ be the {\em net rate\/} at which water is flowing into the reservoir, in millions of gallons per day.  Sometimes $f$ will be positive---when rainfall exceeds evaporation and consumption---and sometimes $f$ will be negative.  The net inflow rate varies from day to day; that is, $f$ is a function of time: $f~=~f(t)$.  Our model of the reservoir is the differential equation
        $$
        \dfrac{dV}{dt} = f(t) \ 
                \mbox{millions of gallons per day}.
        $$

        Suppose $f(t)$ is measured every two days, 
        \mar{The net inflow rate}
and those measurements are recorded in the following table.
        \addtolength{\arraycolsep}{10pt}
        {\small
        $$
\begin{array}{cc}
        \mbox{time $t$} & \mbox{rate $f(t)$} \\ [-2pt]
        \mbox{(days)} & (10^6 \times \ \mbox{gals. per day}) 
                \\ [2pt] \hline
        \begin{array}{r}
        \rule{0pt}{14pt} 0 \\ 2 \\ 4 \\ 6 \\ 8 \\ 10 \\ 12 
        \end{array}
        &
        \begin{array}{r}
        \rule{0pt}{14pt} .34 \\ .11 \\ -.07 \\ -.23 \\ -.14 \\
                .03 \\ .08 
        \end{array}
\end{array}
        $$
}\addtolength{\arraycolsep}{-10pt}%
%
Note that in this table we are able to write down the rate for all
values of $t$ immediately, without having to calculate the
intermediate values of the dependent variable $V$.  This is in marked
contrast with most of the examples we've looked at in this course
where we had to know the values of all the variables for any time $t$
before we could calculate the new rate value at that time.  It is this
simplification that gives differential equations of the form $y' =
f(t)$ their special structure.

If we assume the value of $f(t)$ remains constant for the two days
after each measurement is made, we can approximate the total change in
$V$ over these 14 days.  The following table does this; it tells us
how much $V$ changes over each two-day period, and
also\index{change!accumulated}
%
\mar{The accumulated change in $V$}
%
the total \emph{accumulated} change in $V$ by the end that
period. Since $\Delta t = 2$ days, we calculate $\Delta V$ by
       $$
        \Delta V = V' \cdot \Delta t = f(t) \cdot \Delta t 
                = 2 \cdot f(t).
        $$
\addtolength{\arraycolsep}{6pt}
{\small 
        $$
\begin{array}{cccc}
        \mbox{starting} & \mbox{current} & \mbox{accumulated} &
                \mbox{ending} \\ [-2pt]
        t & \Delta V &  \Delta V & t \\[2pt] \hline
        \begin{array}{r}
        \rule{0pt}{14pt} 0 \\ 2 \\ 4 \\ 6 \\ 8 \\ 10 \\ 12 
        \end{array}
        &
        \begin{array}{r}
        \rule{0pt}{14pt} .68 \\ .22 \\ -.14 \\ -.46 \\ -.28 \\
                .06 \\ .16 
        \end{array}
        &
        \begin{array}{r}
        \rule{0pt}{14pt} .68 \\ .90 \\ .76 \\ .30 \\ .02 \\
                .08 \\ .24 
        \end{array}
        &
        \begin{array}{r}
        \rule{0pt}{14pt} 2 \\ 4 \\ 6 \\ 8 \\ 10 \\ 12 \\ 14
        \end{array}
\end{array}
        $$
}\addtolength{\arraycolsep}{-6pt}

At the end of the 14 days, $V$ has accumulated a total change of .24
million gallons.  Notice this does not depend on the initial size of
$V$.  If $V$ had been 92.64 million gallons at the start, it would be
92.64 + .24 = 92.88 million gallons at the end.  If it had been only 2
million gallons at the start, it would be 2 + .24 = 2.24 million
gallons at the end.  Other models do not behave this way: in two
weeks, a rabbit population of 900 will change much more than a
population of 90.  The total change in $V$ is independent of $V$
because the {\em rate\/} at which $V$ changes is independent of $V$.

\enlargethispage*{40pt}
We can therefore use Euler's method 
%
\mar{Calculating just\\ the accumulated\\ change in~$y$}
%
to solve any differential equation of the form $dy/dt = f(t)$ {\em
  independently\/} of an initial value for $y$.  We just calculate the
total accumulated change in $y$, and add that total to any given
initial $y$.  Here is how it works when the initial value of $t$ is
$a$, and the time step is $\Delta t$.  \addtolength{\arraycolsep}{6pt}
{\small
        $$
\begin{array}{cccc}
        \mbox{starting} & \mbox{current} & \mbox{accumulated} &
                \mbox{ending} \\ [-2pt]
        t & \Delta y &  \Delta y & t \\[2pt] \hline
        \begin{array}{c}
                \rule{0pt}{14pt} a \rule{0pt}{14pt} \\
                a + \Delta t \\
                a + 2\Delta t \\
                a + 3\Delta t \\
                \vdots \\[3pt]
                a + (n-1)\Delta t 
        \end{array}
        &
        \begin{array}{c}
                \rule{0pt}{14pt} f(a) \cdot \Delta t \\
                f(a + \Delta t) \cdot \Delta t  \\
                f(a + 2\Delta t) \cdot \Delta t  \\
                f(a + 3\Delta t) \cdot \Delta t  \\
                \vdots \\[3pt]
                %\hphantom{\Delta t) \cdot \Delta t}
                f(a + (n-1)\Delta t) \cdot \Delta t  
        \end{array}
        &
        %\left\{
        \begin{array}{c}
                \rule{0pt}{14pt} \mbox{previously}
                        \rule{0pt}{14pt}\\
                \mbox{accumulated} \\
                \Delta y \\
                \vphantom{\vdots} + \vphantom{\vdots} \\
                \mbox{current} \\
                \Delta y
        \end{array}
        %\right\}
        &
        \begin{array}{c}
                \rule{0pt}{14pt} a + \Delta t \rule{0pt}{14pt} \\
                a + 2\Delta t \\
                a + 3\Delta t \\
                a + 4\Delta t \\
                \vdots \\[3pt]
                a + n\,\Delta t 
        \end{array}
\end{array}
        $$
%       \addtolength{\arraycolsep}{-6pt}
}


The third column is too small to hold the 
%
\mar{The accumulated change in $y$}
%
values of ``accumulated $\Delta y$.''  Instead, it contains the
instructions for obtaining those values.  It says: to get the current
value of ``accumulated $\Delta y$,'' add the ``current $\Delta y$'' to
the previous value of ``accumulated $\Delta y$.''

Let's use Euler's method to find the accumulated $\Delta y$ when $t =
4$, given that\label{accum delta y}
\[
\dfrac{dy}{dt} = \cos(t^2) 
\]
and $t$ is initially 0.  If we use 8 steps, then $\Delta t = .5$ and
we obtain the following: 
\addtolength{\arraycolsep}{2pt} {\small
        $$
\begin{array}{cccc}
        \mbox{starting} & \mbox{current} & \mbox{accumulated} &
                \mbox{ending} \\ [-2pt]
        t & \Delta y &  \Delta y & t \\[2pt] \hline
        \begin{array}{r}
        \rule{0pt}{14pt} 0 \\ .5 \\ 1.0 \\ 1.5 \\ 2.0 \\ 2.5
                \\ 3.0 \\ 3.5
        \end{array}
        &
        \begin{array}{r}
        \rule{0pt}{14pt} .5000\\ .4845 \\ .2702 \\ -.3141
                \\ -.3268 \\ .4997 \\ -.4556 \\ .4752 
        \end{array}
        &
        \begin{array}{r}
        \rule{0pt}{14pt} .5000 \\ .9845 \\ 1.2546 \\ .9405 
                \\ .6137 \\ 1.1134 \\ .6579 \\ 1.1330
        \end{array}
        &
        \begin{array}{r}
        \rule{0pt}{14pt} .5 \\ 1.0 \\ 1.5 \\ 2.0 \\ 2.5
                \\ 3.0 \\ 3.5 \\ 4.0
        \end{array}
\end{array}
        $$
}\addtolength{\arraycolsep}{-2pt}%
%
The following program generated the last three columns of this
table.\index{program!TABLE}
\begin{center}
        {\bf Program: TABLE}\\
        \end{center}
\begin{verbatim}
  DEF fnf (t) = COS(t ^ 2)
  tinitial = 0
  tfinal = 4
  numberofsteps = 2 ^ 3
  deltat = (tfinal - tinitial) / numberofsteps
  t = tinitial
  accumulation = 0
  FOR k = 1 TO numberofsteps
          deltay = fnf(t) * deltat
          accumulation = accumulation + deltay
          t = t + deltat
          PRINT deltay, accumulation, t
  NEXT k
\end{verbatim}
TABLE is a modification of the program SIRVALUE
(page~\pageref{sirvalue}).  To emphasize the fact that it is the
accumulated change that matters rather than the actual value of $y$,
we have modified the program accordingly.  Note that {\tt
  accumulation} always starts at 0, no matter what the initial value
of $y$ is.  The first line of the program takes advantage of a
capacity most programming languages have to define functions which can
then be referred to elsewhere in the program.

As usual, to find the exact value of the accumulated $\Delta y$, it is
necessary to recalculate, using more steps and smaller step sizes
$\Delta t$.  If we use TABLE to do this, we find
        $$
        \begin{array}{cc}
        \mbox{number} & \mbox{accumulated} \\[-2pt]
        \mbox{of steps} & \Delta y \\[2pt] \hline
        \rule{0pt}{15pt} 2^{3\hphantom{0}} \rule{0pt}{15pt} 
                & 1.13304 \\
        2^{6\hphantom{0}} & \hphantom{0}.65639 \\
        2^{9\hphantom{0}} & \hphantom{0}.60212 \\
        2^{12} & \hphantom{0}.59542 \\
        2^{15} & \hphantom{0}.59458 \\
        2^{18} & \hphantom{0}.59448 
        \end{array}
        $$ 
Thus we can say that if $dy/dt = \cos(t^2)$, then $y$ increases by
.594\ldots when $t$ increases from 0 to 4.

In the same way we changed SIRVALUE to produce the program SIRPLOT
(page~\pageref{sirplot}), we can change the program TABLE into one
that will {\em plot\/} the values of $y$.  In the following program
all those changes are made, and one more besides: we have increased
the number of steps to 400 to get a closer approximation to the true
values of $y$.\index{program!PLOT}  The output of PLOT is shown
immediately below.

\vspace{135pt}

\special{psfile=../ps/calc1/pic260.ps}

\begin{center}
The accumulated $\Delta y$ when $dy/dt = \cos(t^2)$
\end{center}



\begin{center}
{\bf Program: PLOT}\label{program PLOT}
\end{center}

\noindent
\rule{12pt}{0pt}{\sl Set up GRAPHICS}

\vspace{-11pt}
\begin{verbatim}
  DEF fnf (t) = COS(t ^ 2)
  tinitial = 0
  tfinal = 4
  numberofsteps = 400
  deltat = (tfinal - tinitial) / numberofsteps
  t = tinitial
  accumulation = 0
  FOR k = 1 TO numberofsteps
        deltay = fnf(t) * deltat
\end{verbatim}

\vspace{-12pt}
\rule{27pt}{0pt} {\sl Plot the line from {\tt (t, accumulation)}}\\[-1pt]
\rule{60pt}{0pt} {\sl to {\tt (t + deltat, accumulation + deltay)}}

\vspace{-13pt}
\begin{verbatim}
        accumulation = accumulation + deltay
        t = t + deltat
  NEXT k
\end{verbatim}

Let's compare our reservoir model with population growth.  The rate at
which a population grows depends, in an obvious way, on the size of
the population.  By contrast, the rate at which the reservoir fills
does {\em not\/} depend on how much water there is in the reservoir.
%
\mar{Exogenous and endogenous factors}
%
It depends on factors {\em outside\/} the reservoir: rainfall and
consumption.  These factors are said to be {\bf exogenous} (from the
Greek {\em exo-}, ``outside'' and {\em -gen\/}, ``produced,'' or
``born'').\index{exogenous factor} The opposite is called an {\bf
  endogenous} factor (from the Greek {\em endo-},
``within'').\index{endogenous factor} Evaporation is an endogenous
factor for the reservoir model; population size is certainly an
endogenous factor for a population model.

Precisely because exogenous factors are ``outside the system,'' we
need to be {\em given\/} the information on how they vary over time.
In the reservoir model, this information appears in the function
$f(t)$ that describes the rate at which $V$ changes.  In general, if
$y$ depends on exogenous factors that vary over time, we can expect
the differential equation for $y$ to involve a function of time:
        $$
        \dfrac{dy}{dt} = f(t)
        $$
Thus, we can view this section as dealing with models that involve
exogenous factors.

\newpage

\aside{The differential equation of motion for a falling body, $dx/dt
  = -gt + v_0$, indicates that gravity is an exogenous factor.  In
  Greek and medieval European science, the reason an object fell to
  the ground was assumed to lie within the object itself---it was the
  object's ``heaviness.''  By making the cause of motion exogenous,
  rather than endogenous, Galileo and Newton started a scientific
  revolution.}

\subsection{Exercises}

\vspace{-10pt}

\num Find a formula $y = F(t)$ for a solution to the differential
equation $dy/dt = f(t)$ when $f(t)$ is

\vspace{-10pt}
\noindent
\begin{minipage}[t]{2.4in}
\sub $5t - 3$

\sub $t^6 - 8t^5 + 22\pi^3$

\sub $5e^t - 3\sin{t}$
\end{minipage}
\hfill
\begin{minipage}[t]{2.4in}

\sub $12\sqrt{t}$

\sub $2^t + 7/t^9$
\sub $5e^{4t}-1/t$
\end{minipage}

\num  Find $G(5)$ if $y = G(x)$ is the solution to the initial value problem
        $$
        \dfrac{dy}{dx} = \dfrac{1}{x^2} \qquad y(2) = 3.
        $$

\num  Find $F(2)$ if $y = F(x)$ is the solution to the initial value problem
        $$
        \dfrac{dy}{dx} = \dfrac{1}{x} \qquad y(1) = 5.
        $$

\num  Find $H(3)$ if $y = H(x)$ is the solution to the initial value problem
        $$
        \dfrac{dy}{dx} = x^3-7x^2+19 \qquad y(-1) = 5.
        $$

\num  Find $L(-2)$ if $y = L(x)$ is the solution to the initial value problem
        $$
        \dfrac{dy}{dx} = e^{3x} \qquad y(1) = 6.
        $$

\numa Sketch the graph of the solution to the initial value problem
        $$
        \dfrac{dy}{dx} = \sin{x} \qquad y(0) = 1
        $$
over the interval $0 \leq x \leq 4\pi$.
\sub By finding a suitable antiderivative, evaluate $y(2)$.

\numa  Sketch the graph of the solution to the initial value problem
        $$
        \dfrac{dy}{dx} = \sin(x^2) \qquad y(0) = 0
        $$ 
over the interval $0 \leq x \leq 5$.  (This one can't be done by
finding a formula for the antiderivative.)

\sub  What is the slope of the solution graph at $x = 0$?  Does your graph show this?

\sub  How many peaks (local maxima) does the solution have on the interval $0 \leq x \leq 5$?

\sub  What is the maximum value that the solution achieves on the interval $0 \leq x \leq 5$?  For which value of $x$ does this happen?

\sub What is $y(6)$?


\numa  What is the accumulated change in $y$ if $dy/dt = 3t^2 - 2t$ and $t$ increases from 0 to 1?  What if $t$ increases from 1 to 2?  What if $t$ increases from 0 to 2?  

\sub  Sketch the graph of the accumulated change in $y$ as a function of $t$.  Let $0 \leq t \leq 2$.


\numa  Here's another problem for which there is no formula for an antiderivative. Sketch the graph of the solution to the initial value problem 
        $$
        \dfrac{dy}{dx} = \dfrac{\sin{x}}{x} \qquad y(0) = 0
        $$
on the interval $0 \leq x \leq 40$.  [Note: $\sin{x}/x$ is not defined when $x = 0$, so take the initial value of $x$ to be .00001.  That is, use $y(.00001) = 0$.]

\sub  How many peaks (local maxima) does the solution have on the interval $0 \leq x \leq 40$?  

\sub  What is the maximum value of the solution on the interval $0 \leq x \leq 40$?  For which $x$ is this maximum achieved?



\newpage

%%%%%%%%%%%%%%%%%%%%%%%%%%%%%%%%%%%%%%%%%%%%%%%%%%%%%%%%%%%%%%%%%%%%%%%%%%%
%                                                                         %
%                                Section 6                                %
%                                                                         %
%%%%%%%%%%%%%%%%%%%%%%%%%%%%%%%%%%%%%%%%%%%%%%%%%%%%%%%%%%%%%%%%%%%%%%%%%%%

\section{Chapter Summary}


\subsection{The Main Ideas}

\begin{itemize}

\item  A {\bf system of differential equations} expresses the derivatives of a set of functions in terms of those  functions and the input variable.  

\item  An {\bf initial value problem} is a system of differential equations together with values of the functions for some specified value of the input variable.  

\item  Many processes in the physical, biological, and social sciences are {\bf modelled} as initial value problems.

\item  A {\bf solution to a system of differential equations} is a set of functions which make the equations {\em true\/} when they and their derivatives are substituted into the equations.

\item  A {\bf solution to an initial value problem} is a set of functions that solve the differential equations and satisfy the initial conditions.  Typically, a solution is unique.

\item  {\bf Euler's method} provides a recipe to find the solution to an initial value problem.   

\item  In special circumstances it is possible to find {\bf formulas} for the solution to a system of differential equations.  If the differential equations involve {\bf parameters}, the solutions will, too.

\item  Systems of differential equations define functions as their solutions.  Among the most important are the {\bf exponential} and {\bf logarithm functions}.

\item  The {\bf natural logarithm} function is the inverse of the exponential function.  

\item  The {\bf graphs} and the {\bf derivatives} of a function and its inverse are connected geometrically to each other by {\bf reflection}.  
 
\item  Exponential functions $b^{\displaystyle x}$ {\bf grow to infinity} faster than any power of~$x$.

\item  The solution to $dy/dx = f(x)$ is an {\bf antiderivative} of $f$---that is, a function whose derivative is $f$.


\end{itemize}


\subsection{Expectations}

\begin{itemize}

\item  You should be able to use computer programs to produce tables and graphs of solutions to initial value problems.

\item  You should be able to check whether a system of differential equations reflects the hypotheses being made in constructing a model of a process.

\item  You should be able to verify whether a set of functions given by formulas is a solution to a system of differential equations.

\item  You should be familiar with the basic properties of the exponential and logarithm functions.  

\item  You should be able to express solutions to initial value problems involving exponential growth or decay in terms of the exponential function.

\item  You should be able to solve $dy/dx = f(x)$ by antidifferentiation when $f(x)$ is a basic function or a simple combination of them.  

\item  You should be able to analyze and graph the inverse of a given function.


\end{itemize}

\cleardoublepage
